\documentclass[12pt,a4paper]{article}
\usepackage[utf8]{inputenc}
\usepackage[spanish]{babel}
\usepackage{geometry}
\usepackage{enumitem}
\usepackage{hyperref}
\usepackage{fancyhdr}
\usepackage{titlesec}
\usepackage{setspace}
\usepackage{times}
\usepackage{longtable}
\usepackage{array}
\usepackage{xcolor}
\usepackage{colortbl}
\usepackage{lscape}
\usepackage{afterpage}
\usepackage{caption}
\usepackage{graphicx}
\usepackage{float}
\usepackage{pdflscape}

\usepackage{indentfirst}  % Sangría en primer párrafo después de título (APA 7)

% Configuración de márgenes APA 7 (1 pulgada = 2.54 cm)
\geometry{
    left=2.54cm,
    right=2.54cm,
    top=2.54cm,
    bottom=2.54cm,
    headheight=15pt
}

% Interlineado 1.5 (compromiso entre legibilidad y economía de espacio)
\onehalfspacing

% Sangría de primera línea (0.5 pulgadas = 1.27 cm) - APA 7
\setlength{\parindent}{1.27cm}
\setlength{\parskip}{0pt}

% Configuración de encabezados APA 7
\pagestyle{fancy}
\fancyhf{}
\fancyhead[R]{\small\thepage}
\renewcommand{\headrulewidth}{0pt}

% Configuración de títulos APA 7
% Nivel 1: Centrado, Negrita, Título en Mayúsculas y Minúsculas
\titleformat{\section}
    {\normalfont\fontsize{12}{15}\bfseries\centering}
    {\thesection}{1em}{}
\titlespacing*{\section}{0pt}{12pt}{0pt}

% Nivel 2: Alineado a la Izquierda, Negrita, Título en Mayúsculas y Minúsculas
\titleformat{\subsection}
    {\normalfont\fontsize{12}{15}\bfseries\raggedright}
    {\thesubsection}{1em}{}
\titlespacing*{\subsection}{0pt}{12pt}{0pt}

% Nivel 3: Alineado a la Izquierda, Negrita, Cursiva, Título en Mayúsculas y Minúsculas
\titleformat{\subsubsection}
    {\normalfont\fontsize{12}{15}\bfseries\itshape\raggedright}
    {\thesubsubsection}{1em}{}
\titlespacing*{\subsubsection}{0pt}{12pt}{0pt}

% Colores para tablas
\definecolor{headercolor}{RGB}{41, 128, 185}
\definecolor{rowgray}{RGB}{245, 245, 245}
\definecolor{headerblue}{RGB}{68, 114, 196}
\definecolor{columnblue}{RGB}{189, 215, 238}
\definecolor{headergray}{RGB}{217, 217, 217}

% Espaciado entre filas de tablas (1.2 para compensar doble espaciado del texto)
\renewcommand{\arraystretch}{1.3}

% Configuración de listas según APA 7
\setlist{noitemsep, topsep=6pt, parsep=0pt, partopsep=0pt}

% Ajuste de ancho de tablas para evitar desbordamiento
\setlength{\tabcolsep}{4pt}

\begin{document}

% ============================================
% CARÁTULA APA 7
% ============================================
\begin{titlepage}
    \centering
    \vspace*{1cm}
    
    % Título del documento (centrado, negrita, doble espacio)
    {\bfseries\large
    Sistema de Gestión de Eventos Deportivos (SIGED)\\[12pt]
    Especificación de Requisitos de Software\\[12pt]
    Versión 1.0
    \par}
      \vspace{2cm}

% LOGO
    \begin{center}
    \includegraphics[width=8cm]{sigetlogo1.png}
    \end{center}

    \vspace{2cm}
    
    
    
    % Autores (centrado)
    {\normalsize
    Equipo de Desarrollo SIGED\\[6pt]
    \par}
    
  
    
    % Afiliación institucional (centrado)
    {\normalsize
    Facultad de Ingeniería de Sistemas\\[6pt]
    Universidad Nacional Agraria de la Selva\\[6pt]
    Tingo María, Perú
    \par}
    
    \vfill
    
    % Fecha (centrada al final)
    {\normalsize
    Febrero 2026
    \par}
    
\end{titlepage}

% Resetear numeración de página después de la portada
\setcounter{page}{2}

\newpage

% Tabla de contenido (opcional en APA 7, pero útil para documentos extensos)
\tableofcontents

\newpage

\section{Introducción}

\subsection{Propósito del Documento}

El presente documento constituye la Especificación de Requisitos de Software (ERS) del Sistema de Gestión de Eventos Deportivos (SIGED), en adelante denominado "el Sistema". Este documento ha sido elaborado siguiendo los lineamientos establecidos por el estándar IEEE 830-1998 (IEEE Recommended Practice for Software Requirements Specifications) y las mejores prácticas recomendadas por SWEBOK (Software Engineering Body of Knowledge).

La finalidad de este documento es proporcionar una descripción completa, precisa y verificable de los requisitos tanto del proyecto como del producto de software, estableciendo las bases contractuales y técnicas para el desarrollo, verificación y validación del Sistema. Este documento está dirigido a:

\begin{itemize}[noitemsep]
    \item \textbf{Stakeholders y patrocinadores:} Para validar la alineación con objetivos estratégicos institucionales.
    \item \textbf{Equipo de desarrollo:} Como especificación técnica detallada para implementación.
    \item \textbf{Equipo de aseguramiento de calidad:} Como base para diseño de casos de prueba y validación.
    \item \textbf{Gestores de proyecto:} Para planificación, seguimiento y control del proyecto.
    \item \textbf{Arquitectos de software:} Para diseño de arquitectura y componentes del sistema.
\end{itemize}

\subsection{Contexto y Motivación del Proyecto}

\subsubsection{Problemática Actual}

Las instituciones educativas de nivel superior en Perú enfrentan desafíos significativos en la organización y gestión de eventos deportivos interuniversitarios. Actualmente, la mayoría de universidades gestionan estos procesos mediante métodos tradicionales que incluyen:

\begin{itemize}[noitemsep]
    \item \textbf{Documentación manual en papel:} Inscripciones, actas de partidos y reportes elaborados manualmente, propensos a errores y pérdida de información.
    \item \textbf{Hojas de cálculo desarticuladas:} Uso de múltiples archivos Excel no integrados para gestionar equipos, calendarios y resultados.
    \item \textbf{Comunicación fragmentada:} Dependencia de WhatsApp, correos electrónicos y llamadas telefónicas sin trazabilidad formal.
    \item \textbf{Ausencia de información en tiempo real:} Imposibilidad de consultar resultados, tablas de posiciones o calendarios actualizados de manera inmediata.
    \item \textbf{Procesos lentos y burocráticos:} Trámites administrativos que requieren múltiples aprobaciones presenciales y firmas físicas.
\end{itemize}

Esta situación genera consecuencias negativas que afectan la calidad y eficiencia de los eventos deportivos:

\begin{itemize}[noitemsep]
    \item Retrasos constantes en la publicación de resultados oficiales.
    \item Errores en el cálculo de tablas de posiciones y clasificaciones.
    \item Falta de transparencia en procesos disciplinarios y arbitrales.
    \item Duplicación de esfuerzos administrativos y uso ineficiente de recursos humanos.
    \item Baja participación estudiantil debido a procesos de inscripción complejos.
    \item Imposibilidad de generar estadísticas históricas y analítica deportiva.
\end{itemize}

\subsubsection{Oportunidad de Mejora}

La Universidad Nacional Agraria de la Selva (UNAS), consciente de la importancia del deporte universitario como parte integral de la formación académica y el bienestar estudiantil, ha identificado la necesidad estratégica de modernizar sus procesos deportivos mediante la implementación de soluciones tecnológicas.

El proyecto SIGED surge como una iniciativa institucional que busca aprovechar las capacidades tecnológicas actuales para:

\begin{itemize}[noitemsep]
    \item \textbf{Digitalizar completamente} los procesos administrativos deportivos, eliminando el uso de papel y documentación física.
    \item \textbf{Centralizar la información} en una plataforma unificada accessible desde cualquier dispositivo con conexión a internet.
    \item \textbf{Automatizar tareas repetitivas} como la generación de fixtures, cálculo de resultados y notificaciones.
    \item \textbf{Democratizar el acceso} a información deportiva mediante un portal público de consulta.
    \item \textbf{Garantizar transparencia} en todos los procesos mediante trazabilidad digital completa.
    \item \textbf{Facilitar la toma de decisiones} con reportes analíticos y métricas de desempeño en tiempo real.
\end{itemize}

\subsubsection{Alineación Estratégica}

El desarrollo de SIGED se alinea con los siguientes objetivos estratégicos institucionales:

\begin{enumerate}[noitemsep]
    \item \textbf{Transformación Digital:} Impulsar la adopción de tecnologías de información en procesos administrativos universitarios.
    \item \textbf{Calidad Educativa:} Fortalecer la formación integral del estudiante mediante la mejora de servicios deportivos.
    \item \textbf{Innovación y Desarrollo:} Promover proyectos de desarrollo tecnológico ejecutados por estudiantes de Ingeniería de Sistemas.
    \item \textbf{Transparencia y Eficiencia:} Optimizar el uso de recursos institucionales mediante procesos digitales auditables.
    \item \textbf{Bienestar Estudiantil:} Mejorar la experiencia deportiva universitaria, incentivando la participación activa de la comunidad estudiantil.
\end{enumerate}

\subsubsection{Beneficiarios del Sistema}

El Sistema beneficiará directamente a múltiples actores de la comunidad universitaria:

\begin{itemize}[noitemsep]
    \item \textbf{Estudiantes deportistas (aprox. 5,000):} Inscripciones simplificadas, acceso a información en tiempo real, seguimiento de estadísticas personales.
    \item \textbf{Personal administrativo deportivo (12 personas):} Reducción del 70\% en carga administrativa manual, herramientas digitales eficientes.
    \item \textbf{Docentes coordinadores (30 personas):} Mayor control sobre equipos de sus facultades, reportes automáticos de participación.
    \item \textbf{Comunidad universitaria (8,500 personas):} Acceso público a información deportiva, fomento del espíritu universitario.
    \item \textbf{Autoridades institucionales:} Métricas y analítica para evaluación de políticas deportivas y asignación de presupuestos.
\end{itemize}

\subsection{Alcance del Producto}

El Sistema de Gestión de Eventos Deportivos (SIGED) es una plataforma integral de software diseñada para automatizar, optimizar y digitalizar la gestión completa del ciclo de vida de competencias deportivas interuniversitarias. El Sistema abarca desde la fase de planificación estratégica de torneos hasta la generación de analítica avanzada post-evento.

El producto contempla 16 módulos funcionales principales, arquitecturados de manera modular y escalable, que proporcionan capacidades para:

\begin{enumerate}[noitemsep]
    \item Gestión de identidad, acceso y seguridad
    \item Administración de usuarios con control basado en roles (RBAC)
    \item Gestión de información institucional multi-facultad
    \item Administración del catálogo de disciplinas deportivas
    \item Planificación y configuración de torneos competitivos
    \item Gestión de equipos participantes y plantillas de jugadores
    \item Administración de recursos físicos (escenarios) y programación de horarios
    \item Generación automatizada de fixtures mediante algoritmos configurables
    \item Ejecución digital de partidos con actas electrónicas
    \item Cálculo automático de resultados, clasificaciones y tablas de posiciones
    \item Gestión de procesos disciplinarios y justicia deportiva
    \item Sistema de notificaciones multi-canal
    \item Portal de información pública para usuarios no autenticados
    \item Generación de reportes ejecutivos y exportación de datos
    \item Dashboard de analítica con indicadores clave de desempeño (KPIs)
    \item Configuración general del sistema y personalización institucional
\end{enumerate}

\subsection{Definiciones, Acrónimos y Abreviaturas}

\subsubsection{Definiciones}

\begin{itemize}[noitemsep]
    \item \textbf{Usuario Autenticado:} Persona que ha completado exitosamente el proceso de autenticación y posee una sesión activa válida en el Sistema.
    \item \textbf{Rol:} Conjunto de permisos y capacidades asignadas a un tipo específico de usuario dentro del Sistema.
    \item \textbf{Fixture:} Calendario completo de partidos programados para un torneo, generado algorítmicamente considerando restricciones de equipos, escenarios y horarios.
    \item \textbf{Acta Digital:} Documento electrónico que registra formalmente todos los eventos y datos relevantes de un partido deportivo.
    \item \textbf{Stakeholder:} Cualquier individuo, grupo u organización que puede afectar, ser afectado o percibirse afectado por el desarrollo o implementación del Sistema.
\end{itemize}

\subsubsection{Acrónimos y Abreviaturas}

\begin{itemize}[noitemsep]
    \item \textbf{SIGED:} Sistema de Gestión de Eventos Deportivos (denominación del producto).
    \item \textbf{ERS:} Especificación de Requisitos de Software.
    \item \textbf{RF:} Requisito Funcional.
    \item \textbf{RNF:} Requisito No Funcional.
    \item \textbf{RP:} Requisito de Proyecto.
    \item \textbf{CRUD:} Create, Read, Update, Delete (Operaciones básicas de persistencia).
    \item \textbf{API:} Application Programming Interface (Interfaz de Programación de Aplicaciones).
    \item \textbf{REST:} Representational State Transfer (Estilo arquitectónico de servicios web).
    \item \textbf{JWT:} JSON Web Token (Estándar de autenticación basado en tokens).
    \item \textbf{RBAC:} Role-Based Access Control (Control de acceso basado en roles).
    \item \textbf{KPI:} Key Performance Indicator (Indicador Clave de Desempeño).
    \item \textbf{UI:} User Interface (Interfaz de Usuario).
    \item \textbf{UX:} User Experience (Experiencia de Usuario).
    \item \textbf{SLA:} Service Level Agreement (Acuerdo de Nivel de Servicio).
    \item \textbf{HTTPS:} HyperText Transfer Protocol Secure (Protocolo seguro de transferencia de hipertexto).
    \item \textbf{SQL:} Structured Query Language (Lenguaje de Consulta Estructurado).
    \item \textbf{WCAG:} Web Content Accessibility Guidelines (Pautas de Accesibilidad de Contenido Web).
\end{itemize}


\subsection{Visión General del Documento}

Este documento se estructura en las siguientes secciones principales:

\begin{itemize}[noitemsep]
    \item \textbf{Sección 1 - Introducción:} Presenta el propósito, alcance, audiencia y terminología del documento.
    \item \textbf{Sección 2 - Requisitos del Proyecto:} Define aspectos de gestión, recursos, plazos y restricciones del proyecto de desarrollo.
    \item \textbf{Sección 3 - Descripción General del Producto:} Contextualiza el producto, sus usuarios y restricciones generales.
    \item \textbf{Sección 4 - Requisitos Funcionales del Producto:} Especifica todas las funcionalidades y comportamientos esperados del Sistema.
    \item \textbf{Sección 5 - Requisitos No Funcionales del Producto:} Define atributos de calidad, rendimiento, seguridad y otros aspectos transversales.
    \item \textbf{Sección 6 - Matriz de Trazabilidad:} Establece vínculos entre requisitos, casos de uso, componentes y pruebas.
    \item \textbf{Sección 7 - Apéndices:} Información complementaria, diagramas adicionales y documentación de soporte.
\end{itemize}

\newpage

\section{Requisitos del Proyecto}

Esta sección establece los requisitos relacionados con la gestión del proyecto de desarrollo del Sistema, incluyendo aspectos de planificación, recursos, cronograma, presupuesto, metodología y gestión de riesgos.

\subsection{RP-1. Objetivos del Proyecto}

\subsubsection{RP-1.1 Objetivo General}

Desarrollar e implementar un Sistema de Gestión de Eventos Deportivos (SIGED) completo, escalable y de alta disponibilidad que permita automatizar y optimizar la gestión integral de torneos deportivos interuniversitarios, cumpliendo con estándares internacionales de calidad de software y mejores prácticas de ingeniería.

\subsubsection{RP-1.2 Objetivos Específicos}

\begin{enumerate}[noitemsep]
    \item Reducir en un 70\% el tiempo invertido en tareas administrativas relacionadas con gestión de torneos.
    \item Eliminar el uso de documentación física, digitalizando el 100\% de los procesos mediante actas electrónicas.
    \item Proporcionar acceso público en tiempo real a fixtures, resultados y estadísticas para usuarios externos.
    \item Garantizar disponibilidad del sistema del 99.5\% durante períodos de competencia activa.
    \item Implementar mecanismos de seguridad conforme a estándares ISO 27001 para protección de datos sensibles.
    \item Generar capacidades de analítica avanzada con al menos 15 KPIs deportivos relevantes.
\end{enumerate}

\subsection{RP-2. Alcance del Proyecto}

\subsubsection{RP-2.1 Entregables del Proyecto}

\begin{enumerate}[noitemsep]
    \item \textbf{Documentación de Requisitos:} ERS completa (este documento).
    \item \textbf{Diseño Arquitectónico:} Documento de arquitectura de software (SAD).
    \item \textbf{Diseño Detallado:} Diagramas UML (casos de uso, secuencia, clases, componentes).
    \item \textbf{Código Fuente:} Repositorio completo con código documentado y versionado.
    \item \textbf{Base de Datos:} Scripts de creación, migración y población de datos de prueba.
    \item \textbf{Aplicación Web:} Sistema completamente funcional desplegado en ambiente de producción.
    \item \textbf{Manuales de Usuario:} Documentación para cada perfil de usuario (Admin, Comité, Delegado).
    \item \textbf{Manual Técnico:} Guía de instalación, configuración y mantenimiento del sistema.
    \item \textbf{Plan de Pruebas:} Especificación de casos de prueba y resultados de ejecución.
    \item \textbf{Informe de Seguridad:} Análisis de vulnerabilidades y medidas implementadas.
    \item \textbf{Material de Capacitación:} Videos tutoriales y presentaciones para formación de usuarios.
\end{enumerate}

\subsubsection{RP-2.2 Exclusiones del Proyecto}

El proyecto NO incluye:

\begin{itemize}[noitemsep]
    \item Desarrollo de aplicaciones móviles nativas (iOS/Android).
    \item Integración con sistemas de transmisión de video en vivo.
    \item Integración con plataformas de pago electrónico para venta de entradas.
    \item Gestión de recursos humanos (arbitraje, personal médico) más allá del registro básico.
    \item Sistema de reservas para uso recreativo de escenarios deportivos.
\end{itemize}

\subsection{RP-3. Recursos del Proyecto}

\subsubsection{RP-3.1 Recursos Humanos}

El equipo del proyecto está conformado por 6 integrantes que operan bajo modalidad de "Aporte Académico" como parte de su proyecto de tesis. La siguiente tabla detalla los roles, cargas horarias y valorización del esfuerzo según tarifas de mercado peruano para proyectos similares.

\begin{longtable}{|>{\raggedright\arraybackslash}p{3.2cm}|>{\centering\arraybackslash}p{1cm}|>{\centering\arraybackslash}p{1.8cm}|>{\centering\arraybackslash}p{2cm}|>{\centering\arraybackslash}p{2cm}|>{\raggedright\arraybackslash}p{4cm}|}
\hline
\rowcolor{headercolor}
\textcolor{white}{\textbf{Rol}} & \textcolor{white}{\textbf{Cant.}} & \textcolor{white}{\textbf{Horas Total}} & \textcolor{white}{\textbf{Tarifa/\newline Hora}} & \textcolor{white}{\textbf{Costo Total}} & \textcolor{white}{\textbf{Responsabilidades}} \\
\hline
\endfirsthead
\hline
\rowcolor{headercolor}
\textcolor{white}{\textbf{Rol}} & \textcolor{white}{\textbf{Cant.}} & \textcolor{white}{\textbf{Horas Total}} & \textcolor{white}{\textbf{Tarifa/\newline Hora}} & \textcolor{white}{\textbf{Costo Total}} & \textcolor{white}{\textbf{Responsabilidades}} \\
\hline
\endhead
\hline
\endlastfoot

Gerente de Proyecto / Arquitecto & 1 & 240 hrs & S/. 50.00 & S/. 12,000 & Planificación, control, arquitectura, decisiones técnicas \\
\hline

Desarrollador Full-Stack Senior & 1 & 300 hrs & S/. 40.00 & S/. 12,000 & Backend .NET, diseño DB, APIs REST \\
\hline
Desarrollador Full-Stack & 2 & 420 hrs c/u & S/. 30.00 & S/. 25,200 & Frontend React, backend, integración \\
\hline

Analista / QA & 1 & 120 hrs & S/. 30.00 & S/. 3,600 & Análisis de requisitos, pruebas \\
\hline
Diseñador UI/UX & 1 & 80 hrs & S/. 35.00 & S/. 2,800 & Prototipos, diseño de interfaz \\
\hline

\textbf{SUBTOTAL} & \textbf{6} & \textbf{1,580 hrs} & - & \textbf{S/. 55,600} & \textbf{Aporte valorizado del equipo} \\
\hline
\caption{Recursos Humanos - Valorización del Equipo}
\end{longtable}

\textbf{Nota:} Estas tarifas corresponden a valores de mercado para profesionales junior-mid en Perú (2024-2025). El equipo NO recibirá remuneración monetaria; los costos representan el aporte valorizado académico.

\subsubsection{RP-3.2 Recursos Tecnológicos e Infraestructura}

\textbf{A. Software y Licencias}

\begin{longtable}{|>{\raggedright\arraybackslash}p{4.5cm}|>{\centering\arraybackslash}p{3cm}|>{\centering\arraybackslash}p{2.5cm}|>{\centering\arraybackslash}p{2cm}|>{\raggedright\arraybackslash}p{2cm}|}
\hline
\rowcolor{headercolor}
\textcolor{white}{\textbf{Herramienta / Servicio}} & \textcolor{white}{\textbf{Tipo}} & \textcolor{white}{\textbf{Licencia}} & \textcolor{white}{\textbf{Costo Mensual}} & \textcolor{white}{\textbf{Costo Total (6 meses)}} \\
\hline
\endfirsthead
\hline
\rowcolor{headercolor}
\textcolor{white}{\textbf{Herramienta / Servicio}} & \textcolor{white}{\textbf{Tipo}} & \textcolor{white}{\textbf{Licencia}} & \textcolor{white}{\textbf{Costo Mensual}} & \textcolor{white}{\textbf{Costo Total (6 messes)}} \\
\hline
\endhead
\hline
\endlastfoot

Visual Studio Code & IDE & Gratuita & S/. 0 & S/. 0 \\
\hline
\rowcolor{rowgray}
.NET 8 SDK & Framework & Open Source & S/. 0 & S/. 0 \\
\hline
Node.js + React & Framework & Open Source & S/. 0 & S/. 0 \\
\hline
\rowcolor{rowgray}
PostgreSQL & Base de Datos & Open Source & S/. 0 & S/. 0 \\
\hline
Git + GitHub & Control Versiones & Estudiante (gratuito) & S/. 0 & S/. 0 \\
\hline
\rowcolor{rowgray}
Docker Desktop & Contenedores & Gratuita uso académico & S/. 0 & S/. 0 \\
\hline
Postman & Testing API & Plan gratuito & S/. 0 & S/. 0 \\
\hline
\rowcolor{rowgray}
Figma & Diseño UI/UX & Plan gratuito & S/. 0 & S/. 0 \\
\hline
\textbf{SUBTOTAL Software} & & & & \textbf{S/. 0} \\
\hline
\caption{Software y Herramientas de Desarrollo}
\end{longtable}

\textbf{B. Infraestructura Cloud}

\begin{longtable}{|>{\raggedright\arraybackslash}p{4cm}|>{\centering\arraybackslash}p{2.5cm}|>{\centering\arraybackslash}p{3cm}|>{\centering\arraybackslash}p{2cm}|>{\raggedright\arraybackslash}p{2cm}|}
\hline
\rowcolor{headercolor}
\textcolor{white}{\textbf{Servicio}} & \textcolor{white}{\textbf{Proveedor}} & \textcolor{white}{\textbf{Especificación}} & \textcolor{white}{\textbf{Costo Mensual}} & \textcolor{white}{\textbf{Costo Total}} \\
\hline
\endfirsthead
\hline
\rowcolor{headercolor}
\textcolor{white}{\textbf{Servicio}} & \textcolor{white}{\textbf{Proveedor}} & \textcolor{white}{\textbf{Especificación}} & \textcolor{white}{\textbf{Costo Mensual}} & \textcolor{white}{\textbf{Costo Total (6 meses)}} \\
\hline
\endhead
\hline
\endlastfoot

\multicolumn{5}{|l|}{\cellcolor{columnblue}\textbf{Opción 1: Infraestructura Cloud Moderna}} \\
\hline
Frontend & Netlify/Vercel & 100GB bandwidth/mes & Incluido & Incluido \\
\hline
\rowcolor{rowgray}
Base de Datos & PostgreSQL 15 & 10GB (hosting cloud) & Incluido & Incluido \\
\hline
Storage & Cloud Storage & 100GB archivos & Incluido & Incluido \\
\hline
\rowcolor{rowgray}
Certificado SSL & Let's Encrypt & TLS automático & Incluido & Incluido \\
\hline
Monitoreo & UptimeRobot & 50 monitores & Incluido & Incluido \\
\hline
\rowcolor{rowgray}
CDN & Cloudflare & Protección DDoS & Incluido & Incluido \\
\hline

\multicolumn{5}{|l|}{\cellcolor{columnblue}\textbf{Opción 2: VPS Nacional}} \\
\hline
VPS Perú & HostDime/Axarnet & 4 vCPU, 16GB RAM, 500GB SSD & S/. 150 & S/. 900 \\
\hline

\multicolumn{5}{|l|}{\cellcolor{columnblue}\textbf{Opción 3: Infraestructura Universitaria}} \\
\hline
App Service & Azure for Students & B2 (2 cores, 4GB RAM) & Créditos Azure & Créditos Azure \\
\hline
\rowcolor{rowgray}
PostgreSQL & Azure for Students & Basic (2 vCore, 20GB) & Créditos Azure & Créditos Azure \\
\hline

\rowcolor{rowgray}
\textbf{SUBTOTAL Cloud} & & \textbf{Variable según opción} & & \textbf{S/. 0 - 900} \\
\hline
\caption{Infraestructura Cloud - Múltiples Opciones}
\end{longtable}

\textbf{Consideraciones sobre Infraestructura:} 
\begin{itemize}[noitemsep]
    \item \textbf{Opción 1:} Utiliza servicios cloud modernos con planes educativos. Ideal para comenzar y escalar progresivamente según la demanda real del sistema
    \item \textbf{Opción 2:} VPS nacional con mejor control y latencia local. Útil si se requiere personalización avanzada o integración con sistemas internos de la universidad
    \item \textbf{Opción 3:} Aprovecha beneficios para estudiantes de Azure. La universidad cuenta con créditos renovables anuales vía GitHub Education
    \item La elección final dependerá de la proyección de uso, políticas institucionales y preferencias de la OTI de la UNAS
\end{itemize}

\textbf{C. Hardware}

\begin{longtable}{|>{\raggedright\arraybackslash}p{4cm}|>{\centering\arraybackslash}p{2.5cm}|>{\centering\arraybackslash}p{2cm}|>{\centering\arraybackslash}p{2cm}|>{\raggedright\arraybackslash}p{3.5cm}|}
\hline
\rowcolor{headercolor}
\textcolor{white}{\textbf{Equipo}} & \textcolor{white}{\textbf{Cantidad}} & \textcolor{white}{\textbf{Costo Unit.}} & \textcolor{white}{\textbf{Costo Total}} & \textcolor{white}{\textbf{Observación}} \\
\hline
\endfirsthead
\hline
\rowcolor{headercolor}
\textcolor{white}{\textbf{Equipo}} & \textcolor{white}{\textbf{Cantidad}} & \textcolor{white}{\textbf{Costo Unit.}} & \textcolor{white}{\textbf{Costo Total}} & \textcolor{white}{\textbf{Observación}} \\
\hline
\endhead
\hline
\endlastfoot

Laptops de Desarrollo & 6 & - & - & Equipo propio de cada integrante \\
\hline
\rowcolor{rowgray}
Desgaste/Depreciación Laptops & 6 & S/. 300 & S/. 1,800 & Valorización 6 meses uso intensivo \\
\hline
Tablets para Pruebas Móviles & 2 & - & - & Dispositivos personales del equipo \\
\hline
\rowcolor{rowgray}
\textbf{SUBTOTAL Hardware} & & & \textbf{S/. 1,800} & \textbf{Costo valorizado} \\
\hline
\caption{Hardware - Equipos de Desarrollo}
\end{longtable}

\subsubsection{RP-3.3 Otros Gastos Operativos}

\begin{longtable}{|>{\raggedright\arraybackslash}p{5cm}|>{\centering\arraybackslash}p{2cm}|>{\centering\arraybackslash}p{2cm}|>{\raggedright\arraybackslash}p{5cm}|}
\hline
\rowcolor{headercolor}
\textcolor{white}{\textbf{Concepto}} & \textcolor{white}{\textbf{Costo Unit.}} & \textcolor{white}{\textbf{Costo Total}} & \textcolor{white}{\textbf{Descripción}} \\
\hline
\endfirsthead
\hline
\rowcolor{headercolor}
\textcolor{white}{\textbf{Concepto}} & \textcolor{white}{\textbf{Costo Unit.}} & \textcolor{white}{\textbf{Costo Total}} & \textcolor{white}{\textbf{Descripción}} \\
\hline
\endhead
\hline
\endlastfoot

Internet / Conectividad & S/. 80/mes & S/. 480 & 6 meses para 6 integrantes (valorizado) \\
\hline
\rowcolor{rowgray}
Energía Eléctrica & S/. 60/mes & S/. 360 & Consumo laptops 6 meses (valorizado) \\
\hline
Impresiones / Documentación & S/. 20 & S/. 120 & Documentos técnicos, entregables \\
\hline
\rowcolor{rowgray}
Movilidad / Reuniones & S/. 15 & S/. 90 & Transporte para sesiones presenciales \\
\hline
Materiales de Capacitación & S/. 10 & S/. 60 & Manuales impresos para usuarios \\
\hline
\rowcolor{rowgray}
Contingencia (Misc.) & - & S/. 100 & Gastos imprevistos menores \\
\hline
\textbf{SUBTOTAL Otros Gastos} & & \textbf{S/. 1,210} & \\
\hline
\caption{Gastos Operativos del Proyecto}
\end{longtable}

\subsubsection{RP-3.4 Presupuesto Consolidado del Proyecto}

\begin{longtable}{|>{\raggedright\arraybackslash}p{7cm}|>{\centering\arraybackslash}p{4cm}|>{\centering\arraybackslash}p{3.5cm}|}
\hline
\rowcolor{headerblue}
\textcolor{white}{\textbf{Categoría}} & \textcolor{white}{\textbf{Costo (S/.)}} & \textcolor{white}{\textbf{Porcentaje}} \\
\hline
\endfirsthead
\hline
\rowcolor{headerblue}
\textcolor{white}{\textbf{Categoría}} & \textcolor{white}{\textbf{Costo (S/.)}} & \textcolor{white}{\textbf{Porcentaje}} \\
\hline
\endhead
\hline
\endlastfoot

\textbf{1. Recursos Humanos (valorizado)} & S/. 55,600 & 94.5\% \\
\hline
\rowcolor{rowgray}
\textbf{2. Software y Licencias} & S/. 0 & 0.0\% \\
\hline
\textbf{3. Infraestructura Cloud} & S/. 0 & 0.0\% \\
\hline
\rowcolor{rowgray}
\textbf{4. Hardware (depreciación)} & S/. 1,800 & 3.1\% \\
\hline
\textbf{5. Otros Gastos Operativos} & S/. 1,210 & 2.1\% \\
\hline
\rowcolor{headergray}
\textbf{SUBTOTAL} & \textbf{S/. 58,610} & \textbf{99.7\%} \\
\hline
\textbf{Reserva de Contingencia (5\%)} & S/. 2,930 & 0.3\% \\
\hline
\rowcolor{headerblue}
\textcolor{white}{\textbf{TOTAL GENERAL DEL PROYECTO}} & \textcolor{white}{\textbf{S/. 61,540}} & \textcolor{white}{\textbf{100\%}} \\
\hline
\caption{Presupuesto Consolidado del Proyecto SIGDE}
\end{longtable}

\textbf{Fuente de Financiamiento:}
\begin{itemize}[noitemsep]
    \item \textbf{Aporte Valorizado (equipo):} S/. 58,400 (94.9\%) - Horas de desarrollo, equipos, conectividad
    \item \textbf{Aporte Monetario (autofinanciado):} S/. 210 (0.3\%) - Impresiones, movilidad, materiales
    \item \textbf{Aporte Institucional (UNAS):} S/. 0 (0\%) - Infraestructura cloud mediante convenios
    \item \textbf{Reserva de Contingencia:} S/. 2,930 (4.8\%) - Para imprevistos y cambios de alcance
\end{itemize}

\textbf{Notas Importantes:}
\begin{itemize}[noitemsep]
    \item El proyecto es autofinanciado por el equipo como requisito de grado (modalidad tesis).
    \item Las tarifas de recursos humanos reflejan valores de mercado peruano para desarrolladores junior-mid (no son pagos reales).
    \item La infraestructura cloud es gratuita mediante Azure for Students (GitHub Education Pack).
    \item Todo el software utilizado es open source o cuenta con licencia educativa gratuita.
    \item El único gasto monetario real es de S/. 210 para operaciones logísticas.
\end{itemize}

\subsection{RP-4. Cronograma y Fases del Proyecto}

El cronograma del proyecto ha sido estructurado para ejecutarse en un horizonte temporal de \textbf{24 semanas (6 meses)}. Esta duración permite no solo el desarrollo del software, sino también una fase crítica de estabilización y marcha blanca (piloto) en el entorno de producción, asegurando la transferencia operativa a la OTI.

La estrategia de gestión del tiempo utiliza el enfoque de \textbf{Planificación Gradual} (Rolling Wave Planning). Las fases de Análisis y Diseño se detallan a nivel de paquetes de trabajo, mientras que la Construcción y el Despliegue se gestionan por hitos. Se han establecido \textbf{Puertas de Fase} (Phase Gates) mensuales para la validación formal de entregables por el Comité de Interesados (Sponsor, OTI y Departamento de Deportes).

\subsubsection{RP-4.1 Hitos de Alto Nivel del Proyecto}

\begin{longtable}{|>{\centering\arraybackslash}p{2cm}|>{\raggedright\arraybackslash}p{9cm}|>{\centering\arraybackslash}p{3cm}|}
\hline
\rowcolor{headerblue}
\textcolor{white}{\textbf{Hito}} & \textcolor{white}{\textbf{Entregable Clave}} & \textcolor{white}{\textbf{Fecha Estimada}} \\
\hline
\endfirsthead
\hline
\rowcolor{headerblue}
\textcolor{white}{\textbf{Hito}} & \textcolor{white}{\textbf{Entregable Clave}} & \textcolor{white}{\textbf{Fecha Estimada}} \\
\hline
\endhead
\hline
\endlastfoot

Hito 1 & Acta de Constitución y Alcance Aprobados & Mes 1 \\
\hline
\rowcolor{rowgray}
Hito 2 & Arquitectura y Prototipos UI Validados & Mes 2 \\
\hline
Hito 3 & Backend Core y Motor de Reglas Funcional & Mes 3 \\
\hline
\rowcolor{rowgray}
Hito 4 & Integración Frontend + Tiempo Real (SignalR) & Mes 4 \\
\hline
Hito 5 & Despliegue en Servidores OTI (Producción) & Mes 5 \\
\hline
\rowcolor{rowgray}
Hito 6 & Fin de Marcha Blanca y Cierre del Proyecto & Mes 6 \\
\hline
\caption{Hitos de Alto Nivel - SIGED}
\end{longtable}

\subsubsection{RP-4.2 Cronograma Detallado por Fases}

\begin{longtable}{|>{\centering\arraybackslash}p{1cm}|>{\raggedright\arraybackslash}p{6cm}|>{\centering\arraybackslash}p{2cm}|>{\centering\arraybackslash}p{2cm}|>{\raggedright\arraybackslash}p{3cm}|}
\hline
\rowcolor{headerblue}
\textcolor{white}{\textbf{ID}} & \textcolor{white}{\textbf{Fase / Actividad}} & \textcolor{white}{\textbf{Duración}} & \textcolor{white}{\textbf{Semanas}} & \textcolor{white}{\textbf{Responsable}} \\
\hline
\endfirsthead
\hline
\rowcolor{headerblue}
\textcolor{white}{\textbf{ID}} & \textcolor{white}{\textbf{Fase / Actividad}} & \textcolor{white}{\textbf{Duración}} & \textcolor{white}{\textbf{Semanas}} & \textcolor{white}{\textbf{Responsable}} \\
\hline
\endhead
\hline
\endfoot
\hline
\caption{Cronograma Detallado del Proyecto SIGED (24 semanas)}
\endlastfoot

\multicolumn{5}{|l|}{\cellcolor{headerblue}\textcolor{white}{\textbf{1.0 GESTIÓN E INICIO}}} \\
\hline
1.0 & \textbf{GESTIÓN E INICIO} & \textbf{3 Sem} & \textbf{1-3} & \textbf{Project Manager} \\
\hline
\rowcolor{rowgray}
1.1 & Reuniones con Interesados (Deportes/OTI) & 1 sem & Sem 1 & PM \\
\hline
1.2 & Elaboración y Firma del Acta de Constitución & 1 sem & Sem 2 & PM \\
\hline
\rowcolor{rowgray}
1.3 & Configuración de Herramientas y Repositorios & 1 sem & Sem 3 & Tech Lead \\
\hline

\multicolumn{5}{|l|}{\cellcolor{headerblue}\textcolor{white}{\textbf{2.0 PLANIFICACIÓN}}} \\
\hline
2.0 & \textbf{PLANIFICACIÓN} & \textbf{2 Sem} & \textbf{3-4} & \textbf{PM / Equipo} \\
\hline
\rowcolor{rowgray}
2.1 & Definición del Alcance, EDT y Matriz de Riesgos & 1 sem & Sem 3 & PM \\
\hline
2.2 & Elaboración del Cronograma y Plan de Comunicaciones & 1 sem & Sem 4 & PM \\
\hline

\multicolumn{5}{|l|}{\cellcolor{headerblue}\textcolor{white}{\textbf{3.0 ANÁLISIS Y DISEÑO (INGENIERÍA)}}} \\
\hline
3.0 & \textbf{ANÁLISIS Y DISEÑO (INGENIERÍA)} & \textbf{4 Sem} & \textbf{5-8} & \textbf{Arquitecto / UX} \\
\hline
\rowcolor{rowgray}
3.1 & Especificación de Requisitos (SRS) y Reglas de Negocio & 2 sem & Sem 5-6 & Analista \\
\hline
3.2 & Prototipos UX/UI (Figma) - Admin y Móvil & 2 sem & Sem 6-7 & Diseñador \\
\hline
\rowcolor{rowgray}
3.3 & Arquitectura C4 y Modelo de Base de Datos & 1 sem & Sem 8 & Arquitecto \\
\hline

\multicolumn{5}{|l|}{\cellcolor{headerblue}\textcolor{white}{\textbf{4.0 CONSTRUCCIÓN (DESARROLLO)}}} \\
\hline
4.0 & \textbf{CONSTRUCCIÓN (DESARROLLO)} & \textbf{12 Sem} & \textbf{9-20} & \textbf{Dev Team} \\
\hline
\rowcolor{rowgray}
4.1 & \textbf{Backend (.NET 8)} &  &  &  \\
\hline
4.1.1 & Configuración Clean Arch + Identity (Auth/Roles) & 1 sem & Sem 9 & Tech Lead \\
\hline
\rowcolor{rowgray}
4.1.2 & API Configuración y Torneos (Core Multi-evento) & 3 sem & Sem 10-12 & Backend Dev \\
\hline
4.1.3 & API Gestión de Reservas (Algoritmo de Horarios) & 3 sem & Sem 13-15 & Backend Dev \\
\hline
\rowcolor{rowgray}
4.1.4 & Motor de Tiempo Real (SignalR \& WebSockets) & 2 sem & Sem 16-17 & Backend Dev \\
\hline
4.2 & \textbf{Frontend (React / Next.js)} &  &  &  \\
\hline
\rowcolor{rowgray}
4.2.1 & Maquetación Panel Admin y Gestión de Usuarios & 2 sem & Sem 10-11 & Frontend Dev \\
\hline
4.2.2 & Módulo de Reservas y Calendario Interactivo & 3 sem & Sem 12-14 & Frontend Dev \\
\hline
\rowcolor{rowgray}
4.2.3 & Portal Público (Integración con SignalR) & 3 sem & Sem 15-17 & Frontend Dev \\
\hline
4.2.4 & App Móvil Mesa de Control (Acta Digital PWA) & 3 sem & Sem 17-19 & Frontend Dev \\
\hline
\rowcolor{rowgray}
4.3 & \textbf{Integración} &  &  &  \\
\hline
4.3.1 & Integración Final API + Clientes + Base de Datos & 1 sem & Sem 20 & Fullstack \\
\hline

\multicolumn{5}{|l|}{\cellcolor{headerblue}\textcolor{white}{\textbf{5.0 ASEGURAMIENTO DE CALIDAD (QA)}}} \\
\hline
5.0 & \textbf{ASEGURAMIENTO DE CALIDAD (QA)} & \textbf{3 Sem} & \textbf{19-21} & \textbf{QA Tester} \\
\hline
\rowcolor{rowgray}
5.1 & Pruebas Unitarias (Backend) & Continuo & Sem 9-20 & Developers \\
\hline
5.2 & Pruebas de Integración y Flujos de Usuario & 1 sem & Sem 19 & QA \\
\hline
\rowcolor{rowgray}
5.3 & Pruebas de Carga (Stress Testing - 500 usuarios) & 1 sem & Sem 20 & QA / OTI \\
\hline
5.4 & Corrección de Bugs (Sprint de Estabilización) & 1 sem & Sem 21 & Dev Team \\
\hline

\multicolumn{5}{|l|}{\cellcolor{headerblue}\textcolor{white}{\textbf{6.0 DESPLIEGUE EN OTI (PRODUCCIÓN)}}} \\
\hline
6.0 & \textbf{DESPLIEGUE EN OTI (PRODUCCIÓN)} & \textbf{2 Sem} & \textbf{21-22} & \textbf{DevOps / OTI} \\
\hline
\rowcolor{rowgray}
6.1 & Configuración de Servidor Linux y Docker & 3 días & Sem 21 & DevOps \\
\hline
6.2 & Despliegue de Contenedores y Base de Datos & 2 días & Sem 21 & DevOps \\
\hline
\rowcolor{rowgray}
6.3 & Configuración de Dominio, SSL y Firewall & 1 sem & Sem 22 & OTI \\
\hline

\multicolumn{5}{|l|}{\cellcolor{headerblue}\textcolor{white}{\textbf{7.0 MARCHA BLANCA Y PILOTO}}} \\
\hline
7.0 & \textbf{MARCHA BLANCA Y PILOTO} & \textbf{2 Sem} & \textbf{22-23} & \textbf{PM / Deportes} \\
\hline
\rowcolor{rowgray}
7.1 & Capacitación a Usuarios (Jueces y Admin) & 3 días & Sem 22 & PM \\
\hline
7.2 & Ejecución de Torneo Piloto (Simulacro Real) & 1 sem & Sem 23 & Todos \\
\hline
\rowcolor{rowgray}
7.3 & Monitoreo de Logs y Ajustes Finales & 2 días & Sem 23 & Dev Team \\
\hline

\multicolumn{5}{|l|}{\cellcolor{headerblue}\textcolor{white}{\textbf{8.0 CIERRE DEL PROYECTO}}} \\
\hline
8.0 & \textbf{CIERRE DEL PROYECTO} & \textbf{1 Sem} & \textbf{24} & \textbf{PM} \\
\hline
\rowcolor{rowgray}
8.1 & Entrega de Manuales (Usuario y Operaciones) & 2 días & Sem 24 & PM \\
\hline
8.2 & Acta de Cierre y Transferencia a OTI & 2 días & Sem 24 & PM \\
\hline

\end{longtable}

\textbf{Notas sobre el Cronograma:}
\begin{itemize}[noitemsep]
    \item El proyecto contempla 24 semanas efectivas de trabajo distribuidas en 6 meses calendario.
    \item Se considera un buffer implícito de 2-3 semanas para recesos académicos y feriados.
    \item Las actividades de pruebas unitarias (5.1) son continuas durante todo el desarrollo.
    \item Las fases de Backend y Frontend se ejecutan en paralelo a partir de la semana 10.
    \item Se establecen 6 hitos mensuales con revisiones formales del Comité de Interesados.
    \item La fase de Marcha Blanca incluye un torneo piloto real para validación integral del sistema.
\end{itemize}
\newpage
\subsubsection{RP-4.3 Matriz RACI de Responsabilidades}

La matriz RACI define claramente las responsabilidades y roles de cada miembro del equipo en las actividades clave del proyecto:

\begin{itemize}[noitemsep]
    \item \textbf{R (Responsible - Responsable):} Quien ejecuta la actividad.
    \item \textbf{A (Accountable - Aprobador):} Quien tiene la autoridad final y rinde cuentas.
    \item \textbf{C (Consulted - Consultado):} Quien proporciona información para la actividad.
    \item \textbf{I (Informed - Informado):} Quien debe ser notificado sobre el progreso/resultado.
\end{itemize}

\textbf{Roles del Equipo:}

\begin{table}[h]
\centering
\begin{tabular}{|>{\centering\arraybackslash}p{2cm}|>{\raggedright\arraybackslash}p{8cm}|}
\hline
\rowcolor{headerblue}
\textcolor{white}{\textbf{Código}} & \textcolor{white}{\textbf{Rol}} \\
\hline
GP & Gerente de Proyecto \\
\hline
\rowcolor{rowgray}
AS & Arquitecto de Software \\
\hline
AQ & Analista/QA \\
\hline
\rowcolor{rowgray}
DS & Desarrollador Full-Stack Senior \\
\hline
DEV & Desarrollador Full-Stack \\
\hline
\rowcolor{rowgray}
DIS & Diseñador UI/UX \\
\hline
DBA & Administrador de Base de Datos \\
\hline
\rowcolor{rowgray}
SEG & Especialista en Seguridad \\
\hline
DOC & Redactor Técnico / Documentación \\
\hline
\end{tabular}
\caption{Códigos de Roles - Matriz RACI}
\end{table}

\clearpage

\begin{landscape}

\begin{longtable}{|>{\raggedright\arraybackslash}p{5cm}|>{\centering\arraybackslash}p{1cm}|>{\centering\arraybackslash}p{1cm}|>{\centering\arraybackslash}p{1cm}|>{\centering\arraybackslash}p{1cm}|>{\centering\arraybackslash}p{1cm}|>{\centering\arraybackslash}p{1cm}|>{\centering\arraybackslash}p{1cm}|>{\centering\arraybackslash}p{1cm}|>{\centering\arraybackslash}p{1cm}|>{\raggedright\arraybackslash}p{4cm}|}
\hline
\rowcolor{headerblue}
\textcolor{white}{\textbf{Actividad}} & \textcolor{white}{\textbf{GP}} & \textcolor{white}{\textbf{AS}} & \textcolor{white}{\textbf{AQ}} & \textcolor{white}{\textbf{DS}} & \textcolor{white}{\textbf{DEV}} & \textcolor{white}{\textbf{DIS}} & \textcolor{white}{\textbf{DBA}} & \textcolor{white}{\textbf{SEG}} & \textcolor{white}{\textbf{DOC}} & \textcolor{white}{\textbf{Observaciones}} \\
\hline
\endfirsthead


\hline
\rowcolor{headerblue}
\textcolor{white}{\textbf{Actividad}} & \textcolor{white}{\textbf{GP}} & \textcolor{white}{\textbf{AS}} & \textcolor{white}{\textbf{AQ}} & \textcolor{white}{\textbf{DS}} & \textcolor{white}{\textbf{DEV}} & \textcolor{white}{\textbf{DIS}} & \textcolor{white}{\textbf{DBA}} & \textcolor{white}{\textbf{SEG}} & \textcolor{white}{\textbf{DOC}} & \textcolor{white}{\textbf{Observaciones}} \\
\hline
\endhead

\hline
\endfoot

\hline
\caption{Matriz RACI - Responsabilidades del Proyecto SIGDE}
\endlastfoot

\multicolumn{11}{|l|}{\cellcolor{columnblue}\textbf{\textcolor{white}{FASE 1: INICIO Y PLANIFICACIÓN}}} \\
\hline

Acta de Constitución del Proyecto & A & C & I & I & I & I & I & I & C & Documento fundacional \\
\hline
\rowcolor{rowgray}
Plan de Gestión del Proyecto & A/R & C & C & I & I & I & I & I & C & Incluye alcance, tiempo, costos \\
\hline
Definición de Alcance (WBS) & A & R & C & C & I & I & I & I & C & Estructura EDT completa \\
\hline
\rowcolor{rowgray}
Plan de Comunicaciones & A/R & I & I & I & I & I & I & I & I & Stakeholders y canales \\
\hline

\multicolumn{11}{|l|}{\cellcolor{columnblue}\textbf{\textcolor{white}{FASE 2: ANÁLISIS DE REQUISITOS}}} \\
\hline

Levantamiento de Requisitos & A & C & R & I & I & C & C & C & C & Entrevistas y talleres \\
\hline
\rowcolor{rowgray}
Especificación de Requisitos (ERS) & A & C & R & C & I & I & C & C & R & Documento IEEE 830 \\
\hline
Casos de Uso del Sistema & I & C & R & C & I & I & I & I & C & Diagramas UML \\
\hline
\rowcolor{rowgray}
Historias de Usuario & A & I & R & R & I & C & I & I & I & Backlog priorizado \\
\hline
Prototipos de Interfaz & I & I & C & I & I & A/R & I & I & I & Wireframes y mockups \\
\hline
\rowcolor{rowgray}
Validación de Requisitos & A & C & R & C & I & C & I & I & I & Sesiones con stakeholders \\
\hline

\multicolumn{11}{|l|}{\cellcolor{columnblue}\textbf{\textcolor{white}{FASE 3: DISEÑO ARQUITECTÓNICO}}} \\
\hline

Arquitectura de Software (SAD) & A & A/R & I & C & I & I & C & C & R & Patrones y capas \\
\hline
\rowcolor{rowgray}
Diseño de Base de Datos & I & A & I & R & I & I & A/R & I & C & Modelo ER y normalización \\
\hline
Diseño de APIs REST & I & A & I & R & C & I & I & C & R & Contratos y documentación \\
\hline
\rowcolor{rowgray}
Estándares de Codificación & A & R & I & C & C & I & I & I & C & Guías y convenciones \\
\hline
Plan de Seguridad & A & C & C & C & I & I & I & A/R & C & Autenticación y autorización \\
\hline
\rowcolor{rowgray}
Diseño UI/UX Detallado & I & I & I & I & I & A/R & I & I & C & Sistema de diseño completo \\
\hline

\multicolumn{11}{|l|}{\cellcolor{columnblue}\textbf{\textcolor{white}{FASE 4: DESARROLLO - MÓDULOS CORE}}} \\
\hline

Módulo de Autenticación JWT & I & C & I & A/R & R & I & C & R & I & Login, registro, tokens \\
\hline
\rowcolor{rowgray}
Módulo de Gestión de Usuarios & I & C & C & A/R & R & I & R & C & I & CRUD y perfiles \\
\hline
Módulo de Instituciones & I & I & I & A/R & R & I & R & I & I & Facultades y escuelas \\
\hline
\rowcolor{rowgray}
Módulo de Disciplinas Deportivas & I & I & I & R & A/R & I & R & I & I & Catálogo y configuración \\
\hline
Componentes UI Reutilizables & I & C & I & R & R & A & I & I & I & Librería de componentes \\
\hline

\multicolumn{11}{|l|}{\cellcolor{columnblue}\textbf{\textcolor{white}{FASE 5: DESARROLLO - MÓDULOS TORNEOS}}} \\
\hline

\rowcolor{rowgray}
Módulo de Gestión de Torneos & I & C & C & A/R & R & I & R & I & I & Creación y configuración \\
\hline
Módulo de Equipos & I & I & I & A/R & R & I & R & I & I & Registro y convocatorias \\
\hline
\rowcolor{rowgray}
Módulo de Escenarios Deportivos & I & I & I & R & A/R & I & R & I & I & Infraestructura y reservas \\
\hline
Generador de Fixture & I & C & I & A/R & R & I & C & I & I & Algoritmos de emparejamiento \\
\hline
\rowcolor{rowgray}
Panel de Administración Torneos & I & C & I & R & R & A & I & C & I & Interfaz administrativa \\
\hline

\multicolumn{11}{|l|}{\cellcolor{columnblue}\textbf{\textcolor{white}{FASE 6: DESARROLLO - MÓDULOS COMPETENCIA}}} \\
\hline

Módulo de Gestión de Partidos & I & C & I & A/R & R & I & R & I & I & Programación y árbitros \\
\hline
\rowcolor{rowgray}
Módulo de Registro de Resultados & I & I & I & A/R & R & I & R & C & I & Actas y validaciones \\
\hline
Cálculo de Tablas de Posiciones & I & C & I & A/R & R & I & C & I & I & Puntos, goles, DIF \\
\hline
\rowcolor{rowgray}
Módulo de Justicia Deportiva & I & C & C & A/R & R & I & R & C & I & Sanciones y apelaciones \\
\hline
Sistema de Notificaciones & I & C & I & R & A/R & I & I & I & I & Email y push \\
\hline

\multicolumn{11}{|l|}{\cellcolor{columnblue}\textbf{\textcolor{white}{FASE 7: DESARROLLO - PORTAL Y REPORTES}}} \\
\hline

\rowcolor{rowgray}
Portal Público de Consulta & I & I & I & R & A/R & A & C & C & I & Landing y resultados en vivo \\
\hline
Dashboard Analítico & I & C & I & A/R & R & A & R & I & I & Gráficos y estadísticas \\
\hline
\rowcolor{rowgray}
Generador de Reportes & I & I & C & A/R & R & I & C & I & C & PDF y Excel exportables \\
\hline
Módulo de Configuración Global & I & C & I & A/R & R & I & C & C & I & Parámetros del sistema \\
\hline
\rowcolor{rowgray}
Módulo de Auditoría & I & C & C & A/R & R & I & C & A/R & I & Logs y trazabilidad \\
\hline

\multicolumn{11}{|l|}{\cellcolor{columnblue}\textbf{\textcolor{white}{FASE 8: PRUEBAS Y CALIDAD}}} \\
\hline

Plan de Pruebas & A & I & A/R & C & I & I & I & C & R & Estrategia integral de QA \\
\hline
\rowcolor{rowgray}
Pruebas Unitarias & I & I & C & R & A/R & I & I & I & I & Coverage mínimo 80\% \\
\hline
Pruebas de Integración & I & C & A/R & R & R & I & C & C & I & APIs y módulos \\
\hline
\rowcolor{rowgray}
Pruebas de Sistema (E2E) & A & I & A/R & C & C & C & I & C & I & Flujos completos \\
\hline
Pruebas de Usabilidad & A & I & C & I & I & A/R & I & I & R & Testing con usuarios reales \\
\hline
\rowcolor{rowgray}
Pruebas de Seguridad & A & C & C & R & I & I & I & A/R & R & Penetration testing OWASP \\
\hline
Pruebas de Rendimiento & A & C & A/R & R & I & I & R & I & R & Carga y estrés \\
\hline
\rowcolor{rowgray}
Corrección de Defectos Críticos & A & C & I & A/R & R & I & C & C & I & Bugs prioritarios \\
\hline

\multicolumn{11}{|l|}{\cellcolor{columnblue}\textbf{\textcolor{white}{FASE 9: DOCUMENTACIÓN}}} \\
\hline

Manual Técnico del Sistema & I & C & I & C & I & I & C & C & A/R & Arquitectura y APIs \\
\hline
\rowcolor{rowgray}
Manual de Usuario Final & A & I & C & I & I & C & I & I & A/R & Guías ilustradas \\
\hline
Manual de Administrador & I & C & C & C & I & I & C & C & A/R & Configuración y mantenimiento \\
\hline
\rowcolor{rowgray}
Material de Capacitación & A & I & C & I & I & C & I & I & A/R & Presentaciones y videos \\
\hline
Documentación de APIs & I & C & I & R & I & I & I & C & A/R & Swagger/OpenAPI \\
\hline

\multicolumn{11}{|l|}{\cellcolor{columnblue}\textbf{\textcolor{white}{FASE 10: DESPLIEGUE Y CIERRE}}} \\
\hline

\rowcolor{rowgray}
Configuración Servidor Producción & A & C & I & A/R & I & I & R & C & I & Azure App Service \\
\hline
Migración de Datos & A & C & C & R & I & I & A/R & I & I & Carga inicial \\
\hline
\rowcolor{rowgray}
Despliegue a Producción & A/R & C & C & R & I & I & C & C & I & Release definitivo \\
\hline
Capacitación Usuarios Finales & A/R & I & C & I & I & C & I & I & R & Talleres presenciales \\
\hline
\rowcolor{rowgray}
Monitoreo Post-Lanzamiento & A & C & C & R & I & I & R & C & I & Primeros 30 días \\
\hline
Soporte Inicial (Hypercare) & A & I & A/R & R & R & I & R & I & C & Respuesta rápida 24/7 \\
\hline
\rowcolor{rowgray}
Lecciones Aprendidas & A/R & C & C & C & C & C & C & C & R & Retrospectiva final \\
\hline
Cierre Formal del Proyecto & A/R & I & I & I & I & I & I & I & R & Acta de aceptación \\
\hline

\end{longtable}

\end{landscape}

\textbf{Notas sobre la Matriz RACI:}
\begin{itemize}[noitemsep]
    \item Cada actividad debe tener \textbf{un único Accountable (A)}, quien tiene la autoridad final.
    \item Puede haber \textbf{múltiples Responsables (R)}, quienes ejecutan la tarea.
    \item Los roles marcados como \textbf{A/R} indican que la misma persona es tanto responsable como aprobador (común en equipos pequeños).
    \item La matriz se revisará y actualizará al inicio de cada sprint para reflejar cambios organizacionales.
    \item En caso de conflictos de responsabilidad, el Gerente de Proyecto tiene la autoridad para resolver y reasignar.
\end{itemize}

\subsection{RP-5. Metodología de Desarrollo}

\subsubsection{RP-5.1 Marco Metodológico}

El proyecto adoptará metodología ágil basada en \textbf{Scrum}, con sprints de 3 semanas de duración. Se utilizarán prácticas complementarias de:

\begin{itemize}[noitemsep]
    \item \textbf{Kanban:} Para gestión visual del flujo de trabajo.
    \item \textbf{DevOps:} Integración continua y despliegue continuo (CI/CD).
    \item \textbf{Test-Driven Development (TDD):} Pruebas unitarias antes de implementación.
    \item \textbf{Pair Programming:} Para módulos críticos de seguridad y lógica compleja.
    \item \textbf{Code Review:} Revisión obligatoria por pares antes de merge a rama principal.
\end{itemize}

\subsubsection{RP-5.2 Ceremonias Ágiles}

\begin{itemize}[noitemsep]
    \item \textbf{Sprint Planning:} Primer día de cada sprint (4 horas).
    \item \textbf{Daily Stand-up:} Diario, 15 minutos.
    \item \textbf{Sprint Review:} Último día del sprint (2 horas).
    \item \textbf{Sprint Retrospective:} Último día del sprint (1.5 horas).
    \item \textbf{Backlog Refinement:} Mitad de sprint (2 horas).
\end{itemize}

\subsection{RP-6. Gestión de Riesgos}

\begin{longtable}{|>{\centering\arraybackslash}p{1.5cm}|>{\raggedright\arraybackslash}p{3.5cm}|>{\centering\arraybackslash}p{1.5cm}|>{\centering\arraybackslash}p{1.5cm}|>{\raggedright\arraybackslash}p{5.5cm}|}
\hline
\rowcolor{headercolor}
\textcolor{white}{\textbf{ID}} & \textcolor{white}{\textbf{Riesgo}} & \textcolor{white}{\textbf{Prob.}} & \textcolor{white}{\textbf{Imp.}} & \textcolor{white}{\textbf{Mitigación}} \\
\hline
\endfirsthead
\hline
\rowcolor{headercolor}
\textcolor{white}{\textbf{ID}} & \textcolor{white}{\textbf{Riesgo}} & \textcolor{white}{\textbf{Prob.}} & \textcolor{white}{\textbf{Imp.}} & \textcolor{white}{\textbf{Mitigación}} \\
\hline
\endhead
\hline
\endlastfoot

R-01 & Cambios frecuentes en requisitos & Alta & Alta & Metodología ágil con sprints cortos, validación continua con stakeholders. \\
\hline
\rowcolor{rowgray}
R-02 & Rotación de personal técnico & Media & Alta & Documentación exhaustiva, pair programming, transferencia de conocimiento. \\
\hline
R-03 & Problemas de rendimiento con carga concurrente & Media & Alta & Pruebas de carga tempranas, optimización proactiva, escalabilidad horizontal. \\
\hline
\rowcolor{rowgray}
R-04 & Vulnerabilidades de seguridad & Media & Crítica & Auditorías de seguridad periódicas, pruebas de penetración, cumplimiento OWASP. \\
\hline
R-05 & Retrasos en entrega de módulos & Media & Media & Buffer de tiempo del 15\%, priorización clara, seguimiento semanal. \\
\hline
\rowcolor{rowgray}
R-06 & Incompatibilidad con navegadores antiguos & Baja & Media & Definir navegadores soportados, testing cross-browser automatizado. \\
\hline
R-07 & Falta de disponibilidad de stakeholders clave & Media & Alta & Definir ventanas de revisión obligatorias, proxies de negocio. \\
\hline
\rowcolor{rowgray}
R-08 & Problemas con integración de terceros & Baja & Media & Abstraer dependencias externas, mocks para desarrollo y pruebas. \\
\hline
\caption{Matriz de Riesgos del Proyecto}
\end{longtable}

\subsection{RP-7. Presupuesto y Costos}

\subsubsection{RP-7.1 Estimación de Costos del Proyecto (6 meses de desarrollo)}

El presupuesto consolidado del proyecto se detalla en la sección RP-3.4 con un total de \textbf{S/. 61,540} mayormente como aporte valorizado del equipo de desarrollo.

\begin{itemize}[noitemsep]
    \item \textbf{Recursos Humanos (valorizado):} 94.5\% del presupuesto (S/. 55,600)
    \item \textbf{Hardware (depreciación):} 3.1\% (S/. 1,800)
    \item \textbf{Gastos operativos:} 2.1\% (S/. 1,210)
    \item \textbf{Infraestructura Cloud:} 0\% - Uso de servicios gratuitos
    \item \textbf{Software y Licencias:} 0\% - Todo software libre/open source
\end{itemize}

\subsubsection{RP-7.2 Costos de Operación Post-Implementación (Realidad UNAS)}

Una vez finalizado el desarrollo, la UNAS deberá asumir costos recurrentes de operación según la opción de despliegue elegida:

\textbf{OPCIÓN 1: Cloud Moderno (Netlify + PostgreSQL en la nube) - RECOMENDADO}

\begin{longtable}{|>{\raggedright\arraybackslash}p{5cm}|>{\centering\arraybackslash}p{3cm}|>{\raggedright\arraybackslash}p{6cm}|}
\hline
\rowcolor{headercolor}
\textcolor{white}{\textbf{Concepto}} & \textcolor{white}{\textbf{Costo Mensual}} & \textcolor{white}{\textbf{Observaciones}} \\
\hline
\endfirsthead
\hline
\rowcolor{headercolor}
\textcolor{white}{\textbf{Concepto}} & \textcolor{white}{\textbf{Costo Mensual}} & \textcolor{white}{\textbf{Observaciones}} \\
\hline
\endhead
\hline
\endlastfoot

Hosting Frontend (Netlify/Vercel) & Incluido & 100GB bandwidth/mes, SSL automático \\
\hline
\rowcolor{rowgray}
Base de Datos PostgreSQL & Incluido & 10GB, suficiente para 2-3 años de operación \\
\hline
Almacenamiento archivos & Incluido & 100GB storage para imágenes y documentos \\
\hline
\rowcolor{rowgray}
SSL/TLS Certificado & S/. 0 & Automático con Let's Encrypt \\
\hline
Monitoreo (UptimeRobot) & S/. 0 & Plan gratuito, 50 monitores \\
\hline
\rowcolor{rowgray}
\textbf{TOTAL MENSUAL} & \textbf{S/. 0} & \textbf{Sin costo recurrente} \\
\hline
\textbf{TOTAL ANUAL} & \textbf{S/. 0} & \\
\hline
\caption{Costos de Operación - Opción Gratuita}
\end{longtable}

\textbf{Limitación a considerar:} Si el sistema crece mucho (+ 500 usuarios simultáneos diarios) y se agota el plan gratuito, se deberá migrar a Opción 2.

\textbf{OPCIÓN 2: VPS Nacional (HostDime/Axarnet Perú)}

\begin{longtable}{|>{\raggedright\arraybackslash}p{5cm}|>{\centering\arraybackslash}p{3cm}|>{\raggedright\arraybackslash}p{6cm}|}
\hline
\rowcolor{headercolor}
\textcolor{white}{\textbf{Concepto}} & \textcolor{white}{\textbf{Costo Mensual}} & \textcolor{white}{\textbf{Observaciones}} \\
\hline
\endfirsthead
\hline
\rowcolor{headercolor}
\textcolor{white}{\textbf{Concepto}} & \textcolor{white}{\textbf{Costo Mensual}} & \textcolor{white}{\textbf{Observaciones}} \\
\hline
\endhead
\hline
\endlastfoot

Servidor VPS & S/. 150 - 200 & 4 vCPU, 16GB RAM, 500GB SSD \\
\hline
\rowcolor{rowgray}
Backup automático (proveedor) & Incluido & Backups semanales incluidos \\
\hline
Certificado SSL & S/. 0 & Let's Encrypt gratuito \\
\hline
\rowcolor{rowgray}
Soporte técnico del proveedor & Incluido & Chat/ticket incluido en plan \\
\hline
\textbf{TOTAL MENSUAL} & \textbf{S/. 100-150} & \\
\hline
\textbf{TOTAL ANUAL} & \textbf{S/. 1,200-1,800} & Pago anual con 10-15\% descuento \\
\hline
\caption{Costos de Operación - VPS Nacional}
\end{longtable}

\textbf{OPCIÓN 3: Servidor Institucional (OTI)}

\begin{longtable}{|>{\raggedright\arraybackslash}p{5cm}|>{\centering\arraybackslash}p{3cm}|>{\raggedright\arraybackslash}p{6cm}|}
\hline
\rowcolor{headercolor}
\textcolor{white}{\textbf{Concepto}} & \textcolor{white}{\textbf{Costo Mensual}} & \textcolor{white}{\textbf{Observaciones}} \\
\hline
\endfirsthead
\hline
\rowcolor{headercolor}
\textcolor{white}{\textbf{Concepto}} & \textcolor{white}{\textbf{Costo Mensual}} & \textcolor{white}{\textbf{Observaciones}} \\
\hline
\endhead
\hline
\endlastfoot

Energía eléctrica servidor & S/. 50 - 80 & 24/7, aprox. 150-200W consumo \\
\hline
\rowcolor{rowgray}
Disco USB para backups & S/. 10 mensual & Amortización disco 1TB (S/. 200) \\
\hline
Internet institucional & S/. 0 & Ya incluido en presupuesto UNAS \\
\hline
\rowcolor{rowgray}
Personal OTI (mantenimiento) & S/. 0* & *4-6 hrs/sem, personal existente \\
\hline
\textbf{TOTAL MENSUAL} & \textbf{S/. 60-90} & \\
\hline
\textbf{TOTAL ANUAL} & \textbf{S/. 720-1,080} & Más económico que VPS \\
\hline
\caption{Costos de Operación - Servidor Institucional}
\end{longtable}

\textit{*El costo de personal de OTI no se incluye porque ya están en planilla institucional.}

\subsubsection{RP-7.3 Comparativa de Opciones y Recomendación}

\begin{longtable}{|>{\raggedright\arraybackslash}p{3.5cm}|>{\centering\arraybackslash}p{2cm}|>{\centering\arraybackslash}p{2.5cm}|>{\raggedright\arraybackslash}p{6cm}|}
\hline
\rowcolor{headerblue}
\textcolor{white}{\textbf{Opción}} & \textcolor{white}{\textbf{Costo Anual}} & \textcolor{white}{\textbf{Complejidad}} & \textcolor{white}{\textbf{Recomendación}} \\
\hline
\endfirsthead
\hline
\rowcolor{headerblue}
\textcolor{white}{\textbf{Opción}} & \textcolor{white}{\textbf{Costo Anual}} & \textcolor{white}{\textbf{Complejidad}} & \textcolor{white}{\textbf{Recomendación}} \\
\hline
\endhead
\hline
\endlastfoot

Hosting Gratuito & S/. 0 & Muy Baja & \textbf{IDEAL para inicio}. Migrar si crece mucho. \\
\hline
\rowcolor{rowgray}
VPS Nacional & S/. 1,200-1,800 & Media & Si hosting gratuito queda pequeño, buena opción. \\
\hline
Servidor OTI & S/. 720-1,080 & Alta & Solo si OTI tiene servidor disponible y personal capacitado. \\
\hline
\caption{Comparativa de Costos de Operación}
\end{longtable}

\textbf{Recomendación oficial del equipo de desarrollo:}
\begin{enumerate}[noitemsep]
    \item \textbf{Año 1-2:} Opción 1 (Hosting Gratuito) - Sin costo, configuración simple
    \item \textbf{Año 3+:} Si el sistema tiene éxito y crece, evaluar migración a Opción 2 (VPS) con presupuesto institucional
    \item Opción 3 solo si la OTI ya tiene infraestructura y experiencia en administración de servidores
\end{enumerate}

\textit{\textbf{Importante:} Para una universidad pública regional con presupuesto ajustado, la Opción 1 es perfectamente viable y profesional. Muchas startups y empresas pequeñas operan exitosamente con servicios gratuitos durante años.}

\subsection{RP-8. Criterios de Aceptación del Proyecto}

El proyecto se considerará exitoso y aceptado cuando:

\begin{enumerate}[noitemsep]
    \item Todos los requisitos funcionales de prioridad Alta y Media estén implementados y validados.
    \item Todos los requisitos no funcionales críticos cumplan con umbrales definidos (rendimiento, seguridad, disponibilidad).
    \item El sistema supere satisfactoriamente todas las pruebas de aceptación de usuario (UAT).
    \item La documentación técnica y manuales de usuario estén completos y aprobados.
    \item El sistema esté desplegado en ambiente de producción y operativo.
    \item Los usuarios clave hayan sido capacitados y certifiquen competencia en uso del sistema.
    \item El sistema cumpla con al menos 95\% de cobertura de código en pruebas automatizadas.
    \item No existan defectos críticos o bloqueantes sin resolver.
\end{enumerate}

\newpage

\section{Estructura de Desglose del Trabajo (EDT / WBS)}

\subsection{Introducción a la EDT}

La Estructura de Desglose del Trabajo (EDT), conocida internacionalmente como Work Breakdown Structure (WBS), es la descomposición jerárquica orientada a entregables del trabajo que debe ejecutar el equipo del proyecto para cumplir con los objetivos y crear los entregables requeridos. Cada nivel descendente de la EDT representa una definición cada vez más detallada del trabajo del proyecto.

La EDT del proyecto SIGDE se ha estructurado siguiendo las mejores prácticas del PMI (Project Management Institute) y se organiza en 8 paquetes de trabajo principales, cada uno subdividido en entregables específicos. Esta estructura facilita la asignación de responsabilidades, estimación de recursos, seguimiento del progreso y control del alcance del proyecto.

\subsection{Principios de la EDT}

\begin{itemize}[noitemsep]
    \item \textbf{Regla del 100\%:} La EDT incluye el 100\% del trabajo definido por el alcance del proyecto.
    \item \textbf{Orientación a entregables:} Cada elemento representa un entregable verificable, no actividades o funciones.
    \item \textbf{Nivel de detalle apropiado:} Los paquetes de trabajo son lo suficientemente pequeños para ser estimados, asignados y controlados.
    \item \textbf{Mutua exclusividad:} No hay duplicación de trabajo entre componentes de la EDT.
\end{itemize}

\subsection{Desglose Jerárquico del Trabajo}

\subsubsection{Nivel 1: Proyecto SIGED}

\textbf{1.0 PROYECTO SIGDE - Sistema de Gestión Deportiva}

\textit{Nota: La EDT está alineada con el cronograma del proyecto establecido en el Acta de Constitución. Ver archivo adjunto \texttt{edt\_diagrama.tex} para visualización gráfica de la estructura jerárquica.}

\subsubsection{Nivel 2: Paquetes de Trabajo Principales}

\begin{longtable}{|>{\centering\arraybackslash}p{1.5cm}|>{\raggedright\arraybackslash}p{5cm}|>{\raggedright\arraybackslash}p{8cm}|}
\hline
\rowcolor{headercolor}
\textcolor{white}{\textbf{ID EDT}} & \textcolor{white}{\textbf{Paquete de Trabajo}} & \textcolor{white}{\textbf{Descripción}} \\
\hline
\endfirsthead
\hline
\rowcolor{headercolor}
\textcolor{white}{\textbf{ID EDT}} & \textcolor{white}{\textbf{Paquete de Trabajo}} & \textcolor{white}{\textbf{Descripción}} \\
\hline
\endhead
\hline
\endlastfoot

1.0 & Gestión e Inicio & Reuniones con stakeholders, elaboración del acta de constitución, configuración de herramientas. \\
\hline
\rowcolor{rowgray}
2.0 & Planificación & Definición de alcance, EDT, matriz de riesgos, cronograma y plan de comunicaciones. \\
\hline
3.0 & Análisis y Diseño (Ingeniería) & Especificación de requisitos, prototipos UI/UX, arquitectura y modelo de base de datos. \\
\hline
\rowcolor{rowgray}
4.0 & Construcción (Desarrollo) & Implementación de backend, frontend, base de datos e integración de componentes. \\
\hline
5.0 & Aseguramiento de Calidad (QA) & Pruebas unitarias, de integración, de sistema, de rendimiento y aceptación de usuario. \\
\hline
\rowcolor{rowgray}
6.0 & Despliegue en Producción (OTI) & Preparación de infraestructura, instalación de plataforma de contenedores, configuración de base de datos, despliegue de aplicaciones, seguridad y certificados SSL, monitoreo, validación y documentación completa de operaciones. \\
\hline
7.0 & Marcha Blanca y Piloto & Capacitación a usuarios, ejecución de torneo piloto, monitoreo y ajustes finales. \\
\hline
\rowcolor{rowgray}
8.0 & Cierre del Proyecto & Entrega de manuales, acta de aceptación final y transferencia formal a OTI. \\
\hline
\caption{EDT Nivel 2 - Paquetes de Trabajo Principales}
\end{longtable}

\noindent\textit{\textbf{Nota:} El diagrama EDT está disponible en múltiples formatos (PlantUML, Mermaid y LaTeX/TikZ). Ver documentación en README\_EDT.md}

\subsection{Desglose Detallado por Paquete de Trabajo}

\subsubsection{1.0 Gestión e Inicio}

\begin{longtable}{|>{\centering\arraybackslash}p{2.5cm}|>{\raggedright\arraybackslash}p{5cm}|>{\raggedright\arraybackslash}p{7cm}|}
\hline
\rowcolor{headercolor}
\textcolor{white}{\textbf{ID}} & \textcolor{white}{\textbf{Entregable}} & \textcolor{white}{\textbf{Descripción}} \\
\hline
\endfirsthead
\hline
\rowcolor{headercolor}
\textcolor{white}{\textbf{ID}} & \textcolor{white}{\textbf{Entregable}} & \textcolor{white}{\textbf{Descripción}} \\
\hline
\endhead
\hline
\endlastfoot

1.1 & Reuniones con Interesados & Documentación de entrevistas con Deportes/OTI, identificación de necesidades. \\
\hline
\rowcolor{rowgray}
1.2 & Acta de Constitución & Autorización formal del proyecto, designación del PM, objetivos estratégicos. \\
\hline
1.3 & Configuración de Herramientas & Setup de repositorios Git, herramientas CI/CD, ambiente de desarrollo. \\
\hline
\caption{EDT 1.0 - Gestión e Inicio}
\end{longtable}

\subsubsection{2.0 Planificación}

\begin{longtable}{|>{\centering\arraybackslash}p{2.5cm}|>{\raggedright\arraybackslash}p{5cm}|>{\raggedright\arraybackslash}p{7cm}|}
\hline
\rowcolor{headercolor}
\textcolor{white}{\textbf{ID}} & \textcolor{white}{\textbf{Entregable}} & \textcolor{white}{\textbf{Descripción}} \\
\hline
\endfirsthead
\hline
\rowcolor{headercolor}
\textcolor{white}{\textbf{ID}} & \textcolor{white}{\textbf{Entregable}} & \textcolor{white}{\textbf{Descripción}} \\
\hline
\endhead
\hline
\endlastfoot

2.1 & Definición del Alcance & Enunciado detallado del alcance, EDT y diccionario de EDT. \\
\hline
\rowcolor{rowgray}
2.1 & Matriz de Riesgos & Identificación, análisis y estrategias de respuesta a riesgos. \\
\hline
2.2 & Cronograma del Proyecto & Diagrama de Gantt con actividades, duraciones, dependencias, ruta crítica. \\
\hline
\rowcolor{rowgray}
2.2 & Plan de Comunicaciones & Estrategia de comunicación con stakeholders, frecuencia de reportes. \\
\hline
\caption{EDT 2.0 - Planificación}
\end{longtable}

\subsubsection{3.0 Análisis y Diseño (Ingeniería)}

\begin{longtable}{|>{\centering\arraybackslash}p{2.5cm}|>{\raggedright\arraybackslash}p{5cm}|>{\raggedright\arraybackslash}p{7cm}|}
\hline
\rowcolor{headercolor}
\textcolor{white}{\textbf{ID}} & \textcolor{white}{\textbf{Entregable}} & \textcolor{white}{\textbf{Descripción}} \\
\hline
\endfirsthead
\hline
\rowcolor{headercolor}
\textcolor{white}{\textbf{ID}} & \textcolor{white}{\textbf{Entregable}} & \textcolor{white}{\textbf{Descripción}} \\
\hline
\endhead
\hline
\endlastfoot

3.1 & Especificación de Requisitos (SRS) & Documento completo de requisitos funcionales, no funcionales y reglas de negocio. \\
\hline
\rowcolor{rowgray}
3.2 & Prototipos UX/UI (Figma) & Wireframes y mockups de alta fidelidad para Admin, Delegado y Portal Público. \\
\hline
3.3 & Arquitectura C4 & Diagramas de contexto, contenedores, componentes y código (Clean Architecture). \\
\hline
\rowcolor{rowgray}
3.3 & Modelo de Base de Datos & Diagrama entidad-relación, diccionario de datos, esquema físico. \\
\hline
\caption{EDT 3.0 - Análisis y Diseño}
\end{longtable}

\subsubsection{4.0 Construcción (Desarrollo)}

\begin{longtable}{|>{\centering\arraybackslash}p{2.5cm}|>{\raggedright\arraybackslash}p{5cm}|>{\raggedright\arraybackslash}p{7cm}|}
\hline
\rowcolor{headercolor}
\textcolor{white}{\textbf{ID}} & \textcolor{white}{\textbf{Entregable}} & \textcolor{white}{\textbf{Descripción}} \\
\hline
\endfirsthead
\hline
\rowcolor{headercolor}
\textcolor{white}{\textbf{ID}} & \textcolor{white}{\textbf{Entregable}} & \textcolor{white}{\textbf{Descripción}} \\
\hline
\endhead
\hline
\endlastfoot

\multicolumn{3}{|l|}{\cellcolor{headercolor}\textcolor{white}{\textbf{4.1 Backend (.NET 8)}}} \\
\hline
4.1.1 & Configuración Clean Arch + Identity & Setup de arquitectura limpia, autenticación JWT, autorización RBAC. \\
\hline
\rowcolor{rowgray}
4.1.2 & API Configuración y Torneos & Core multi-evento: creación de torneos, inscripciones, configuraciones. \\
\hline
4.1.3 & API Gestión de Reservas & Algoritmo de programación de horarios, validación de disponibilidad de escenarios. \\
\hline
\rowcolor{rowgray}
4.1.4 & Motor de Tiempo Real (SignalR) & Implementación de WebSockets para actualización en vivo de resultados. \\
\hline

\multicolumn{3}{|l|}{\cellcolor{headercolor}\textcolor{white}{\textbf{4.2 Frontend (React / Next.js)}}} \\
\hline
4.2.1 & Panel Admin y Gestión de Usuarios & Maquetación de interfaces administrativas, CRUD de usuarios y roles. \\
\hline
\rowcolor{rowgray}
4.2.2 & Módulo de Reservas y Calendario & Interfaz interactiva de programación de partidos y gestión de escenarios. \\
\hline
4.2.3 & Portal Público (SignalR) & Interfaces públicas con actualización en tiempo real de fixtures y resultados. \\
\hline
\rowcolor{rowgray}
4.2.4 & Acta Digital PWA & Aplicación web progresiva móvil para mesa de control en campo. \\
\hline

\multicolumn{3}{|l|}{\cellcolor{headercolor}\textcolor{white}{\textbf{4.3 Integración}}} \\
\hline
4.3.1 & Integración Final & Conexión completa de API, clientes frontend y base de datos PostgreSQL. \\
\hline
\caption{EDT 4.0 - Construcción (Desarrollo)}
\end{longtable}

\subsubsection{5.0 Aseguramiento de Calidad (QA)}

\begin{longtable}{|>{\centering\arraybackslash}p{2.5cm}|>{\raggedright\arraybackslash}p{5cm}|>{\raggedright\arraybackslash}p{7cm}|}
\hline
\rowcolor{headercolor}
\textcolor{white}{\textbf{ID}} & \textcolor{white}{\textbf{Entregable}} & \textcolor{white}{\textbf{Descripción}} \\
\hline
\endfirsthead
\hline
\rowcolor{headercolor}
\textcolor{white}{\textbf{ID}} & \textcolor{white}{\textbf{Entregable}} & \textcolor{white}{\textbf{Descripción}} \\
\hline
\endhead
\hline
\endlastfoot

5.1 & Pruebas Unitarias (Backend) & Tests automatizados con xUnit/NUnit, cobertura mínima 80\%. \\
\hline
\rowcolor{rowgray}
5.2 & Pruebas de Integración y Flujos & Validación de interacción entre módulos, flujos end-to-end de usuario. \\
\hline
5.3 & Pruebas de Carga (Stress Testing) & Simulación de 500 usuarios concurrentes, identificación de cuellos de botella. \\
\hline
\rowcolor{rowgray}
5.4 & Corrección de Bugs & Sprint de estabilización, resolución de defectos críticos y mayores. \\
\hline
\caption{EDT 5.0 - Aseguramiento de Calidad}
\end{longtable}

\subsubsection{6.0 Despliegue en Producción (OTI)}

\begin{longtable}{|>{\centering\arraybackslash}p{2.5cm}|>{\raggedright\arraybackslash}p{5cm}|>{\raggedright\arraybackslash}p{7cm}|}
\hline
\rowcolor{headercolor}
\textcolor{white}{\textbf{ID}} & \textcolor{white}{\textbf{Entregable}} & \textcolor{white}{\textbf{Descripción}} \\
\hline
\endfirsthead
\hline
\rowcolor{headercolor}
\textcolor{white}{\textbf{ID}} & \textcolor{white}{\textbf{Entregable}} & \textcolor{white}{\textbf{Descripción}} \\
\hline
\endhead
\hline
\endlastfoot

\multicolumn{3}{|l|}{\cellcolor{headerblue}\textcolor{white}{\textbf{6.1 Preparación de Infraestructura}}} \\
\hline
6.1.1 & Aprovisionamiento de Servidor & Servidor físico o virtual con especificaciones mínimas: 8GB RAM, 4 vCPU, 100GB SSD. \\
\hline
\rowcolor{rowgray}
6.1.2 & Instalación de Sistema Operativo & Ubuntu Server 22.04 LTS con actualizaciones de seguridad aplicadas. \\
\hline
6.1.3 & Configuración de Red y Firewall & Asignación de IP estática, configuración UFW con reglas de entrada/salida restrictivas. \\
\hline
\rowcolor{rowgray}
6.1.4 & Hardening de Seguridad & Configuración SSH sin root, fail2ban, actualizaciones automáticas, AppArmor. \\
\hline

\multicolumn{3}{|l|}{\cellcolor{headerblue}\textcolor{white}{\textbf{6.2 Instalación de Plataforma de Contenedores}}} \\
\hline
6.2.1 & Instalación de Docker Engine & Docker Engine 24.0+ Community Edition con systemd habilitado. \\
\hline
\rowcolor{rowgray}
6.2.2 & Instalación de Docker Compose & Docker Compose Plugin v2.20+ para orquestación multi-contenedor. \\
\hline
6.2.3 & Configuración de Registry Privado & Docker Hub privado o registry local para almacenamiento de imágenes. \\
\hline
\rowcolor{rowgray}
6.2.4 & Verificación de Contenedores & Prueba de despliegue de contenedores de prueba para validar configuración. \\
\hline

\multicolumn{3}{|l|}{\cellcolor{headerblue}\textcolor{white}{\textbf{6.3 Configuración de Base de Datos}}} \\
\hline
6.3.1 & Despliegue de PostgreSQL 15 & Contenedor Docker con PostgreSQL 15, volúmenes persistentes configurados. \\
\hline
\rowcolor{rowgray}
6.3.2 & Configuración de Parámetros & Ajuste de max\_connections, shared\_buffers, effective\_cache\_size. \\
\hline
6.3.3 & Creación de Base de Datos & Base de datos \texttt{siged\_prod} con usuario dedicado y permisos restringidos. \\
\hline
\rowcolor{rowgray}
6.3.4 & Ejecución de Migraciones & Aplicación de todas las migraciones de Entity Framework Core. \\
\hline
6.3.5 & Carga de Datos Semilla & Inserción de datos iniciales: roles, permisos, configuraciones generales. \\
\hline
\rowcolor{rowgray}
6.3.6 & Configuración de Backups & Script de respaldo automático diario con pg\_dump, retención de 30 días. \\
\hline

\multicolumn{3}{|l|}{\cellcolor{headerblue}\textcolor{white}{\textbf{6.4 Despliegue de Aplicaciones}}} \\
\hline
6.4.1 & Build de Imágenes Docker & Construcción de imágenes optimizadas: API Backend (.NET 8) y Frontend (React). \\
\hline
\rowcolor{rowgray}
6.4.2 & Configuración de Variables de Entorno & Archivo .env con secretos: ConnectionString, JWT\_SECRET, SMTP\_CONFIG. \\
\hline
6.4.3 & Despliegue con Docker Compose & Ejecución de \texttt{docker compose up -d} con servicios: API, Frontend, DB, Nginx. \\
\hline
\rowcolor{rowgray}
6.4.4 & Verificación de Health Checks & Comprobación de endpoints /health y /ready respondiendo con status 200 OK. \\
\hline
6.4.5 & Configuración de Reverse Proxy & Nginx como proxy inverso: balance de carga, compresión gzip, caché estático. \\
\hline

\multicolumn{3}{|l|}{\cellcolor{headerblue}\textcolor{white}{\textbf{6.5 Seguridad y Certificados}}} \\
\hline
6.5.1 & Registro de Dominio Institucional & Solicitud y configuración de subdominio siged.unas.edu.pe ante OTI. \\
\hline
\rowcolor{rowgray}
6.5.2 & Configuración de DNS & Registro DNS tipo A apuntando IP del servidor, TTL 3600 segundos. \\
\hline
6.5.3 & Instalación de Certbot & Herramienta para obtención y renovación automática de certificados Let's Encrypt. \\
\hline
\rowcolor{rowgray}
6.5.4 & Obtención de Certificado SSL & Generación de certificado SSL/TLS válido con autorenovación cada 90 días. \\
\hline
6.5.5 & Configuración HTTPS en Nginx & Redirección forzada HTTP → HTTPS, HSTS header, TLS 1.3. \\
\hline
\rowcolor{rowgray}
6.5.6 & Pruebas de Seguridad SSL & Validación con SSL Labs obteniendo calificación A+ en seguridad. \\
\hline

\multicolumn{3}{|l|}{\cellcolor{headerblue}\textcolor{white}{\textbf{6.6 Monitoreo y Logging}}} \\
\hline
6.6.1 & Configuración de Logs Centralizados & Docker logging driver para captura de logs de todos los contenedores. \\
\hline
\rowcolor{rowgray}
6.6.2 & Implementación de Uptime Monitoring & Servicio externo (UptimeRobot) verificando disponibilidad cada 5 minutos. \\
\hline
6.6.3 & Dashboard de Métricas & Instalación opcional de Prometheus + Grafana para métricas en tiempo real. \\
\hline
\rowcolor{rowgray}
6.6.4 & Configuración de Alertas & Notificaciones automáticas a OTI ante caídas o errores críticos. \\
\hline

\multicolumn{3}{|l|}{\cellcolor{headerblue}\textcolor{white}{\textbf{6.7 Validación y Pruebas en Producción}}} \\
\hline
6.7.1 & Smoke Tests & Verificación de funcionalidades críticas: Login, creación de torneo, registro resultados. \\
\hline
\rowcolor{rowgray}
6.7.2 & Pruebas de Carga en Producción & Simulación de 500 usuarios concurrentes, validación de tiempos de respuesta. \\
\hline
6.7.3 & Pruebas de Failover & Validación de recuperación automática ante caída de contenedores. \\
\hline
\rowcolor{rowgray}
6.7.4 & Pruebas de Backup y Restore & Restauración exitosa de backup en entorno aislado. \\
\hline

\multicolumn{3}{|l|}{\cellcolor{headerblue}\textcolor{white}{\textbf{6.8 Documentación de Despliegue}}} \\
\hline
6.8.1 & Manual de Despliegue & Documento paso a paso con comandos exactos para replicar despliegue. \\
\hline
\rowcolor{rowgray}
6.8.2 & Runbook de Operaciones & Procedimientos para tareas comunes: restart, actualizaciones, rollback. \\
\hline
6.8.3 & Plan de Rollback & Procedimiento documentado para revertir a versión anterior en <5 minutos. \\
\hline
\rowcolor{rowgray}
6.8.4 & Diagrama de Arquitectura Desplegada & Diagrama de red con todos los componentes, puertos, flujos de conexión. \\
\hline
6.8.5 & Checklist de Aceptación & Verificación de todos los requisitos de infraestructura cumplidos. \\
\hline

\caption{EDT 6.0 - Despliegue en Producción (Expandido)}
\end{longtable}

\subsubsection{7.0 Marcha Blanca y Piloto}

\begin{longtable}{|>{\centering\arraybackslash}p{2.5cm}|>{\raggedright\arraybackslash}p{5cm}|>{\raggedright\arraybackslash}p{7cm}|}
\hline
\rowcolor{headercolor}
\textcolor{white}{\textbf{ID}} & \textcolor{white}{\textbf{Entregable}} & \textcolor{white}{\textbf{Descripción}} \\
\hline
\endfirsthead
\hline
\rowcolor{headercolor}
\textcolor{white}{\textbf{ID}} & \textcolor{white}{\textbf{Entregable}} & \textcolor{white}{\textbf{Descripción}} \\
\hline
\endhead
\hline
\endlastfoot

7.1 & Capacitación a Usuarios & Formación de jueces, administradores y delegados en uso del sistema. \\
\hline
\rowcolor{rowgray}
7.2 & Ejecución de Torneo Piloto & Simulacro real de competencia para validación del sistema en producción. \\
\hline
7.3 & Monitoreo y Ajustes Finales & Análisis de logs, corrección de incidencias, optimizaciones post-piloto. \\
\hline
\caption{EDT 7.0 - Marcha Blanca y Piloto}
\end{longtable}

\subsubsection{8.0 Cierre del Proyecto}

\begin{longtable}{|>{\centering\arraybackslash}p{2.5cm}|>{\raggedright\arraybackslash}p{5cm}|>{\raggedright\arraybackslash}p{7cm}|}
\hline
\rowcolor{headercolor}
\textcolor{white}{\textbf{ID}} & \textcolor{white}{\textbf{Entregable}} & \textcolor{white}{\textbf{Descripción}} \\
\hline
\endfirsthead
\hline
\rowcolor{headercolor}
\textcolor{white}{\textbf{ID}} & \textcolor{white}{\textbf{Entregable}} & \textcolor{white}{\textbf{Descripción}} \\
\hline
\endhead
\hline
\endlastfoot

8.1 & Manuales de Usuario y Operaciones & Documentación completa para usuarios finales y equipo de soporte OTI. \\
\hline
\rowcolor{rowgray}
8.2 & Acta de Cierre y Transferencia & Documento de aceptación final, transferencia formal del sistema a OTI. \\
\hline
\caption{EDT 8.0 - Cierre del Proyecto}
\end{longtable}

\newpage



\begin{landscape}
\subsection{Diagrama Visual de la EDT}
\begin{figure}[H]
\centering
\includegraphics[width=1\linewidth]{edtSIGED.png}
\caption{Estructura de Desglose del Trabajo (EDT) - Vista Jerárquica Completa}
\label{fig:edt_diagrama}
\end{figure}

\end{landscape}

\newpage

\begin{longtable}{|>{\centering\arraybackslash}p{2.5cm}|>{\raggedright\arraybackslash}p{5cm}|>{\raggedright\arraybackslash}p{7cm}|}
\hline
\rowcolor{headercolor}
\textcolor{white}{\textbf{ID}} & \textcolor{white}{\textbf{Entregable}} & \textcolor{white}{\textbf{Descripción}} \\
\hline
\endfirsthead
\hline
\rowcolor{headercolor}
\textcolor{white}{\textbf{ID}} & \textcolor{white}{\textbf{Entregable}} & \textcolor{white}{\textbf{Descripción}} \\
\hline
\endhead
\hline
\endlastfoot

1.1.1 & Acta de Constitución & Autorización formal del proyecto, designación del PM, objetivos de alto nivel. \\
\hline
\rowcolor{rowgray}
1.1.2 & Plan de Gestión del Proyecto & Plan integrado con gestión de alcance, tiempo, costos, calidad, RRHH, comunicaciones, riesgos. \\
\hline
1.1.3 & Cronograma del Proyecto & Diagrama de Gantt con actividades, duraciones, dependencias y ruta crítica. \\
\hline
\rowcolor{rowgray}
1.1.4 & Plan de Gestión de Riesgos & Identificación, análisis y estrategias de respuesta a riesgos. \\
\hline
1.1.5 & Informes de Seguimiento & Reportes semanales/mensuales de avance, KPIs, problemas identificados. \\
\hline
\rowcolor{rowgray}
1.1.6 & Gestión de Cambios & Control de solicitudes de cambio, análisis de impacto, aprobaciones. \\
\hline
1.1.7 & Actas de Reuniones & Registro de reuniones con stakeholders, decisiones tomadas, acuerdos. \\
\hline
\caption{Entregables de Gestión del Proyecto}
\end{longtable}

\subsection{Diccionario de la EDT}

El Diccionario de la EDT proporciona información detallada adicional sobre cada paquete de trabajo, incluyendo:

\begin{itemize}[noitemsep]
    \item \textbf{Identificador único:} Código numérico jerárquico para cada elemento.
    \item \textbf{Descripción del trabajo:} Definición clara del alcance del paquete.
    \item \textbf{Criterios de aceptación:} Condiciones que deben cumplirse para considerar completo el entregable.
    \item \textbf{Supuestos:} Factores que se consideran verdaderos para la planificación.
    \item \textbf{Restricciones:} Limitaciones que afectan la ejecución del paquete de trabajo.
    \item \textbf{Responsable:} Rol o persona asignada al paquete de trabajo.
    \item \textbf{Recursos necesarios:} Personas, equipos, herramientas requeridas.
    \item \textbf{Estimación de esfuerzo:} Horas-persona necesarias para completar el trabajo.
    \item \textbf{Dependencias:} Paquetes de trabajo predecesores requeridos.
\end{itemize}


\newpage

\section{Descripción General del Producto}

\subsection{Perspectiva del Producto}

El Sistema SIGED es un producto de software completamente nuevo, desarrollado específicamente para cubrir las necesidades de gestión deportiva interuniversitaria. No reemplaza sistemas existentes sino que llena un vacío en la digitalización de procesos deportivos académicos actualmente gestionados de forma manual o mediante herramientas genéricas inadecuadas.

El producto se diseña como una aplicación web moderna de arquitectura cliente-servidor, accesible desde navegadores web estándar, sin requerir instalación de software adicional por parte de los usuarios finales. Se integrará con servicios de correo electrónico institucionales para notificaciones y con sistemas de almacenamiento cloud para respaldos.

\subsection{Funciones del Producto}

Las funciones principales que el producto proporcionará son:

\begin{enumerate}[noitemsep]
    \item \textbf{Gestión de Identidad y Acceso:} Autenticación segura, control basado en roles, auditoría de acciones.
    \item \textbf{Administración de Torneos:} Creación, configuración y gestión del ciclo de vida completo de competencias.
    \item \textbf{Gestión de Participantes:} Registro y administración de instituciones, equipos y jugadores.
    \item \textbf{Automatización de Programación:} Generación algorítmica de fixtures considerando múltiples restricciones.
    \item \textbf{Ejecución Digital de Eventos:} Actas electrónicas de partidos, registro de eventos en tiempo real.
    \item \textbf{Procesamiento de Resultados:} Cálculo automático de tablas de posiciones, estadísticas y clasificaciones.
    \item \textbf{Gestión Disciplinaria:} Seguimiento de sanciones, suspensiones y procesos de justicia deportiva.
    \item \textbf{Comunicación Multi-canal:} Sistema de notificaciones por correo electrónico y en plataforma.
    \item \textbf{Transparencia Pública:} Portal de información abierta para consulta de fixtures y resultados.
    \item \textbf{Business Intelligence:} Dashboard de analítica con KPIs deportivos y generación de reportes ejecutivos.
\end{enumerate}

\subsection{Características de los Usuarios}

El Sistema contempla cinco perfiles de usuario con diferentes niveles de acceso y responsabilidades:

\begin{longtable}{|>{\raggedright\arraybackslash}p{2.8cm}|>{\raggedright\arraybackslash}p{3cm}|>{\raggedright\arraybackslash}p{2.5cm}|>{\raggedright\arraybackslash}p{5cm}|}
\hline
\rowcolor{headercolor}
\textcolor{white}{\textbf{Perfil}} & \textcolor{white}{\textbf{Nivel Técnico}} & \textcolor{white}{\textbf{Frecuencia de Uso}} & \textcolor{white}{\textbf{Responsabilidades Principales}} \\
\hline
\endfirsthead
\hline
\rowcolor{headercolor}
\textcolor{white}{\textbf{Perfil}} & \textcolor{white}{\textbf{Nivel Técnico}} & \textcolor{white}{\textbf{Frecuencia de Uso}} & \textcolor{white}{\textbf{Responsabilidades Principales}} \\
\hline
\endhead
\hline
\endlastfoot

Super Administrador & Avanzado & Ocasional & Gestión global del sistema, creación de administradores, configuración crítica. \\
\hline
\rowcolor{rowgray}
Administrador & Intermedio-Avanzado & Frecuente & Gestión de torneos, creación de comités, supervisión general. \\
\hline
Comité Deportivo & Intermedio & Muy Frecuente & Operación de torneos, gestión de disciplinas, equipos y fixtures. \\
\hline
\rowcolor{rowgray}
Delegado de Facultad & Básico-Intermedio & Frecuente & Inscripción de equipos propios, gestión de jugadores de su facultad. \\
\hline
Usuario Público & Básico & Variable & Consulta pasiva de información pública (fixtures, resultados, estadísticas). \\
\hline
\caption{Características de Usuarios del Sistema}
\end{longtable}

\subsection{Restricciones}

\subsubsection{Restricciones Regulatorias y Legales}

\begin{itemize}[noitemsep]
    \item El sistema debe cumplir con legislación de protección de datos personales vigente.
    \item Debe implementar mecanismos de consentimiento para tratamiento de datos personales de jugadores menores de edad.
    \item Debe mantener registros de auditoría para fines de transparencia y rendición de cuentas.
    \item Las actas digitales deben tener validez legal equivalente a documentos físicos firmados.
\end{itemize}

\subsubsection{Restricciones de Hardware}

\begin{itemize}[noitemsep]
    \item El sistema debe ser accesible desde dispositivos con pantallas de al menos 1024x768 píxeles para funcionalidades administrativas.
    \item Debe ser responsivo y funcional en tablets (768px) y smartphones (320px) para consultas públicas.
    \item No requiere hardware especial en cliente, solo navegador web moderno.
\end{itemize}

\subsubsection{Restricciones de Interfaces con Otras Aplicaciones}

\begin{itemize}[noitemsep]
    \item Debe integrarse con servidores SMTP para envío de correos electrónicos.
    \item Puede integrarse con servicios de almacenamiento cloud (AWS S3, Azure Blob Storage).
    \item Debe proporcionar API REST documentada para futuras integraciones con otros sistemas institucionales.
\end{itemize}

\subsubsection{Restricciones de Operaciones Paralelas}

\begin{itemize}[noitemsep]
    \item El sistema debe soportar al menos 500 usuarios concurrentes sin degradación de rendimiento.
    \item Debe manejar correctamente concurrencia en actualización de resultados de partidos simultáneos.
    \item Debe implementar bloqueos optimistas para prevenir inconsistencias en datos compartidos.
\end{itemize}

\subsubsection{Restricciones de Tecnología}

\begin{itemize}[noitemsep]
    \item Backend debe implementarse con tecnologías modernas (Node.js, Python o Java).
    \item Frontend debe utilizar framework web contemporáneo (React, Angular o Vue.js).
    \item Base de datos relacional (PostgreSQL o MySQL) para persistencia principal.
    \item Compatibilidad con navegadores: Chrome 90+, Firefox 88+, Safari 14+, Edge 90+.
\end{itemize}

\subsection{Suposiciones y Dependencias}

\subsubsection{Suposiciones}

\begin{itemize}[noitemsep]
    \item Los usuarios administrativos dispondrán de computadoras con conexión estable a internet.
    \item La institución proporcionará infraestructura de hosting con características adecuadas.
    \item Los usuarios tienen conocimientos básicos de navegación web y uso de formularios digitales.
    \item La institución cuenta con servicio de correo electrónico corporativo funcionando.
    \item Los datos de torneos históricos pueden no estar disponibles para migración inicial.
\end{itemize}

\subsubsection{Dependencias}

\begin{itemize}[noitemsep]
    \item Disponibilidad de servicios de hosting cloud (AWS, Azure, GCP o similar).
    \item Proveedor de servicios SMTP para envío de correos electrónicos.
    \item Certificados SSL/TLS válidos para comunicación segura HTTPS.
    \item Disponibilidad de stakeholders clave para validación de requisitos y pruebas de aceptación.
    \item Acceso a información institucional necesaria (logos, reglamentos deportivos, estructura facultades).
\end{itemize}

\subsection{Arquitectura del Sistema}

El sistema se diseñó utilizando \textbf{Clean Architecture} (propuesta por Robert C. Martin) organizado en 16 módulos funcionales. Esta arquitectura separa claramente las responsabilidades y facilita el mantenimiento futuro del código.

\subsubsection{Organización por Capas}

El sistema está dividido en 4 capas que se comunican siguiendo la regla de dependencia (las capas externas dependen de las internas, nunca al revés):

\paragraph{1. Capa de Dominio (Núcleo)}

Contiene la lógica de negocio pura sin dependencias externas:

\begin{itemize}[noitemsep]
    \item \textbf{Entidades:} Torneo, Equipo, Partido, Usuario, Disciplina, Escenario
    \item \textbf{Objetos de Valor:} Marcador, Resultado, Estado del Partido
    \item \textbf{Servicios de Dominio:} Cálculo de tablas, validación de reglas deportivas
    \item \textbf{Eventos:} PartidoIniciado, GolAnotado, TorneoFinalizado
\end{itemize}

\paragraph{2. Capa de Aplicación}

Implementa los casos de uso del sistema:

\begin{itemize}[noitemsep]
    \item \textbf{Casos de Uso:} CrearTorneo, InscribirEquipo, RegistrarResultado, GenerarFixture
    \item \textbf{DTOs:} Modelos de transferencia de datos entre capas
    \item \textbf{Interfaces:} Contratos de repositorios y servicios externos
\end{itemize}

\paragraph{3. Capa de Infraestructura}

Implementa las interfaces definidas en la capa de aplicación:

\begin{itemize}[noitemsep]
    \item \textbf{Repositorios:} Acceso a datos con Entity Framework Core + PostgreSQL
    \item \textbf{Servicios Externos:} Email (SMTP), notificaciones, almacenamiento de archivos
    \item \textbf{Autenticación:} JWT + ASP.NET Identity para login y control de acceso
\end{itemize}

\paragraph{4. Capa de Presentación}

Expone la funcionalidad del sistema:

\begin{itemize}[noitemsep]
    \item \textbf{API REST:} Endpoints con ASP.NET Core Web API
    \item \textbf{Frontend:} React 18 + TypeScript + TailwindCSS
    \item \textbf{Tiempo Real:} SignalR para actualizaciones en vivo de marcadores
    \item \textbf{Validaciones:} FluentValidation en entrada de datos
\end{itemize}

\subsubsection{Módulos Funcionales}

El sistema se divide en 16 módulos independientes:

\begin{enumerate}[noitemsep]
    \item \textbf{Autenticación y Seguridad:} Login, permisos, auditoría
    \item \textbf{Usuarios:} Administración de cuentas con roles (RBAC)
    \item \textbf{Gestión Institucional:} Facultades y escuelas
    \item \textbf{Disciplinas Deportivas:} Catálogo de deportes
    \item \textbf{Torneos:} Creación y configuración de competencias
    \item \textbf{Equipos y Participantes:} Registro de equipos y jugadores
    \item \textbf{Escenarios y Horarios:} Infraestructura deportiva
    \item \textbf{Generación de Fixture:} Calendario automático de partidos
    \item \textbf{Ejecución de Partidos:} Actas digitales en tiempo real
    \item \textbf{Resultados y Tablas:} Clasificaciones y estadísticas
    \item \textbf{Justicia Deportiva:} Sanciones y multas
    \item \textbf{Notificaciones:} Comunicación con usuarios
    \item \textbf{Portal Público:} Consulta de información abierta
    \item \textbf{Reportes:} Exportación de datos y documentos
    \item \textbf{Analítica:} Dashboard con indicadores (KPIs)
    \item \textbf{Configuración:} Parámetros del sistema
\end{enumerate}

\subsubsection{Tecnologías Utilizadas}

\begin{itemize}[noitemsep]
    \item \textbf{Frontend:} React 18 + TypeScript + Vite + TailwindCSS 3 + Radix UI
    \item \textbf{Backend:} .NET 8 + ASP.NET Core Web API + MediatR (CQRS)
    \item \textbf{Base de Datos:} PostgreSQL 15 + Entity Framework Core 8
    \item \textbf{Tiempo Real:} SignalR para marcadores en vivo
    \item \textbf{Autenticación:} JWT Bearer Tokens + ASP.NET Identity
    \item \textbf{Contenedores:} Docker + Docker Compose
\end{itemize}

\subsection{Requisitos de Infraestructura y Despliegue}

Esta sección especifica los requisitos técnicos, de hardware, software y configuración necesarios para el despliegue exitoso del sistema SIGED en el ambiente de producción de la Oficina de Tecnologías de la Información (OTI) de la UNAS. Los requisitos han sido adaptados a la realidad de una universidad pública regional, priorizando soluciones de bajo costo y alta eficiencia.


\subsubsection{Opciones de Despliegue (por Orden de Preferencia)}

El sistema SIGED puede desplegarse en múltiples escenarios según los recursos disponibles:

\paragraph{Opción 1: Cloud Moderno (Recomendada)}

\textbf{Plataforma:} Netlify/Vercel (frontend) + PostgreSQL cloud (backend)\\
\textbf{Ventajas:} Configuración rápida, escalabilidad automática, SSL incluido, actualizaciones transparentes\\
\textbf{Consideraciones:} Límites de 100GB transferencia/mes en frontend, funciones con timeout de 10s\\
\textbf{Uso apropiado:} Sistema completo en producción, recomendado para universidades\\
\textbf{Costo estimado:} Incluido en planes educativos

\paragraph{Opción 2: VPS Nacional}

\textbf{Proveedor:} HostDime Perú, Axarnet Perú\\
\textbf{Especificaciones:} 4 vCPU, 16GB RAM, 500GB SSD, 3TB transferencia\\
\textbf{Ventajas:} Control total del servidor, latencia baja, soporte local\\
\textbf{Uso apropiado:} Si se requiere integración con infraestructura existente de la universidad\\
\textbf{Costo estimado:} S/. 150 - S/. 200 mensuales

\paragraph{Opción 3: Servidor Institucional}

\textbf{Requisitos:} Ver especificaciones en tabla siguiente\\
\textbf{Ventajas:} Control total, datos on-premise\\
\textbf{Requisitos:} Personal TI capacitado, UPS, respaldos\\
\textbf{Uso apropiado:} Si OTI cuenta con servidor y equipo de administración\\
\textbf{Costo estimado:} Electricidad + mantenimiento (S/. 80-120 mensuales)

\subsubsection{Requisitos de Hardware (Servidor Físico - Opción 3)}

Solo aplicable si se decide hospedar en servidores propios de la OTI:

\begin{longtable}{|>{\raggedright\arraybackslash}p{4cm}|>{\raggedright\arraybackslash}p{4.5cm}|>{\raggedright\arraybackslash}p{5cm}|}
\hline
\rowcolor{headercolor}
\textcolor{white}{\textbf{Componente}} & \textcolor{white}{\textbf{Mínimo Aceptable}} & \textcolor{white}{\textbf{Recomendado (si aplica)}} \\
\hline
\endfirsthead
\hline
\rowcolor{headercolor}
\textcolor{white}{\textbf{Componente}} & \textcolor{white}{\textbf{Mínimo Aceptable}} & \textcolor{white}{\textbf{Recomendado (si aplica)}} \\
\hline
\endhead
\hline
\endlastfoot

Procesador & 2 núcleos (2 vCPU) @ 2.0 GHz & 4 núcleos @ 2.4 GHz \\
\hline
\rowcolor{rowgray}
Memoria RAM & 4 GB (suficiente para 200-300 usuarios) & 8 GB (para 500+ usuarios) \\
\hline
Almacenamiento & 40 GB disco duro (HDD aceptable) & 80 GB SSD (mejor rendimiento) \\
\hline
\rowcolor{rowgray}
Red & 10 Mbps dedicado (típico UNAS) & 20 Mbps o más \\
\hline
Disponibilidad & 95\% uptime (aceptable para uso interno) & 98\% con UPS institucional \\
\hline
\caption{Requisitos de Hardware }
\end{longtable}

\textit{\textbf{Nota importante:} Muchas universidades regionales peruanas operan exitosamente sistemas similares con estas especificaciones modestas. La clave está en optimización de código y uso eficiente de recursos.}

\paragraph{Sistema Operativo y Software Base (Para servidor propio)}

\begin{itemize}[noitemsep]
    \item \textbf{Sistema Operativo:} Ubuntu Server 22.04 LTS (gratuito) o Windows Server (si ya tienen licencia)
    \item \textbf{Justificación:} Ubuntu es gratuito, estable y ampliamente documentado. Sin costos de licenciamiento
    \item \textbf{Alternativa económica:} Linux Mint Server o Debian 11+ (también gratuitos)
    \item \textbf{Arquitectura:} x86\_64 (64-bit) - Compatible con PCs de escritorio reutilizadas
\end{itemize}

\paragraph{Estrategia de Despliegue Simplificada}

\textbf{Para Hosting Gratuito (Opción 1):}
\begin{itemize}[noitemsep]
    \item Frontend: Archivos estáticos en Netlify/Vercel
    \item Backend: PostgreSQL 15 en cloud (10GB incluidos en planes educativos)
    \item Archivos: Cloud storage (100GB para archivos e imágenes)
    \item Sin necesidad de Docker ni configuración de servidor
\end{itemize}

\textbf{Para VPS Nacional (Opción 2):}
\begin{itemize}[noitemsep]
    \item Panel de control cPanel o Plesk incluido (facilita administración)
    \item Node.js + PostgreSQL instalados mediante panel
    \item Sin contenedores Docker (simplifica operación para personal con menos experiencia)
    \item Backups automatizados por el proveedor
\end{itemize}

\textbf{Para Servidor Institucional (Opción 3):}
\begin{itemize}[noitemsep]
    \item \textbf{Opción Simple:} Instalación directa (Node.js + PostgreSQL + Nginx)
    \item \textbf{Opción Moderna:} Docker Compose (si el personal de OTI está capacitado)
    \item \textbf{Justificación opción simple:} Menor complejidad, troubleshooting más sencillo
\end{itemize}

\paragraph{Base de Datos}

\begin{itemize}[noitemsep]
    \item \textbf{Motor:} PostgreSQL 15+ (puede ser autohospedado o en servicios cloud)
    \item \textbf{Almacenamiento para BD:} 10-20 GB suficiente para 3-5 años de datos 
    \item \textbf{Conexiones concurrentes:} 50-100 conexiones (más que suficiente para carga esperada)
    \item \textbf{Backups:} 
    \begin{itemize}
        \item Cloud: Backups automáticos incluidos en mayoría de proveedores
        \item VPS: Backups incluidos por proveedor 
        \item Servidor propio: Script cron diario con pg\_dump guardado en disco externo USB
    \end{itemize}
    \item \textbf{Replicación:} No necesaria para esta escala (complejidad innecesaria)
\end{itemize}

\subsubsection{Requisitos de Red y Conectividad}

\paragraph{Dominio y Acceso Web}

\begin{itemize}[noitemsep]
    \item \textbf{Opción 1 (Hosting gratuito):} Dominio generado automáticamente (ej: siged-unas.netlify.app) - Gratuito
    \item \textbf{Opción 2 (Subdominio institucional):} Coordinación con OTI para crear siged.unas.edu.pe
    \item \textbf{Requisito mínimo de red:} Conexión a internet estable de al menos 5 Mbps (disponible en UNAS)
    \item \textbf{IP estática:} Solo necesaria si se usa servidor físico institucional (coordinar con OTI)
\end{itemize}

\paragraph{Certificados SSL/TLS (HTTPS)}

\begin{itemize}[noitemsep]
    \item \textbf{Netlify/Vercel:} Certificado SSL automático incluido (Let's Encrypt) - Gratuito
    \item \textbf{VPS Nacional:} Certificado SSL incluido o Let's Encrypt vía Certbot - Gratuito
    \item \textbf{Servidor propio:} Let's Encrypt vía Certbot (gratuito, renovación automática)
    \item \textbf{Requerimiento:} Conexión HTTPS obligatoria para proteger credenciales de usuarios
\end{itemize}

\paragraph{Seguridad de Red Básica}

\begin{itemize}[noitemsep]
    \item \textbf{Firewall:} UFW (Ubuntu Firewall) con configuración básica - Ya incluido en Ubuntu
    \item \textbf{Protección DDoS básica:} Cloudflare Free (plan gratuito) - Opcional pero recomendado
    \item \textbf{Puertos abiertos:} Solo 80 (HTTP) y 443 (HTTPS) - Puerto 22 (SSH) solo desde red UNAS
    \item \textbf{Fail2ban:} Recomendado para servidor propio (protección anti fuerza bruta) - Gratuito
    \item \textbf{Rate Limiting:} 30-50 peticiones/minuto por IP (suficiente para uso normal)
\end{itemize}

\textit{\textbf{Realidad UNAS:} La infraestructura de red actual de la universidad es suficiente para soportar el sistema. No se requieren upgrades de conectividad.}

\subsubsection{Requisitos de Seguridad Básica}

\paragraph{Seguridad del Sistema (Solo para servidor propio)}

\begin{itemize}[noitemsep]
    \item \textbf{Actualizaciones:} Script semanal de actualizaciones de seguridad (cron job simple)
    \item \textbf{Acceso SSH:} Cambiar puerto SSH por defecto (22 → otro), usar contraseñas fuertes
    \item \textbf{Usuario root:} Crear usuario administrativo, evitar usar root directamente
    \item \textbf{Logs del sistema:} Logs estándar de Linux, revisión mensual manual
\end{itemize}

\textit{\textbf{Nota:} Para hosting en Netlify/Vercel o VPS administrado, la seguridad del servidor es responsabilidad del proveedor.}

\paragraph{Seguridad de la Aplicación}

\begin{itemize}[noitemsep]
    \item \textbf{Credenciales:} Variables de entorno nunca subidas a GitHub (archivo .env en .gitignore)
    \item \textbf{Contraseñas de BD:} Generadas aleatoriamente, almacenadas en archivo .env del servidor
    \item \textbf{JWT Secret:} Clave aleatoria de 32+ caracteres, regenerada por entorno
    \item \textbf{CORS:} Permitir solo el dominio donde está el frontend
    \item \textbf{Validación de entrada:} Validar todos los datos de formularios antes de guardar en BD
\end{itemize}

\subsubsection{Monitoreo Simplificado}

\paragraph{Monitoreo Básico Gratuito}

\begin{itemize}[noitemsep]
    \item \textbf{Uptime monitoring:} UptimeRobot plan gratuito (chequeo cada 5 minutos, 50 monitores gratis)
    \item \textbf{Notificaciones:} Email o WhatsApp cuando el sitio cae (configurar en UptimeRobot)
    \item \textbf{Logs de aplicación:} 
    \begin{itemize}
        \item Netlify/Vercel: Panel integrado con logs en tiempo real
        \item Servidor propio: Archivos de log en /var/log/, revisión cuando hay problemas
    \end{itemize}
    \item \textbf{Métricas:} Revisar uso de recursos (CPU, RAM, disco) una vez por semana
\end{itemize}

\textit{\textbf{Pragmatismo:} Para una universidad pequeña, monitoreo sofisticado con Prometheus/Grafana es sobredimensionado. El uptime monitoring gratuito es suficiente.}

\subsubsection{Backups y Recuperación de Datos}

\paragraph{Estrategia de Respaldos por Escenario}

\textbf{Opción 1 - PostgreSQL en Cloud:}
\begin{itemize}[noitemsep]
    \item Backups automáticos diarios incluidos en plan gratuito
    \item Retención: 7 días en plan Free, 30 días si se actualiza a Pro (US\$ 25/mes)
    \item Recuperación: Un clic desde el panel de administración del proveedor
    \item \textbf{Backup manual adicional:} Exportar SQL una vez por mes y guardar en Google Drive UNAS
\end{itemize}

\textbf{Opción 2 - VPS Nacional:}
\begin{itemize}[noitemsep]
    \item Backups semanales automáticos por el proveedor (verificar que esté incluido)
    \item Backup manual mensual adicional: pg\_dump exportado y descargado a PC de OTI
    \item Almacenamiento: Disco externo USB en oficina de OTI
\end{itemize}

\textbf{Opción 3 - Servidor Institucional:}
\begin{itemize}[noitemsep]
    \item \textbf{Script de backup diario:} Cron job que ejecuta pg\_dump a las 3:00 AM
    \item \textbf{Almacenamiento local:} Carpeta /backups/ en el servidor (retención 15 días)
    \item \textbf{Backup externo:} Copiar manualmente a disco USB cada semana
    \item \textbf{Alternativa cloud:} Subir a Google Drive institucional (15 GB gratis por cuenta)
\end{itemize}

\paragraph{Plan de Recuperación Realista}

\begin{itemize}[noitemsep]
    \item \textbf{RTO (Recovery Time Objective):} 24 horas máximo para restauración completa (realista para OTI)
    \item \textbf{RPO (Recovery Point Objective):} Pérdida máxima de 1-7 días según estrategia de backup elegida
    \item \textbf{Responsable:} Jefe de OTI o personal técnico designado
    \item \textbf{Procedimiento:} Documento simple de 1-2 páginas con pasos de restauración
    \item \textbf{Prueba anual:} Realizar prueba de restauración una vez al año (no trimestral)
\end{itemize}

\textit{\textbf{Importante:} Para una universidad sin presupuesto para soluciones enterprise, combinar backups automáticos del proveedor + exportaciones manuales mensuales es una estrategia sólida y de costo cero.}

\subsubsection{Estrategia de Despliegue y Actualizaciones}

\paragraph{Despliegue Inicial}

\textbf{Opción Recomendada para UNAS - Hosting Gratuito:}
\begin{enumerate}[noitemsep]
    \item Configurar PostgreSQL en servicio cloud (Railway, Render, Supabase, etc.)
    \item Crear cuenta en Netlify o Vercel (frontend + backend)
    \item Conectar repositorio GitHub al servicio de hosting
    \item Configurar variables de entorno en el dashboard
    \item Deploy automático al hacer push a rama main
    \item \textbf{Tiempo total de setup:} 2-3 horas para alguien con experiencia básica
\end{enumerate}

\paragraph{Actualizaciones del Sistema}

\begin{itemize}[noitemsep]
    \item \textbf{Hosting gratuito:} Actualización automática al subir código a GitHub (CI/CD incluido)
    \item \textbf{VPS:} FTP/SFTP para subir archivos actualizados, reiniciar servicio manualmente
    \item \textbf{Servidor propio:} Git pull + pm2 restart (si se usa Node.js) o systemctl restart
    \item \textbf{Frecuencia:} Actualizaciones cada 2-3 meses (no cada semana como en empresas grandes)
    \item \textbf{Ventana:} Domingos por la noche (menos uso del sistema)
    \item \textbf{Rollback:} Mantener versión anterior por 1 semana por si hay problemas
\end{itemize}

\subsubsection{Capacitación y Soporte Técnico}

\paragraph{Requisitos de Personal}

\begin{itemize}[noitemsep]
    \item \textbf{OTI - Administrador del sistema:} 
    \begin{itemize}
        \item Conocimientos: Básicos de Linux/Windows Server, MySQL/PostgreSQL
        \item Tiempo dedicado: 4-6 horas semanales para mantenimiento rutinario
        \item Capacitación: 1 día de training sobre el sistema SIGED
    \end{itemize}
    \item \textbf{Departamento de Deportes - Usuarios administradores:}
    \begin{itemize}
        \item Conocimientos: Uso básico de computadora y navegadores web
        \item Capacitación: 2 sesiones de 3 horas (gestión de torneos y reportes)
    \end{itemize}
\end{itemize}

\paragraph{Documentación Entregable}

\begin{itemize}[noitemsep]
    \item Manual de usuario (PDF, 15-20 páginas)
    \item Manual de administrador (PDF, 10 páginas)
    \item Guía rápida de troubleshooting (1 página, imprimir y pegar en pared de OTI)
    \item Video tutoriales cortos (5-10 minutos) subidos a YouTube institucional
\end{itemize}

\subsubsection{Checklist de Aceptación del Entorno de Producción}

Antes de considerar el entorno de producción como apto para el despliegue del sistema SIGED, se debe verificar el cumplimiento de los siguientes criterios adaptados a la realidad de la UNAS:

\begin{longtable}{|>{\centering\arraybackslash}p{0.8cm}|>{\raggedright\arraybackslash}p{9cm}|>{\centering\arraybackslash}p{1.5cm}|>{\raggedright\arraybackslash}p{2.5cm}|}
\hline
\rowcolor{headercolor}
\textcolor{white}{\textbf{\#}} & \textcolor{white}{\textbf{Criterio de Aceptación}} & \textcolor{white}{\textbf{Estado}} & \textcolor{white}{\textbf{Responsable}} \\
\hline
\endfirsthead
\hline
\rowcolor{headercolor}
\textcolor{white}{\textbf{\#}} & \textcolor{white}{\textbf{Criterio de Aceptación}} & \textcolor{white}{\textbf{Estado}} & \textcolor{white}{\textbf{Responsable}} \\
\hline
\endhead
\hline
\endlastfoot

\multicolumn{4}{|l|}{\cellcolor{columnblue}\textbf{OPCIÓN 1: Hosting Gratuito (Recomendado para inicio)}} \\
\hline
1 & Base de datos PostgreSQL configurada y funcional en cloud & ☐ & Desarrollador \\
\hline
\rowcolor{rowgray}
2 & Cuenta creada en Netlify o Vercel con repositorio GitHub conectado & ☐ & Desarrollador \\
\hline
3 & Variables de entorno configuradas correctamente en dashboard & ☐ & Desarrollador \\
\hline
\rowcolor{rowgray}
4 & Deploy exitoso y sitio accesible públicamente & ☐ & Desarrollador \\
\hline
5 & Certificado SSL automático activo (candado verde en navegador) & ☐ & Desarrollador \\
\hline
\rowcolor{rowgray}
6 & Sistema de login funcionando correctamente & ☐ & QA/OTI \\
\hline
7 & Backup manual inicial realizado y descargado & ☐ & OTI \\
\hline
\rowcolor{rowgray}
8 & UptimeRobot configurado con alertas a email de OTI & ☐ & Desarrollador \\
\hline
9 & Manuales de usuario entregados (PDF) & ☐ & Desarrollador \\
\hline
\rowcolor{rowgray}
10 & Capacitación realizada a 2 personas de Dpto. de Deportes & ☐ & Desarrollador \\
\hline

\multicolumn{4}{|l|}{\cellcolor{columnblue}\textbf{OPCIÓN 2: VPS Nacional}} \\
\hline
11 & Contrato firmado con proveedor de VPS peruano & ☐ & OTI \\
\hline
\rowcolor{rowgray}
12 & Accesos de administración (cPanel/SSH) recibidos y probados & ☐ & OTI \\
\hline
13 & PostgreSQL instalado y accesible desde aplicación & ☐ & Desarrollador \\
\hline
\rowcolor{rowgray}
14 & Dominio siged.unas.edu.pe apuntando al VPS (DNS configurado) & ☐ & OTI \\
\hline
15 & Certificado SSL Let's Encrypt instalado y activo & ☐ & OTI/Proveedor \\
\hline
\rowcolor{rowgray}
16 & Aplicación desplegada y accesible vía HTTPS & ☐ & Desarrollador \\
\hline
17 & Backup semanal automático del proveedor confirmado & ☐ & OTI \\
\hline
\rowcolor{rowgray}
18 & Backup manual mensual programado en calendario & ☐ & OTI \\
\hline

\multicolumn{4}{|l|}{\cellcolor{columnblue}\textbf{OPCIÓN 3: Servidor Institucional}} \\
\hline
19 & Servidor físico o virtual asignado con specs mínimas cumplidas & ☐ & OTI \\
\hline
\rowcolor{rowgray}
20 & Sistema operativo Ubuntu Server 22.04 LTS instalado & ☐ & OTI \\
\hline
21 & IP estática asignada al servidor & ☐ & OTI \\
\hline
\rowcolor{rowgray}
22 & Firewall UFW configurado (solo puertos 80, 443, 22 abiertos) & ☐ & OTI \\
\hline
23 & PostgreSQL instalado y configurado & ☐ & Desarrollador \\
\hline
\rowcolor{rowgray}
24 & Aplicación desplegada y funcionando & ☐ & Desarrollador \\
\hline
25 & Subdominio siged.unas.edu.pe creado y apuntando al servidor & ☐ & OTI \\
\hline
\rowcolor{rowgray}
26 & Certificado SSL Let's Encrypt instalado con Certbot & ☐ & Desarrollador \\
\hline
27 & Script de backup automático creado y probado exitosamente & ☐ & OTI \\
\hline
\rowcolor{rowgray}
28 & Disco USB externo adquirido para backups offline & ☐ & OTI \\
\hline
29 & Personal de OTI capacitado en procedimientos básicos & ☐ & Desarrollador \\
\hline

\multicolumn{4}{|l|}{\cellcolor{columnblue}\textbf{VALIDACIÓN GENERAL (Todas las opciones)}} \\
\hline
30 & Pruebas de funcionalidad crítica completadas exitosamente & ☐ & QA \\
\hline
\rowcolor{rowgray}
31 & Al menos 20 usuarios de prueba creados (diferentes roles) & ☐ & Desarrollador \\
\hline
32 & Torneo de prueba completo ejecutado de inicio a fin & ☐ & Dpto. Deportes \\
\hline
\rowcolor{rowgray}
33 & Documentación técnica entregada a OTI & ☐ & Desarrollador \\
\hline
34 & Acta de entrega firmada por Dpto. de Deportes y OTI & ☐ & Todos \\
\hline

\caption{Checklist de Aceptación Adaptado a la Realidad UNAS}
\end{longtable}

\textit{\textbf{Nota Importante:} Este checklist tiene 3 secciones según la opción de despliegue elegida. Solo debe completarse la sección correspondiente a la opción seleccionada, más la validación general que aplica a todas.}

\newpage

Los siguientes diagramas de casos de uso ilustran las principales interacciones entre los actores del sistema y las funcionalidades críticas del SIGED. Estos diagramas representan los flujos de trabajo fundamentales que soportan la gestión deportiva universitaria, desde la inscripción de equipos hasta la oficialización de resultados.

\subsection{Inscripción de Equipo por Promoción}

Este diagrama representa el proceso mediante el cual un delegado registra un equipo conformado por estudiantes de la misma promoción. El sistema valida automáticamente que cada integrante esté activo y pertenezca al año correspondiente antes de permitir su incorporación al equipo. Finalmente, verifica que se cumpla la cantidad mínima de jugadores y genera una constancia oficial de inscripción.

\begin{figure}[H]
    \centering
    \includegraphics[width=0.85\textwidth]{Promoción Inscripción-2026-02-13-231458.png}
    \caption{Diagrama de Caso de Uso: Inscripción de Equipo por Promoción}
    \label{fig:caso_uso_promocion}
\end{figure}
\newpage
\subsection{Inscripción de Equipo por Facultad}

Este diagrama muestra el flujo de inscripción cuando un equipo representa a una facultad específica. El sistema valida que cada estudiante esté matriculado en el semestre actual y que pertenezca a la misma facultad del delegado. Una vez cumplidas las reglas deportivas y académicas, se confirma la inscripción oficial y se genera el documento correspondiente.

\begin{figure}[H]
    \centering
    \includegraphics[width=0.85\textwidth]{Facultad Enrollment-2026-02-13-233801.png}
    \caption{Diagrama de Caso de Uso: Inscripción de Equipo por Facultad}
    \label{fig:caso_uso_facultad}
\end{figure}
\newpage
\subsection{Acta Digital – Desarrollo del Partido (Futsal)}

Este diagrama describe la gestión operativa de un partido de futsal desde la perspectiva del juez de mesa. Incluye el control del cronómetro, registro de goles, validación de faltas acumuladas y actualización automática del marcador en tiempo real. El sistema garantiza que las reglas del juego se apliquen correctamente y que la información se refleje inmediatamente para el público.

\begin{figure}[H]
    \centering
    \includegraphics[width=0.85\textwidth]{Ejecución Inter-Promociones-futsal-2026-02-14-000055.png}
    \caption{Diagrama de Caso de Uso: Acta Digital – Desarrollo del Partido (Futsal)}
    \label{fig:caso_uso_futsal}
\end{figure}
\newpage
\subsection{Acta Digital – Gestión de Partido con Tiempo Extra (Básquet)}

Este diagrama representa la administración de un partido de básquet, incluyendo la validación del marcador final y la activación automática de prórroga en caso de empate. También contempla el cierre administrativo del acta, la asignación de puntos según reglamento y la actualización de la tabla de posiciones.

\begin{figure}[H]
    \centering
    \includegraphics[width=0.85\textwidth]{Básquet Inter-Promociones-cierre-2026-02-14-010728.png}
    \caption{Diagrama de Caso de Uso: Acta Digital – Gestión de Partido con Tiempo Extra (Básquet)}
    \label{fig:caso_uso_basquet}
\end{figure}
\newpage
\subsection{Acta Digital – Proceso de W.O. (Walk Over)}

Este diagrama muestra el procedimiento que se ejecuta cuando un equipo no se presenta dentro del tiempo reglamentario. El sistema detecta la ausencia, aplica automáticamente el resultado por W.O., actualiza las estadísticas y notifica el estado final del partido.

\begin{figure}[H]
    \centering
    \includegraphics[width=0.85\textwidth]{Ejecucion Inter-Facultades-futsal-2026.png}
    \caption{Diagrama de Caso de Uso: Acta Digital – Proceso de W.O. (Walk Over)}
    \label{fig:caso_uso_walkover}
\end{figure}
\newpage
\subsection{Cierre y Oficialización del Partido}

Este diagrama describe el proceso final de validación y cierre oficial de un partido. Incluye la verificación de integridad de datos, captura de firmas digitales, bloqueo del acta para evitar modificaciones posteriores, ejecución del motor de puntuación y generación del acta oficial en formato digital.

\begin{figure}[H]
    \centering
    \includegraphics[width=0.85\textwidth]{Inter-facultades-CIERRE-Futsal.png}
    \caption{Diagrama de Caso de Uso: Cierre y Oficialización del Partido}
    \label{fig:caso_uso_cierre}
\end{figure}
\newpage
\subsection{Acta Digital – Gestión de Partido con Tiempo Extra (Vóley)}

Este diagrama representa la administración de un partido de vóley, incluyendo la validación de sets ganados, el registro de puntos por set y la gestión de tiempos técnicos. El sistema verifica las reglas específicas del vóley como la diferencia de puntos mínima para ganar un set, controla los cambios de campo y maneja la lógica de sets decisivos. También contempla el cierre administrativo del acta y la actualización de estadísticas del torneo.

\begin{figure}[H]
    \centering
    \includegraphics[width=0.85\textwidth]{Vóley Inter-Promoción-cierre-2026-02-14-011921.png}
    \caption{Diagrama de Caso de Uso: Acta Digital – Gestión de Partido con Tiempo Extra (Vóley)}
    \label{fig:caso_uso_voley}
\end{figure}

\newpage

\section{Requisitos Funcionales del Producto}

Esta sección especifica todas las capacidades, comportamientos y funciones que el Sistema debe proporcionar para satisfacer las necesidades de los usuarios y objetivos del proyecto. Los requisitos están organizados por módulos funcionales, priorizados y trazables.

\subsection{RF-1. Módulo de Autenticación y Seguridad}

Este módulo constituye la base de seguridad del sistema, garantizando que solo usuarios autorizados puedan acceder a las funcionalidades según su rol. Implementa mecanismos de autenticación robustos, gestión de sesiones seguras, y auditoría completa de acciones críticas. El sistema de roles jerárquico permite un control granular de permisos, adaptándose a la estructura organizativa de la institución deportiva.

\begin{longtable}{|>{\centering\arraybackslash}p{1.5cm}|>{\raggedright\arraybackslash}p{3.2cm}|>{\raggedright\arraybackslash}p{9cm}|}
\hline
\rowcolor{headercolor}
\textcolor{white}{\textbf{ID}} & \textcolor{white}{\textbf{Requisito}} & \textcolor{white}{\textbf{Descripción}} \\
\hline
\endfirsthead

\multicolumn{3}{c}{\textit{Continuación de la tabla anterior}} \\
\hline
\rowcolor{headercolor}
\textcolor{white}{\textbf{ID}} & \textcolor{white}{\textbf{Requisito}} & \textcolor{white}{\textbf{Descripción}} \\
\hline
\endhead

\hline
\multicolumn{3}{r}{\textit{Continúa en la siguiente página}} \\
\endfoot

\hline
\endlastfoot

RF-001 & Login de Usuario & El sistema debe permitir a los usuarios autenticarse mediante credenciales (correo electrónico y contraseña). \\
\hline
\rowcolor{rowgray}
RF-002 & Logout de Usuario & El sistema debe permitir a los usuarios cerrar sesión de forma segura, invalidando el token de autenticación. \\
\hline
RF-003 & Gestión de Sesiones & El sistema debe mantener sesiones activas mediante tokens JWT con tiempo de expiración configurable. \\
\hline
\rowcolor{rowgray}
RF-004 & Roles y Permisos & El sistema debe implementar un sistema de roles jerárquico (Super Admin \textgreater\ Admin \textgreater Comité \textgreater Delegado) con permisos específicos para cada rol. \\
\hline
RF-005 & Recuperación de Contraseña & El sistema debe permitir la recuperación de contraseña mediante enlace enviado al correo electrónico del usuario. \\
\hline
\rowcolor{rowgray}
RF-006 & Cambio de Contraseña & El sistema debe permitir a los usuarios cambiar su contraseña actual proporcionando la contraseña anterior. \\
\hline
RF-007 & Auditoría de Acciones & El sistema debe registrar todas las acciones críticas realizadas por los usuarios con timestamp, usuario y tipo de acción. \\
\hline
\rowcolor{rowgray}
RF-008 & Bloqueo por Intentos Fallidos & El sistema debe bloquear temporalmente cuentas después de 5 intentos fallidos de inicio de sesión. \\
\hline
RF-009 & Autenticación de Dos Factores & El sistema puede implementar autenticación de dos factores mediante código enviado por correo electrónico. \\
\hline
\rowcolor{rowgray}
RF-010 & Validación de Contraseñas & El sistema debe validar que las contraseñas cumplan con requisitos mínimos (8 caracteres, mayúsculas, minúsculas y números). \\
\hline
\caption{Requisitos Funcionales - Autenticación y Seguridad}
\end{longtable}

\subsection{RF-2. Módulo de Gestión de Usuarios}

Este módulo facilita la administración completa del personal involucrado en la gestión deportiva, respetando la jerarquía organizacional establecida. Permite crear y gestionar cuentas de diferentes niveles (administradores, comité deportivo, delegados de facultad), asignar roles y permisos específicos, y mantener un control detallado sobre el estado y las acciones de cada usuario dentro del sistema.

\begin{longtable}{|>{\centering\arraybackslash}p{1.5cm}|>{\raggedright\arraybackslash}p{3.2cm}|>{\raggedright\arraybackslash}p{9cm}|}
\hline
\rowcolor{headercolor}
\textcolor{white}{\textbf{ID}} & \textcolor{white}{\textbf{Requisito}} & \textcolor{white}{\textbf{Descripción}} \\
\hline
\endfirsthead

\multicolumn{3}{c}{\textit{Continuación de la tabla anterior}} \\
\hline
\rowcolor{headercolor}
\textcolor{white}{\textbf{ID}} & \textcolor{white}{\textbf{Requisito}} & \textcolor{white}{\textbf{Descripción}} \\
\hline
\endhead

\hline
\multicolumn{3}{r}{\textit{Continúa en la siguiente página}} \\
\endfoot

\hline
\endlastfoot

RF-011 & Creación de Administradores & El Super Administrador debe poder crear cuentas de administradores con credenciales iniciales. \\
\hline
\rowcolor{rowgray}
RF-012 & Creación de Comités & Los Administradores deben poder crear cuentas de miembros del comité deportivo. \\
\hline
RF-013 & Creación de Delegados & Los miembros del comité deben poder crear cuentas de delegados deportivos de facultades. \\
\hline
\rowcolor{rowgray}
RF-014 & Asignación de Roles & El sistema debe permitir asignar y modificar roles de usuarios según jerarquía establecida. \\
\hline
RF-015 & Gestión de Permisos & El sistema debe permitir personalizar permisos específicos para cada rol. \\
\hline
\rowcolor{rowgray}
RF-016 & Bloqueo/Desbloqueo de Cuentas & El sistema debe permitir bloquear y desbloquear cuentas de usuario manualmente. \\
\hline
RF-017 & Listado de Usuarios & El sistema debe mostrar un listado de todos los usuarios con filtros por rol, estado y facultad. \\
\hline
\rowcolor{rowgray}
RF-018 & Búsqueda de Usuarios & El sistema debe permitir buscar usuarios por nombre, correo electrónico o rol. \\
\hline
RF-019 & Edición de Datos & El sistema debe permitir editar información básica de usuarios (nombre, correo, teléfono). \\
\hline
\rowcolor{rowgray}
RF-020 & Eliminación Lógica & El sistema debe permitir desactivar usuarios sin eliminar su historial. \\
\hline
RF-021 & Historial de Actividad & El sistema debe mostrar el historial de acciones realizadas por cada usuario. \\
\hline
\caption{Requisitos Funcionales - Gestión de Usuarios}
\end{longtable}

\subsection{RF-3. Módulo de Gestión Institucional}

Este módulo centraliza la información de todas las facultades y escuelas participantes en los eventos deportivos. Permite registrar datos institucionales, gestionar identidad visual (logos y colores), asignar delegados responsables, y mantener un historial estadístico del desempeño deportivo de cada facultad. Es fundamental para la organización multi-institucional de los torneos.

\begin{longtable}{|>{\centering\arraybackslash}p{1.5cm}|>{\raggedright\arraybackslash}p{3.2cm}|>{\raggedright\arraybackslash}p{9cm}|}
\hline
\rowcolor{headercolor}
\textcolor{white}{\textbf{ID}} & \textcolor{white}{\textbf{Requisito}} & \textcolor{white}{\textbf{Descripción}} \\
\hline
\endfirsthead

\multicolumn{3}{c}{\textit{Continuación de la tabla anterior}} \\
\hline
\rowcolor{headercolor}
\textcolor{white}{\textbf{ID}} & \textcolor{white}{\textbf{Requisito}} & \textcolor{white}{\textbf{Descripción}} \\
\hline
\endhead

\hline
\multicolumn{3}{r}{\textit{Continúa en la siguiente página}} \\
\endfoot

\hline
\endlastfoot

RF-022 & Registro de Facultades & El sistema debe permitir registrar facultades con nombre, código y descripción. \\
\hline
\rowcolor{rowgray}
RF-023 & Registro de Escuelas & El sistema debe permitir registrar escuelas asociadas a facultades. \\
\hline
RF-024 & Logos Institucionales & El sistema debe permitir cargar y gestionar logos de facultades en formato PNG o JPG. \\
\hline
\rowcolor{rowgray}
RF-025 & Colores Institucionales & El sistema debe permitir definir colores representativos de cada facultad. \\
\hline
RF-026 & Asignación de Delegados & El sistema debe permitir asignar uno o varios delegados deportivos a cada facultad. \\
\hline
\rowcolor{rowgray}
RF-027 & Datos Institucionales & El sistema debe almacenar información de contacto, dirección y datos administrativos. \\
\hline
RF-028 & Activación/Desactivación & El sistema debe permitir activar o desactivar facultades según participación. \\
\hline
\rowcolor{rowgray}
RF-029 & Listado de Facultades & El sistema debe mostrar un listado completo de facultades con sus datos. \\
\hline
RF-030 & Estadísticas por Facultad & El sistema debe mostrar estadísticas históricas de participación y logros. \\
\hline
\caption{Requisitos Funcionales - Gestión Institucional}
\end{longtable}

\subsection{RF-4. Módulo de Gestión de Disciplinas Deportivas}

Este módulo permite configurar las características específicas de cada deporte que se practica en la institución. Define reglas de juego, sistemas de puntuación, criterios de desempate, tipos de competencia, y eventos registrables (goles, tarjetas, faltas). Su flexibilidad permite adaptar el sistema a cualquier disciplina deportiva, desde deportes tradicionales hasta nuevas modalidades que la institución desee incorporar.

\begin{longtable}{|>{\centering\arraybackslash}p{1.5cm}|>{\raggedright\arraybackslash}p{3.2cm}|>{\raggedright\arraybackslash}p{9cm}|}
\hline
\rowcolor{headercolor}
\textcolor{white}{\textbf{ID}} & \textcolor{white}{\textbf{Requisito}} & \textcolor{white}{\textbf{Descripción}} \\
\hline
\endfirsthead

\multicolumn{3}{c}{\textit{Continuación de la tabla anterior}} \\
\hline
\rowcolor{headercolor}
\textcolor{white}{\textbf{ID}} & \textcolor{white}{\textbf{Requisito}} & \textcolor{white}{\textbf{Descripción}} \\
\hline
\endhead

\hline
\multicolumn{3}{r}{\textit{Continúa en la siguiente página}} \\
\endfoot

\hline
\endlastfoot

RF-031 & Creación de Disciplinas & El sistema debe permitir crear disciplinas deportivas (fútbol, básquet, vóley, etc.). \\
\hline
\rowcolor{rowgray}
RF-032 & Configuración de Reglas & El sistema debe permitir definir reglas específicas por disciplina (duración, jugadores, etc.). \\
\hline
RF-033 & Motor de Puntuación & El sistema debe permitir configurar el sistema de puntuación (victoria, empate, derrota). \\
\hline
\rowcolor{rowgray}
RF-034 & Criterios de Desempate & El sistema debe permitir definir criterios de desempate (diferencia de goles, partidos directos, etc.). \\
\hline
RF-035 & Plantillas Deportivas & El sistema debe gestionar plantillas predefinidas para deportes comunes. \\
\hline
\rowcolor{rowgray}
RF-036 & Categorías por Disciplina & El sistema debe permitir definir categorías (varonil, femenil, mixto). \\
\hline
RF-037 & Tipos de Competencia & El sistema debe soportar liga, eliminatoria directa, grupos + eliminatoria. \\
\hline
\rowcolor{rowgray}
RF-038 & Configuración de Eventos & El sistema debe permitir definir eventos registrables (goles, tarjetas, faltas, etc.). \\
\hline
RF-039 & Configuración de Vestimenta & El sistema debe permitir definir requisitos de vestimenta y equipamiento por disciplina. \\
\hline
\rowcolor{rowgray}
RF-040 & Estado de Disciplinas & El sistema debe permitir activar o desactivar disciplinas según disponibilidad. \\
\hline
\caption{Requisitos Funcionales - Disciplinas Deportivas}
\end{longtable}

\subsection{RF-5. Módulo de Gestión de Torneos}

Este módulo es el núcleo organizativo del sistema, permitiendo planificar, configurar y administrar torneos de principio a fin. Gestiona el ciclo completo: desde la creación con fechas y parámetros específicos, pasando por el periodo de inscripciones, hasta la finalización y archivo histórico. Soporta múltiples formatos de competencia y permite publicar información para consulta pública.

\begin{longtable}{|>{\centering\arraybackslash}p{1.5cm}|>{\raggedright\arraybackslash}p{3.2cm}|>{\raggedright\arraybackslash}p{9cm}|}
\hline
\rowcolor{headercolor}
\textcolor{white}{\textbf{ID}} & \textcolor{white}{\textbf{Requisito}} & \textcolor{white}{\textbf{Descripción}} \\
\hline
\endfirsthead

\multicolumn{3}{c}{\textit{Continuación de la tabla anterior}} \\
\hline
\rowcolor{headercolor}
\textcolor{white}{\textbf{ID}} & \textcolor{white}{\textbf{Requisito}} & \textcolor{white}{\textbf{Descripción}} \\
\hline
\endhead

\hline
\multicolumn{3}{r}{\textit{Continúa en la siguiente página}} \\
\endfoot

\hline
\endlastfoot

RF-041 & Creación de Torneos & El sistema debe permitir crear torneos con nombre, descripción, fechas y disciplina asociada. \\
\hline
\rowcolor{rowgray}
RF-042 & Configuración de Fechas & El sistema debe permitir definir fecha de inicio, fin, inscripciones y reclamos. \\
\hline
RF-043 & Tipos de Torneo & El sistema debe soportar liga, eliminación directa, doble eliminación, grupos + playoff. \\
\hline
\rowcolor{rowgray}
RF-044 & Estados del Torneo & El sistema debe gestionar estados: planificación, inscripciones, en curso, finalizado, cancelado. \\
\hline
RF-045 & Publicación de Torneos & El sistema debe permitir publicar torneos para hacerlos visibles al público. \\
\hline
\rowcolor{rowgray}
RF-046 & Cierre de Inscripciones & El sistema debe cerrar automáticamente las inscripciones según fecha configurada. \\
\hline
RF-047 & Número de Equipos & El sistema debe permitir definir número mínimo y máximo de equipos participantes. \\
\hline
\rowcolor{rowgray}
RF-048 & Comité Organizador & El sistema debe permitir asignar miembros del comité responsables del torneo. \\
\hline
RF-049 & Configuración de Premios & El sistema debe permitir registrar información sobre premios y reconocimientos. \\
\hline
\rowcolor{rowgray}
RF-050 & Duplicación de Torneos & El sistema debe permitir duplicar configuración de torneos para nuevas ediciones. \\
\hline
RF-051 & Cancelación de Torneos & El sistema debe permitir cancelar torneos con registro de motivo. \\
\hline
\rowcolor{rowgray}
RF-052 & Historial de Torneos & El sistema debe mantener un historial completo de todos los torneos realizados. \\
\hline
\caption{Requisitos Funcionales - Gestión de Torneos}
\end{longtable}

\subsection{RF-6. Módulo de Equipos y Participantes}

Este módulo gestiona el registro y administración de equipos y jugadores participantes. Los delegados de cada facultad pueden inscribir sus equipos, agregar jugadores con toda su información personal y académica, designar capitanes, y adjuntar documentación de respaldo. El sistema valida automáticamente que los participantes cumplan con los requisitos establecidos (alumno activo, género según categoría, número de jugadores) y mantiene un historial completo de participación.

\begin{longtable}{|>{\centering\arraybackslash}p{1.5cm}|>{\raggedright\arraybackslash}p{3.2cm}|>{\raggedright\arraybackslash}p{9cm}|}
\hline
\rowcolor{headercolor}
\textcolor{white}{\textbf{ID}} & \textcolor{white}{\textbf{Requisito}} & \textcolor{white}{\textbf{Descripción}} \\
\hline
\endfirsthead

\multicolumn{3}{c}{\textit{Continuación de la tabla anterior}} \\
\hline
\rowcolor{headercolor}
\textcolor{white}{\textbf{ID}} & \textcolor{white}{\textbf{Requisito}} & \textcolor{white}{\textbf{Descripción}} \\
\hline
\endhead

\hline
\multicolumn{3}{r}{\textit{Continúa en la siguiente página}} \\
\endfoot

\hline
\endlastfoot

RF-053 & Registro de Equipos & El sistema debe permitir a delegados registrar equipos de su facultad para torneos. \\
\hline
\rowcolor{rowgray}
RF-054 & Inscripción de Jugadores & El sistema debe permitir inscribir jugadores con datos personales completos. \\
\hline
RF-055 & Validación de Alumnos & El sistema debe validar que los jugadores sean alumnos activos matriculados. \\
\hline
\rowcolor{rowgray}
RF-056 & Asignación de Capitán & El sistema debe permitir designar un capitán por equipo. \\
\hline
RF-057 & Gestión de Delegados & El sistema debe asignar automáticamente el delegado de facultad como responsable. \\
\hline
\rowcolor{rowgray}
RF-058 & Validación Académica & El sistema debe verificar que los jugadores estén matriculados en el semestre actual. \\
\hline
RF-059 & Restricciones de Género & El sistema debe validar que los jugadores cumplan requisitos de género. \\
\hline
\rowcolor{rowgray}
RF-060 & Número de Jugadores & El sistema debe validar número mínimo y máximo de jugadores según disciplina. \\
\hline
RF-061 & Fotografía de Jugadores & El sistema debe permitir cargar fotografías para identificación. \\
\hline
\rowcolor{rowgray}
RF-062 & Documentación & El sistema debe permitir adjuntar documentos de respaldo (credencial, DNI). \\
\hline
RF-063 & Estado de Jugadores & El sistema debe gestionar estados: activo, suspendido, lesionado, inactivo. \\
\hline
\rowcolor{rowgray}
RF-064 & Transferencia de Jugadores & El sistema debe permitir transferencias entre equipos bajo condiciones. \\
\hline
RF-065 & Listado de Equipos & El sistema debe mostrar todos los equipos inscritos en un torneo. \\
\hline
\rowcolor{rowgray}
RF-066 & Búsqueda de Jugadores & El sistema debe permitir buscar jugadores por nombre, código o facultad. \\
\hline
RF-067 & Historial de Participación & El sistema debe mantener historial de participación en torneos anteriores. \\
\hline
\caption{Requisitos Funcionales - Equipos y Participantes}
\end{longtable}

\subsection{RF-7. Módulo de Escenarios y Horarios}

Este módulo administra toda la infraestructura deportiva de la institución: canchas, gimnasios, y demás escenarios donde se realizan las competencias. Permite registrar características técnicas de cada instalación (tipo de superficie, capacidad, equipamiento), configurar horarios de disponibilidad, gestionar mantenimiento, y detectar automáticamente conflictos de agenda. Es esencial para la optimización del uso de recursos deportivos.

\begin{longtable}{|>{\centering\arraybackslash}p{1.5cm}|>{\raggedright\arraybackslash}p{3.2cm}|>{\raggedright\arraybackslash}p{9cm}|}
\hline
\rowcolor{headercolor}
\textcolor{white}{\textbf{ID}} & \textcolor{white}{\textbf{Requisito}} & \textcolor{white}{\textbf{Descripción}} \\
\hline
\endfirsthead

\multicolumn{3}{c}{\textit{Continuación de la tabla anterior}} \\
\hline
\rowcolor{headercolor}
\textcolor{white}{\textbf{ID}} & \textcolor{white}{\textbf{Requisito}} & \textcolor{white}{\textbf{Descripción}} \\
\hline
\endhead

\hline
\multicolumn{3}{r}{\textit{Continúa en la siguiente página}} \\
\endfoot

\hline
\endlastfoot

RF-068 & Registro de Canchas & El sistema debe permitir registrar canchas/escenarios con nombre, ubicación y tipo. \\
\hline
\rowcolor{rowgray}
RF-069 & Horarios Disponibles & El sistema debe permitir configurar horarios de disponibilidad de cada cancha. \\
\hline
RF-070 & Capacidad de Escenarios & El sistema debe registrar capacidad máxima de espectadores por cancha. \\
\hline
\rowcolor{rowgray}
RF-071 & Tipo de Superficie & El sistema debe registrar tipo de superficie (césped, sintético, cemento, parquet). \\
\hline
RF-072 & Conflictos de Agenda & El sistema debe detectar automáticamente conflictos al asignar canchas. \\
\hline
\rowcolor{rowgray}
RF-073 & Reserva de Canchas & El sistema debe permitir reservar canchas para partidos específicos. \\
\hline
RF-074 & Mantenimiento & El sistema debe permitir marcar canchas como no disponibles por mantenimiento. \\
\hline
\rowcolor{rowgray}
RF-075 & Características & El sistema debe registrar características adicionales (iluminación, gradas, vestidores). \\
\hline
RF-076 & Mapa de Ubicación & El sistema debe permitir asociar coordenadas GPS o mapa de ubicación. \\
\hline
\rowcolor{rowgray}
RF-077 & Calendario de Uso & El sistema debe mostrar calendario visual del uso de cada cancha. \\
\hline
\caption{Requisitos Funcionales - Escenarios y Horarios}
\end{longtable}

\subsection{RF-8. Módulo de Generación de Fixture}

Este módulo automatiza la creación del calendario de partidos mediante algoritmos inteligentes que consideran múltiples variables: tipo de torneo (liga, eliminación, grupos), disponibilidad de canchas, horarios, restricciones de equipos, y distribución equitativa de fechas. Reduce significativamente el tiempo de planificación y minimiza conflictos, permitiendo también reprogramaciones manuales cuando sea necesario.

\begin{longtable}{|>{\centering\arraybackslash}p{1.5cm}|>{\raggedright\arraybackslash}p{3.2cm}|>{\raggedright\arraybackslash}p{9cm}|}
\hline
\rowcolor{headercolor}
\textcolor{white}{\textbf{ID}} & \textcolor{white}{\textbf{Requisito}} & \textcolor{white}{\textbf{Descripción}} \\
\hline
\endfirsthead

\multicolumn{3}{c}{\textit{Continuación de la tabla anterior}} \\
\hline
\rowcolor{headercolor}
\textcolor{white}{\textbf{ID}} & \textcolor{white}{\textbf{Requisito}} & \textcolor{white}{\textbf{Descripción}} \\
\hline
\endhead

\hline
\multicolumn{3}{r}{\textit{Continúa en la siguiente página}} \\
\endfoot

\hline
\endlastfoot

RF-079 & Generación Automática & El sistema debe generar automáticamente el calendario de partidos. \\
\hline
\rowcolor{rowgray}
RF-080 & Algoritmo Round-Robin & El sistema debe implementar algoritmo Round-Robin para torneos tipo liga. \\
\hline
RF-081 & Algoritmo de Eliminación & El sistema debe generar llaves de eliminación directa. \\
\hline
\rowcolor{rowgray}
RF-082 & Restricciones de Horario & El sistema debe considerar horarios disponibles de canchas. \\
\hline
RF-083 & Restricciones de Equipos & El sistema debe considerar restricciones de disponibilidad de equipos. \\
\hline
\rowcolor{rowgray}
RF-084 & Distribución Equitativa & El sistema debe distribuir equitativamente partidos entre días y horarios. \\
\hline
RF-085 & Reprogramación & El sistema debe permitir reprogramar partidos individuales manualmente. \\
\hline
\rowcolor{rowgray}
RF-086 & Validación de Conflictos & El sistema debe validar que no haya equipos jugando simultáneamente. \\
\hline
RF-087 & Fechas de Descanso & El sistema debe permitir configurar fechas sin partidos (feriados, exámenes). \\
\hline
\rowcolor{rowgray}
RF-088 & Publicación del Fixture & El sistema debe permitir publicar el fixture completo para consulta pública. \\
\hline
RF-089 & Notificación de Fixture & El sistema debe notificar automáticamente a equipos cuando se publique. \\
\hline
\rowcolor{rowgray}
RF-090 & Fechas Dobles & El sistema debe soportar generación de fechas dobles para acelerar torneos. \\
\hline
RF-091 & Fixture por Grupos & El sistema debe generar fixtures separados para torneos con fase de grupos. \\
\hline
\caption{Requisitos Funcionales - Generación de Fixture}
\end{longtable}

\subsection{RF-9. Módulo de Ejecución de Partidos (Acta Digital)}

Este módulo digitaliza completamente el proceso de registro de eventos durante los partidos, reemplazando las actas físicas tradicionales. Permite registrar en tiempo real goles, tarjetas, sustituciones, lesiones y cualquier evento relevante. Incluye firma digital de delegados y árbitros para validación oficial, cronómetro integrado, y gestión automática de walkovers. El acta digital generada sirve como documento oficial del partido.

\begin{longtable}{|>{\centering\arraybackslash}p{1.5cm}|>{\raggedright\arraybackslash}p{3.2cm}|>{\raggedright\arraybackslash}p{9cm}|}
\hline
\rowcolor{headercolor}
\textcolor{white}{\textbf{ID}} & \textcolor{white}{\textbf{Requisito}} & \textcolor{white}{\textbf{Descripción}} \\
\hline
\endfirsthead

\multicolumn{3}{c}{\textit{Continuación de la tabla anterior}} \\
\hline
\rowcolor{headercolor}
\textcolor{white}{\textbf{ID}} & \textcolor{white}{\textbf{Requisito}} & \textcolor{white}{\textbf{Descripción}} \\
\hline
\endhead

\hline
\multicolumn{3}{r}{\textit{Continúa en la siguiente página}} \\
\endfoot

\hline
\endlastfoot

RF-092 & Registro en Tiempo Real & El sistema debe permitir registrar eventos durante la ejecución del partido. \\
\hline
\rowcolor{rowgray}
RF-093 & Cronómetro del Partido & El sistema debe incluir cronómetro para controlar tiempo de juego. \\
\hline
RF-094 & Registro de Goles & El sistema debe permitir registrar goles con jugador, minuto y tipo. \\
\hline
\rowcolor{rowgray}
RF-095 & Registro de Tarjetas & El sistema debe permitir registrar tarjetas amarillas y rojas con jugador y minuto. \\
\hline
RF-096 & Registro de Cambios & El sistema debe permitir registrar sustituciones de jugadores. \\
\hline
\rowcolor{rowgray}
RF-097 & Registro de Faltas & El sistema debe permitir registrar faltas cometidas según configuración. \\
\hline
RF-098 & Firma Digital & El sistema debe permitir firma digital de delegados y árbitros al finalizar. \\
\hline
\rowcolor{rowgray}
RF-099 & Estados del Partido & El sistema debe gestionar estados: programado, en curso, finalizado, suspendido, walkover. \\
\hline
RF-100 & Alineación Inicial & El sistema debe permitir registrar jugadores titulares y suplentes. \\
\hline
\rowcolor{rowgray}
RF-101 & Observaciones & El sistema debe permitir al árbitro agregar observaciones en el acta. \\
\hline
RF-102 & Registro de Lesiones & El sistema debe permitir registrar lesiones ocurridas durante el partido. \\
\hline
\rowcolor{rowgray}
RF-103 & Tiempo Extra & El sistema debe soportar registro de tiempo extra y penales según disciplina. \\
\hline
RF-104 & Verificación de Asistencia & El sistema debe verificar asistencia de jugadores antes de permitir participación. \\
\hline
\rowcolor{rowgray}
RF-105 & Walkover Automático & El sistema debe registrar walkover si un equipo no se presenta (15 min). \\
\hline
RF-106 & Edición de Acta & El sistema debe permitir editar actas dentro de un periodo configurable. \\
\hline
RF-107 & Notificación de Inicio & El sistema debe notificar a los equipos 30 minutos antes del inicio del partido. \\
\hline
\caption{Requisitos Funcionales - Ejecución de Partidos}
\end{longtable}

\subsection{RF-10. Módulo de Resultados y Tablas}

Este módulo procesa automáticamente todos los datos de los partidos para generar tablas de posiciones, estadísticas de equipos y jugadores, y clasificaciones a siguientes fases. Aplica los criterios de desempate configurados, calcula diferencias de goles, mantiene tablas de goleadores, y genera estadísticas disciplinarias. Toda la información se actualiza en tiempo real conforme se finalizan los partidos.

\begin{longtable}{|>{\centering\arraybackslash}p{1.5cm}|>{\raggedright\arraybackslash}p{3.2cm}|>{\raggedright\arraybackslash}p{9cm}|}
\hline
\rowcolor{headercolor}
\textcolor{white}{\textbf{ID}} & \textcolor{white}{\textbf{Requisito}} & \textcolor{white}{\textbf{Descripción}} \\
\hline
\endfirsthead

\multicolumn{3}{c}{\textit{Continuación de la tabla anterior}} \\
\hline
\rowcolor{headercolor}
\textcolor{white}{\textbf{ID}} & \textcolor{white}{\textbf{Requisito}} & \textcolor{white}{\textbf{Descripción}} \\
\hline
\endhead

\hline
\multicolumn{3}{r}{\textit{Continúa en la siguiente página}} \\
\endfoot

\hline
\endlastfoot

RF-108 & Cálculo Automático & El sistema debe calcular automáticamente puntos y estadísticas. \\
\hline
\rowcolor{rowgray}
RF-109 & Tabla de Posiciones & El sistema debe generar y actualizar automáticamente tablas de posiciones. \\
\hline
RF-110 & Criterios de Desempate & El sistema debe aplicar automáticamente criterios de desempate configurados. \\
\hline
\rowcolor{rowgray}
RF-111 & Clasificación & El sistema debe determinar automáticamente equipos clasificados a siguientes fases. \\
\hline
RF-112 & Estadísticas Generales & El sistema debe calcular partidos jugados, ganados, empatados, perdidos. \\
\hline
\rowcolor{rowgray}
RF-113 & Tabla de Goleadores & El sistema debe generar tabla de goleadores del torneo. \\
\hline
RF-114 & Estadísticas de Tarjetas & El sistema debe generar estadísticas de tarjetas por equipo y jugador. \\
\hline
\rowcolor{rowgray}
RF-115 & Diferencia de Goles & El sistema debe calcular diferencia de goles para cada equipo. \\
\hline
RF-116 & Racha de Resultados & El sistema debe mostrar últimos 5 resultados de cada equipo. \\
\hline
\rowcolor{rowgray}
RF-117 & Resultados Históricos & El sistema debe mantener historial completo de resultados anteriores. \\
\hline
RF-118 & Comparación de Equipos & El sistema debe permitir comparar estadísticas entre dos equipos. \\
\hline
\rowcolor{rowgray}
RF-119 & Predicciones & El sistema puede calcular probabilidades de resultados basadas en historial. \\
\hline
\caption{Requisitos Funcionales - Resultados y Tablas}
\end{longtable}

\subsection{RF-11. Módulo de Justicia Deportiva}

Este módulo facilita el proceso formal de reclamos y sanciones deportivas. Permite presentar reclamos con evidencias adjuntas, gestionar el proceso de revisión por parte del tribunal deportivo, notificar resoluciones, y aplicar sanciones tanto automáticas (por acumulación de tarjetas) como manuales (impuestas por el comité). Controla automáticamente el cumplimiento de suspensiones y mantiene un historial disciplinario completo de jugadores y equipos.

\begin{longtable}{|>{\centering\arraybackslash}p{1.5cm}|>{\raggedright\arraybackslash}p{3.2cm}|>{\raggedright\arraybackslash}p{9cm}|}
\hline
\rowcolor{headercolor}
\textcolor{white}{\textbf{ID}} & \textcolor{white}{\textbf{Requisito}} & \textcolor{white}{\textbf{Descripción}} \\
\hline
\endfirsthead

\multicolumn{3}{c}{\textit{Continuación de la tabla anterior}} \\
\hline
\rowcolor{headercolor}
\textcolor{white}{\textbf{ID}} & \textcolor{white}{\textbf{Requisito}} & \textcolor{white}{\textbf{Descripción}} \\
\hline
\endhead

\hline
\multicolumn{3}{r}{\textit{Continúa en la siguiente página}} \\
\endfoot

\hline
\endlastfoot

RF-120 & Presentación de Reclamos & El sistema debe permitir presentar reclamos formales con descripción detallada. \\
\hline
\rowcolor{rowgray}
RF-121 & Adjuntar Evidencias & El sistema debe permitir adjuntar evidencias (fotos, videos, documentos). \\
\hline
RF-122 & Tribunal Deportivo & El sistema debe facilitar proceso de revisión de reclamos por comité. \\
\hline
\rowcolor{rowgray}
RF-123 & Estados de Reclamos & El sistema debe gestionar estados: presentado, en revisión, resuelto, rechazado. \\
\hline
RF-124 & Notificación de Resoluciones & El sistema debe notificar a las partes sobre resoluciones. \\
\hline
\rowcolor{rowgray}
RF-125 & Sanciones Automáticas & El sistema debe aplicar automáticamente sanciones por acumulación de tarjetas. \\
\hline
RF-126 & Sanciones Manuales & El sistema debe permitir registrar sanciones impuestas por el tribunal. \\
\hline
\rowcolor{rowgray}
RF-127 & Suspensión de Jugadores & El sistema debe impedir participación de jugadores suspendidos. \\
\hline
RF-128 & Suspensión de Equipos & El sistema debe permitir suspender equipos completos bajo condiciones graves. \\
\hline
\rowcolor{rowgray}
RF-129 & Historial Disciplinario & El sistema debe mantener historial de sanciones y reclamos. \\
\hline
RF-130 & Apelaciones & El sistema debe permitir presentar apelaciones a resoluciones. \\
\hline
\rowcolor{rowgray}
RF-131 & Cumplimiento de Sanciones & El sistema debe controlar automáticamente cumplimiento de suspensiones. \\
\hline
RF-132 & Puntos Disciplinarios & El sistema puede implementar sistema de puntos disciplinarios acumulables. \\
\hline
\caption{Requisitos Funcionales - Justicia Deportiva}
\end{longtable}

\subsection{RF-12. Módulo de Notificaciones}

Este módulo mantiene informados a todos los involucrados mediante un sistema automatizado de comunicaciones. Envía notificaciones por correo electrónico y dentro de la plataforma sobre eventos importantes: publicación de fixtures, reprogramaciones, recordatorios de partidos, resultados oficiales, sanciones, y resoluciones de reclamos. Los usuarios pueden personalizar sus preferencias de notificación según sus necesidades.

\begin{longtable}{|>{\centering\arraybackslash}p{1.5cm}|>{\raggedright\arraybackslash}p{3.2cm}|>{\raggedright\arraybackslash}p{9cm}|}
\hline
\rowcolor{headercolor}
\textcolor{white}{\textbf{ID}} & \textcolor{white}{\textbf{Requisito}} & \textcolor{white}{\textbf{Descripción}} \\
\hline
\endfirsthead

\multicolumn{3}{c}{\textit{Continuación de la tabla anterior}} \\
\hline
\rowcolor{headercolor}
\textcolor{white}{\textbf{ID}} & \textcolor{white}{\textbf{Requisito}} & \textcolor{white}{\textbf{Descripción}} \\
\hline
\endhead

\hline
\multicolumn{3}{r}{\textit{Continúa en la siguiente página}} \\
\endfoot

\hline
\endlastfoot

RF-133 & Correo Electrónico & El sistema debe enviar notificaciones importantes por correo. \\
\hline
\rowcolor{rowgray}
RF-134 & Notificaciones Internas & El sistema debe mostrar notificaciones dentro de la plataforma. \\
\hline
RF-135 & Alertas por Cambios & El sistema debe notificar cuando se reprograme un partido. \\
\hline
\rowcolor{rowgray}
RF-136 & Recordatorios & El sistema debe enviar recordatorios 24 horas antes de cada partido. \\
\hline
RF-137 & Notificación de Resultados & El sistema debe notificar cuando se publiquen resultados oficiales. \\
\hline
\rowcolor{rowgray}
RF-138 & Confirmación de Inscripción & El sistema debe notificar confirmación de inscripción de equipos. \\
\hline
RF-139 & Notificación de Sanciones & El sistema debe notificar inmediatamente sobre sanciones aplicadas. \\
\hline
\rowcolor{rowgray}
RF-140 & Notificación de Reclamos & El sistema debe notificar a partes involucradas cuando se presente reclamo. \\
\hline
RF-141 & Preferencias & El sistema debe permitir configurar qué notificaciones desean recibir. \\
\hline
\rowcolor{rowgray}
RF-142 & Historial & El sistema debe mantener historial de todas las notificaciones enviadas. \\
\hline
RF-143 & Notificaciones Push & El sistema puede implementar notificaciones push para móviles. \\
\hline
\rowcolor{rowgray}
RF-144 & Notificaciones de Mantenimiento & El sistema debe notificar con anticipación mantenimientos. \\
\hline
RF-145 & Notificación SMS & El sistema puede enviar notificaciones críticas por mensaje de texto. \\
\hline
\rowcolor{rowgray}
RF-146 & Notificaciones de Logros & El sistema puede notificar cuando un equipo alcance hitos importantes. \\
\hline
\caption{Requisitos Funcionales - Notificaciones}
\end{longtable}

\subsection{RF-13. Módulo de Portal Público}

Este módulo proporciona acceso abierto a la información deportiva para toda la comunidad universitaria sin necesidad de autenticación. Muestra resultados en vivo, fixtures completos, tablas de posiciones, estadísticas de jugadores, noticias oficiales y galerías de fotos. Su diseño responsive garantiza una experiencia óptima en cualquier dispositivo, fomentando el seguimiento y apoyo a los equipos institucionales.

\begin{longtable}{|>{\centering\arraybackslash}p{1.5cm}|>{\raggedright\arraybackslash}p{3.2cm}|>{\raggedright\arraybackslash}p{9cm}|}
\hline
\rowcolor{headercolor}
\textcolor{white}{\textbf{ID}} & \textcolor{white}{\textbf{Requisito}} & \textcolor{white}{\textbf{Descripción}} \\
\hline
\endfirsthead

\multicolumn{3}{c}{\textit{Continuación de la tabla anterior}} \\
\hline
\rowcolor{headercolor}
\textcolor{white}{\textbf{ID}} & \textcolor{white}{\textbf{Requisito}} & \textcolor{white}{\textbf{Descripción}} \\
\hline
\endhead

\hline
\multicolumn{3}{r}{\textit{Continúa en la siguiente página}} \\
\endfoot

\hline
\endlastfoot

RF-147 & Resultados en Vivo & El portal debe mostrar resultados de partidos en curso en tiempo real. \\
\hline
\rowcolor{rowgray}
RF-148 & Consulta de Fixture & El portal debe permitir consultar fixture completo sin autenticación. \\
\hline
RF-149 & Visualización de Tablas & El portal debe mostrar tablas de posiciones actualizadas. \\
\hline
\rowcolor{rowgray}
RF-150 & Sección de Noticias & El portal debe incluir sección de noticias y comunicados oficiales. \\
\hline
RF-151 & Galería de Fotos & El portal puede incluir galería de fotos de partidos y eventos. \\
\hline
\rowcolor{rowgray}
RF-152 & Estadísticas Públicas & El portal debe mostrar estadísticas generales, goleadores y tablas históricas. \\
\hline
RF-153 & Calendario de Eventos & El portal debe mostrar calendario con próximos partidos destacados. \\
\hline
\rowcolor{rowgray}
RF-154 & Búsqueda de Equipos & El portal debe permitir buscar y consultar información de equipos. \\
\hline
RF-155 & Búsqueda de Jugadores & El portal debe permitir consultar estadísticas individuales de jugadores. \\
\hline
\rowcolor{rowgray}
RF-156 & Diseño Responsive & El portal debe ser completamente responsive y funcional en móviles. \\
\hline
RF-157 & Compartir en Redes & El portal debe permitir compartir resultados y noticias en redes sociales. \\
\hline
\rowcolor{rowgray}
RF-158 & Descarga de Fixture & El portal debe permitir descargar fixture en PDF o imagen. \\
\hline
\caption{Requisitos Funcionales - Portal Público}
\end{longtable}

\subsection{RF-14. Módulo de Reportes y Exportación}

Este módulo facilita la generación de documentación oficial y análisis de datos mediante reportes profesionales en diversos formatos. Produce actas oficiales de partidos, tablas de posiciones, fixtures, estadísticas de rendimiento, reportes ejecutivos para autoridades, certificados de participación, y análisis comparativos. Los formatos de exportación incluyen PDF para documentos formales y Excel/CSV para análisis de datos.

\begin{longtable}{|>{\centering\arraybackslash}p{1.5cm}|>{\raggedright\arraybackslash}p{3.2cm}|>{\raggedright\arraybackslash}p{9cm}|}
\hline
\rowcolor{headercolor}
\textcolor{white}{\textbf{ID}} & \textcolor{white}{\textbf{Requisito}} & \textcolor{white}{\textbf{Descripción}} \\
\hline
\endfirsthead

\multicolumn{3}{c}{\textit{Continuación de la tabla anterior}} \\
\hline
\rowcolor{headercolor}
\textcolor{white}{\textbf{ID}} & \textcolor{white}{\textbf{Requisito}} & \textcolor{white}{\textbf{Descripción}} \\
\hline
\endhead

\hline
\multicolumn{3}{r}{\textit{Continúa en la siguiente página}} \\
\endfoot

\hline
\endlastfoot

RF-159 & Generación de PDFs & El sistema debe generar reportes en formato PDF con diseño profesional. \\
\hline
\rowcolor{rowgray}
RF-160 & Exportación a Excel & El sistema debe permitir exportar tablas y estadísticas a Excel/CSV. \\
\hline
RF-161 & Actas de Partido & El sistema debe generar actas oficiales en PDF con firmas digitales. \\
\hline
\rowcolor{rowgray}
RF-162 & Tabla de Posiciones & El sistema debe generar reportes PDF de tablas finales e intermedias. \\
\hline
RF-163 & Fixture Completo & El sistema debe generar calendario completo del torneo en PDF. \\
\hline
\rowcolor{rowgray}
RF-164 & Estadísticas Individuales & El sistema debe generar reportes de rendimiento de jugadores. \\
\hline
RF-165 & Estadísticas por Equipo & El sistema debe generar análisis completo del desempeño de equipos. \\
\hline
\rowcolor{rowgray}
RF-166 & Reportes Institucionales & El sistema debe generar reportes ejecutivos para autoridades. \\
\hline
RF-167 & Reporte de Participación & El sistema debe generar reportes de participación por facultad y semestre. \\
\hline
\rowcolor{rowgray}
RF-168 & Uso de Instalaciones & El sistema debe generar reportes de ocupación de canchas. \\
\hline
RF-169 & Reporte de Sanciones & El sistema debe generar reportes de historial disciplinario por torneo. \\
\hline
\rowcolor{rowgray}
RF-170 & Certificados & El sistema debe generar certificados digitales de participación para jugadores. \\
\hline
RF-171 & Ranking de Goleadores & El sistema debe generar ranking oficial de máximos goleadores. \\
\hline
\rowcolor{rowgray}
RF-172 & Reporte Comparativo & El sistema debe generar reportes comparativos entre torneos o semestres. \\
\hline
RF-169 & Programación de Reportes & El sistema debe permitir programar generación automática de reportes. \\
\hline
\rowcolor{rowgray}
RF-173 & Exportación de Actas Históricas & El sistema debe permitir exportar todas las actas de un torneo en un solo archivo. \\
\hline
\caption{Requisitos Funcionales - Reportes y Exportación}
\end{longtable}

\subsection{RF-15. Módulo de Analítica y KPIs}

Este módulo proporciona inteligencia de negocio mediante el análisis de datos históricos y en tiempo real. Presenta dashboards con indicadores clave de desempeño (KPIs) como tasas de participación por facultad, ocupación de instalaciones, popularidad de disciplinas, cumplimiento de horarios, y tendencias históricas. Esta información permite a los administradores tomar decisiones informadas para mejorar continuamente la gestión deportiva.

\begin{longtable}{|>{\centering\arraybackslash}p{1.5cm}|>{\raggedright\arraybackslash}p{3.5cm}|>{\raggedright\arraybackslash}p{9cm}|}
\hline
\rowcolor{headercolor}
\textcolor{white}{\textbf{ID}} & \textcolor{white}{\textbf{Requisito}} & \textcolor{white}{\textbf{Descripción}} \\
\hline
\endfirsthead

\multicolumn{3}{c}{\textit{Continuación de la tabla anterior}} \\
\hline
\rowcolor{headercolor}
\textcolor{white}{\textbf{ID}} & \textcolor{white}{\textbf{Requisito}} & \textcolor{white}{\textbf{Descripción}} \\
\hline
\endhead

\hline
\multicolumn{3}{r}{\textit{Continúa en la siguiente página}} \\
\endfoot

\hline
\endlastfoot

RF-174 & Dashboard de Indicadores & El sistema debe mostrar dashboard con KPIs principales en tiempo real. \\
\hline
\rowcolor{rowgray}
RF-175 & Participación por Facultad & El sistema debe calcular tasas de participación por facultad. \\
\hline
RF-176 & Uso de Canchas & El sistema debe mostrar estadísticas de ocupación de instalaciones. \\
\hline
\rowcolor{rowgray}
RF-177 & Rendimiento del Torneo & El sistema debe evaluar éxito de torneos según participación y finalización. \\
\hline
RF-178 & Tendencias Históricas & El sistema debe mostrar gráficas de evolución de participación. \\
\hline
\rowcolor{rowgray}
RF-179 & Análisis de Asistencia & El sistema puede registrar y analizar asistencia de público. \\
\hline
RF-180 & Disciplinas Populares & El sistema debe identificar disciplinas con mayor demanda. \\
\hline
\rowcolor{rowgray}
RF-181 & Tiempo Promedio & El sistema debe calcular duraciones promedio de partidos por disciplina. \\
\hline
RF-182 & Tasa de Reclamos & El sistema debe analizar frecuencia y tipos de reclamos presentados. \\
\hline
\rowcolor{rowgray}
RF-183 & Cumplimiento de Horarios & El sistema debe medir puntualidad en inicio de partidos. \\
\hline
RF-184 & Exportación Analítica & El sistema debe permitir exportar datos para análisis externo. \\
\hline
\rowcolor{rowgray}
RF-185 & Comparación Interanual & El sistema debe permitir comparar métricas entre años académicos. \\
\hline
\caption{Requisitos Funcionales - Analítica y KPIs}
\end{longtable}

\subsection{RF-16. Módulo de Configuración General}

Este módulo centraliza todos los parámetros y configuraciones globales del sistema. Permite definir periodos académicos, límites operativos, timeouts de sesión, plantillas de correos, configuración de respaldos automáticos, personalización de interfaz, integración con sistemas externos, y ajustes regionales (zona horaria, idioma, formato de fechas). Facilita la adaptación del sistema a las necesidades específicas de cada institución.

\begin{longtable}{|>{\centering\arraybackslash}p{1.5cm}|>{\raggedright\arraybackslash}p{3.5cm}|>{\raggedright\arraybackslash}p{9cm}|}
\hline
\rowcolor{headercolor}
\textcolor{white}{\textbf{ID}} & \textcolor{white}{\textbf{Requisito}} & \textcolor{white}{\textbf{Descripción}} \\
\hline
\endfirsthead

\multicolumn{3}{c}{\textit{Continuación de la tabla anterior}} \\
\hline
\rowcolor{headercolor}
\textcolor{white}{\textbf{ID}} & \textcolor{white}{\textbf{Requisito}} & \textcolor{white}{\textbf{Descripción}} \\
\hline
\endhead

\hline
\multicolumn{3}{r}{\textit{Continúa en la siguiente página}} \\
\endfoot

\hline
\endlastfoot

RF-186 & Parámetros del Sistema & El sistema debe permitir configurar parámetros globales. \\
\hline
\rowcolor{rowgray}
RF-187 & Semestre Académico & El sistema debe permitir configurar periodos académicos activos. \\
\hline
RF-188 & Límites Globales & El sistema debe permitir establecer límites máximos (equipos, jugadores, torneos). \\
\hline
\rowcolor{rowgray}
RF-189 & Configuración de Tiempos & El sistema debe permitir configurar timeouts de sesión, plazos, etc. \\
\hline
RF-190 & Plantillas de Email & El sistema debe permitir personalizar plantillas de correos automáticos. \\
\hline
\rowcolor{rowgray}
RF-191 & Backup Automático & El sistema debe realizar respaldos automáticos de base de datos. \\
\hline
RF-192 & Restauración de Backup & El sistema debe permitir restaurar base de datos desde respaldos. \\
\hline
\rowcolor{rowgray}
RF-193 & Logs del Sistema & El sistema debe mantener logs detallados de errores y eventos críticos. \\
\hline
RF-194 & Roles Personalizados & El sistema debe permitir crear roles con permisos específicos. \\
\hline
\rowcolor{rowgray}
RF-195 & Mantenimiento Programado & El sistema debe permitir programar ventanas de mantenimiento. \\
\hline
RF-196 & Gestión de Versiones & El sistema debe llevar control de versiones y cambios implementados. \\
\hline
\rowcolor{rowgray}
RF-197 & Configuración de Integraciones & El sistema debe permitir configurar integraciones con sistemas externos. \\
\hline
RF-198 & Personalización de UI & El sistema debe permitir personalizar colores y logos institucionales. \\
\hline
\rowcolor{rowgray}
RF-199 & Configuración Regional & El sistema debe permitir configurar zona horaria, idioma y formato de fechas. \\
\hline
\caption{Requisitos Funcionales - Configuración General}
\end{longtable}

\newpage

\section{Requisitos No Funcionales del Producto}

Los requisitos no funcionales especifican los atributos de calidad, restricciones y propiedades transversales que debe poseer el Sistema para garantizar su correcto funcionamiento, mantenibilidad, seguridad y satisfacción de los usuarios finales. Estos requisitos son igualmente críticos que los funcionales para el éxito del producto y están alineados con el estándar ISO/IEC 25010 (SQuaRE).

\subsection{RNF-1. Rendimiento}

Estos requisitos garantizan que el sistema responda con la velocidad necesaria para proporcionar una experiencia de usuario fluida, incluso bajo condiciones de carga elevada durante eventos deportivos importantes.

\begin{longtable}{|>{\centering\arraybackslash}p{1.5cm}|>{\raggedright\arraybackslash}p{3.5cm}|>{\raggedright\arraybackslash}p{9cm}|}
\hline
\rowcolor{headercolor}
\textcolor{white}{\textbf{ID}} & \textcolor{white}{\textbf{Requisito}} & \textcolor{white}{\textbf{Descripción}} \\
\hline
\endfirsthead

\multicolumn{3}{c}{\textit{Continuación de la tabla anterior}} \\
\hline
\rowcolor{headercolor}
\textcolor{white}{\textbf{ID}} & \textcolor{white}{\textbf{Requisito}} & \textcolor{white}{\textbf{Descripción}} \\
\hline
\endhead

\hline
\multicolumn{3}{r}{\textit{Continúa en la siguiente página}} \\
\endfoot

\hline
\endlastfoot

RNF-001 & Tiempo de Respuesta & El sistema debe responder a solicitudes en menos de 2 segundos en condiciones normales. \\
\hline
\rowcolor{rowgray}
RNF-002 & Carga Concurrente & El sistema debe soportar al menos 500 usuarios concurrentes sin degradación. \\
\hline
RNF-003 & Tiempo de Carga & Las páginas del portal público deben cargar en menos de 3 segundos. \\
\hline
\rowcolor{rowgray}
RNF-004 & Procesamiento de Fixture & La generación de fixture para 20 equipos debe completarse en menos de 10 segundos. \\
\hline
RNF-005 & Actualización en Tiempo Real & Los marcadores en vivo deben actualizarse con latencia máxima de 5 segundos. \\
\hline
\caption{Requisitos No Funcionales - Rendimiento}
\end{longtable}

\subsection{RNF-2. Seguridad}

Estos requisitos establecen los mecanismos de protección necesarios para salvaguardar la información sensible, prevenir accesos no autorizados, y garantizar la integridad de los datos del sistema.

\begin{longtable}{|>{\centering\arraybackslash}p{1.5cm}|>{\raggedright\arraybackslash}p{3.5cm}|>{\raggedright\arraybackslash}p{9cm}|}
\hline
\rowcolor{headercolor}
\textcolor{white}{\textbf{ID}} & \textcolor{white}{\textbf{Requisito}} & \textcolor{white}{\textbf{Descripción}} \\
\hline
\endfirsthead

\multicolumn{3}{c}{\textit{Continuación de la tabla anterior}} \\
\hline
\rowcolor{headercolor}
\textcolor{white}{\textbf{ID}} & \textcolor{white}{\textbf{Requisito}} & \textcolor{white}{\textbf{Descripción}} \\
\hline
\endhead

\hline
\multicolumn{3}{r}{\textit{Continúa en la siguiente página}} \\
\endfoot

\hline
\endlastfoot

RNF-006 & Encriptación de Datos & Todas las comunicaciones deben utilizar HTTPS con certificado SSL/TLS válido. \\
\hline
\rowcolor{rowgray}
RNF-007 & Encriptación de Contraseñas & Las contraseñas deben almacenarse usando algoritmos seguros (bcrypt, Argon2). \\
\hline
RNF-008 & Protección SQL & El sistema debe implementar consultas parametrizadas para prevenir inyección SQL. \\
\hline
\rowcolor{rowgray}
RNF-009 & Protección CSRF & El sistema debe implementar tokens CSRF para prevenir falsificación de solicitudes. \\
\hline
RNF-010 & Validación de Entrada & Todos los datos de entrada deben ser validados y sanitizados en servidor. \\
\hline
\rowcolor{rowgray}
RNF-011 & Control de Acceso & El sistema debe implementar RBAC en todas las operaciones críticas. \\
\hline
RNF-012 & Auditoría de Seguridad & El sistema debe registrar intentos de acceso no autorizado. \\
\hline
\rowcolor{rowgray}
RNF-013 & Sesiones Seguras & Los tokens de sesión deben tener expiración y renovarse periódicamente. \\
\hline
RNF-014 & Protección de Archivos & Los archivos cargados deben ser validados por tipo y tamaño. \\
\hline
\rowcolor{rowgray}
RNF-015 & Backups Encriptados & Las copias de seguridad deben almacenarse encriptadas. \\
\hline
\caption{Requisitos No Funcionales - Seguridad}
\end{longtable}

\subsection{RNF-3. Usabilidad}

Estos requisitos aseguran que el sistema sea intuitivo, accesible y fácil de usar para todos los perfiles de usuario, minimizando la curva de aprendizaje y maximizando la eficiencia operativa.

\begin{longtable}{|>{\centering\arraybackslash}p{1.5cm}|>{\raggedright\arraybackslash}p{3.5cm}|>{\raggedright\arraybackslash}p{9cm}|}
\hline
\rowcolor{headercolor}
\textcolor{white}{\textbf{ID}} & \textcolor{white}{\textbf{Requisito}} & \textcolor{white}{\textbf{Descripción}} \\
\hline
\endfirsthead

\multicolumn{3}{c}{\textit{Continuación de la tabla anterior}} \\
\hline
\rowcolor{headercolor}
\textcolor{white}{\textbf{ID}} & \textcolor{white}{\textbf{Requisito}} & \textcolor{white}{\textbf{Descripción}} \\
\hline
\endhead

\hline
\multicolumn{3}{r}{\textit{Continúa en la siguiente página}} \\
\endfoot

\hline
\endlastfoot

RNF-016 & Interfaz Intuitiva & La interfaz no debe requerir más de 1 hora de capacitación para usuarios básicos. \\
\hline
\rowcolor{rowgray}
RNF-017 & Navegación Clara & Cualquier funcionalidad debe alcanzarse en máximo 3 clics. \\
\hline
RNF-018 & Diseño Responsive & La interfaz debe adaptarse a móviles, tablets y escritorio. \\
\hline
\rowcolor{rowgray}
RNF-019 & Accesibilidad & El sistema debe cumplir con estándares WCAG 2.1 nivel AA. \\
\hline
RNF-020 & Mensajes de Error & Los mensajes deben ser claros, específicos y ofrecer soluciones. \\
\hline
\rowcolor{rowgray}
RNF-021 & Retroalimentación Visual & El sistema debe proporcionar retroalimentación inmediata para acciones. \\
\hline
RNF-022 & Consistencia de Diseño & La interfaz debe mantener consistencia en colores, tipografía y patrones. \\
\hline
\rowcolor{rowgray}
RNF-023 & Ayuda Contextual & El sistema debe proporcionar tooltips en funcionalidades complejas. \\
\hline
\caption{Requisitos No Funcionales - Usabilidad}
\end{longtable}

\subsection{RNF-4. Disponibilidad}

Estos requisitos garantizan que el sistema esté operativo cuando los usuarios lo necesiten, con mecanismos de recuperación ante fallos y mantenimiento planificado que minimice interrupciones.

\begin{longtable}{|>{\centering\arraybackslash}p{1.5cm}|>{\raggedright\arraybackslash}p{3.5cm}|>{\raggedright\arraybackslash}p{9cm}|}
\hline
\rowcolor{headercolor}
\textcolor{white}{\textbf{ID}} & \textcolor{white}{\textbf{Requisito}} & \textcolor{white}{\textbf{Descripción}} \\
\hline
\endfirsthead

\multicolumn{3}{c}{\textit{Continuación de la tabla anterior}} \\
\hline
\rowcolor{headercolor}
\textcolor{white}{\textbf{ID}} & \textcolor{white}{\textbf{Requisito}} & \textcolor{white}{\textbf{Descripción}} \\
\hline
\endhead

\hline
\multicolumn{3}{r}{\textit{Continúa en la siguiente página}} \\
\endfoot

\hline
\endlastfoot

RNF-024 & Uptime & El sistema debe mantener disponibilidad del 99.5\% mensual (máximo 3.6 horas de inactividad). \\
\hline
\rowcolor{rowgray}
RNF-025 & Recuperación ante Fallos & El sistema debe recuperarse de fallos menores en menos de 5 minutos. \\
\hline
RNF-026 & Mantenimiento & Las ventanas de mantenimiento no deben exceder 4 horas en horarios de baja actividad. \\
\hline
\rowcolor{rowgray}
RNF-027 & Backup Automático & El sistema debe realizar respaldos automáticos diarios. \\
\hline
RNF-028 & Plan de Contingencia & Debe existir plan documentado con RTO de 24 horas. \\
\hline
\caption{Requisitos No Funcionales - Disponibilidad}
\end{longtable}

\subsection{RNF-5. Escalabilidad}

Estos requisitos aseguran que el sistema pueda crecer y adaptarse al aumento de usuarios, datos y torneos sin comprometer el rendimiento ni requerir rediseños arquitectónicos significativos.

\begin{longtable}{|>{\centering\arraybackslash}p{1.5cm}|>{\raggedright\arraybackslash}p{3.5cm}|>{\raggedright\arraybackslash}p{9cm}|}
\hline
\rowcolor{headercolor}
\textcolor{white}{\textbf{ID}} & \textcolor{white}{\textbf{Requisito}} & \textcolor{white}{\textbf{Descripción}} \\
\hline
\endfirsthead

\multicolumn{3}{c}{\textit{Continuación de la tabla anterior}} \\
\hline
\rowcolor{headercolor}
\textcolor{white}{\textbf{ID}} & \textcolor{white}{\textbf{Requisito}} & \textcolor{white}{\textbf{Descripción}} \\
\hline
\endhead

\hline
\multicolumn{3}{r}{\textit{Continúa en la siguiente página}} \\
\endfoot

\hline
\endlastfoot

RNF-029 & Crecimiento de Datos & El sistema debe soportar crecimiento de al menos 100GB sin pérdida de rendimiento. \\
\hline
\rowcolor{rowgray}
RNF-030 & Escalabilidad Horizontal & La arquitectura debe permitir agregar servidores adicionales. \\
\hline
RNF-031 & Carga de Archivos & El sistema debe soportar carga de archivos de hasta 50MB. \\
\hline
\rowcolor{rowgray}
RNF-032 & Número de Torneos & El sistema debe soportar al menos 50 torneos activos simultáneamente. \\
\hline
RNF-033 & Almacenamiento & El sistema debe optimizar mediante compresión de imágenes y archivos. \\
\hline
\caption{Requisitos No Funcionales - Escalabilidad}
\end{longtable}

\subsection{RNF-6. Mantenibilidad}

Estos requisitos facilitan el mantenimiento continuo del sistema, permitiendo identificar y corregir problemas rápidamente, así como implementar nuevas funcionalidades de manera eficiente.

\begin{longtable}{|>{\centering\arraybackslash}p{1.5cm}|>{\raggedright\arraybackslash}p{3.5cm}|>{\raggedright\arraybackslash}p{9cm}|}
\hline
\rowcolor{headercolor}
\textcolor{white}{\textbf{ID}} & \textcolor{white}{\textbf{Requisito}} & \textcolor{white}{\textbf{Descripción}} \\
\hline
\endfirsthead

\multicolumn{3}{c}{\textit{Continuación de la tabla anterior}} \\
\hline
\rowcolor{headercolor}
\textcolor{white}{\textbf{ID}} & \textcolor{white}{\textbf{Requisito}} & \textcolor{white}{\textbf{Descripción}} \\
\hline
\endhead

\hline
\multicolumn{3}{r}{\textit{Continúa en la siguiente página}} \\
\endfoot

\hline
\endlastfoot

RNF-034 & Código Documentado & El código fuente debe estar documentado siguiendo estándares. \\
\hline
\rowcolor{rowgray}
RNF-035 & Arquitectura Modular & El sistema debe seguir arquitectura modular para facilitar mantenimiento. \\
\hline
RNF-036 & Logs Detallados & El sistema debe generar logs detallados para diagnóstico. \\
\hline
\rowcolor{rowgray}
RNF-037 & Versionamiento & El código debe mantenerse en sistema de control de versiones (Git). \\
\hline
RNF-038 & Pruebas Automatizadas & Deben existir pruebas automatizadas con cobertura mínima del 70\%. \\
\hline
\caption{Requisitos No Funcionales - Mantenibilidad}
\end{longtable}

\subsection{RNF-7. Compatibilidad}

Estos requisitos garantizan que el sistema funcione correctamente en una amplia variedad de plataformas, navegadores y dispositivos, asegurando accesibilidad universal para todos los usuarios.

\begin{longtable}{|>{\centering\arraybackslash}p{1.5cm}|>{\raggedright\arraybackslash}p{3.5cm}|>{\raggedright\arraybackslash}p{9cm}|}
\hline
\rowcolor{headercolor}
\textcolor{white}{\textbf{ID}} & \textcolor{white}{\textbf{Requisito}} & \textcolor{white}{\textbf{Descripción}} \\
\hline
\endfirsthead

\multicolumn{3}{c}{\textit{Continuación de la tabla anterior}} \\
\hline
\rowcolor{headercolor}
\textcolor{white}{\textbf{ID}} & \textcolor{white}{\textbf{Requisito}} & \textcolor{white}{\textbf{Descripción}} \\
\hline
\endhead

\hline
\multicolumn{3}{r}{\textit{Continúa en la siguiente página}} \\
\endfoot

\hline
\endlastfoot

RNF-039 & Navegadores Web & El sistema debe funcionar en Chrome, Firefox, Safari y Edge (últimas 2 versiones). \\
\hline
\rowcolor{rowgray}
RNF-040 & Sistemas Operativos & El portal debe funcionar en Windows, macOS, Linux, iOS y Android. \\
\hline
RNF-041 & Resoluciones & La interfaz debe adaptarse desde 320px (móvil) hasta 2560px (4K). \\
\hline
\rowcolor{rowgray}
RNF-042 & Formatos de Archivo & El sistema debe soportar PDF, JPG, PNG, Excel, CSV. \\
\hline
\caption{Requisitos No Funcionales - Compatibilidad}
\end{longtable}

\subsection{RNF-8. Portabilidad}

Estos requisitos aseguran que el sistema no dependa de tecnologías propietarias específicas y pueda ser migrado o desplegado en diferentes entornos de infraestructura sin modificaciones significativas.

\begin{longtable}{|>{\centering\arraybackslash}p{1.5cm}|>{\raggedright\arraybackslash}p{3.2cm}|>{\raggedright\arraybackslash}p{9cm}|}
\hline
\rowcolor{headercolor}
\textcolor{white}{\textbf{ID}} & \textcolor{white}{\textbf{Requisito}} & \textcolor{white}{\textbf{Descripción}} \\
\hline
\endfirsthead

\multicolumn{3}{c}{\textit{Continuación de la tabla anterior}} \\
\hline
\rowcolor{headercolor}
\textcolor{white}{\textbf{ID}} & \textcolor{white}{\textbf{Requisito}} & \textcolor{white}{\textbf{Descripción}} \\
\hline
\endhead

\hline
\multicolumn{3}{r}{\textit{Continúa en la siguiente página}} \\
\endfoot

\hline
\endlastfoot

RNF-043 & Base de Datos & El sistema debe ser portable entre diferentes motores de base de datos relacionales. \\
\hline
\rowcolor{rowgray}
RNF-044 & Configuración Externalizada & Toda configuración debe estar en archivos externos. \\
\hline
RNF-045 & Contenedorización & El sistema debe poder ejecutarse en contenedores Docker. \\
\hline
\caption{Requisitos No Funcionales - Portabilidad}
\end{longtable}

\subsection{RNF-9. Cumplimiento Legal}

Estos requisitos garantizan que el sistema cumpla con todas las regulaciones aplicables sobre protección de datos personales, privacidad y derechos de los usuarios, evitando riesgos legales para la institución.

\begin{longtable}{|>{\centering\arraybackslash}p{1.5cm}|>{\raggedright\arraybackslash}p{3.5cm}|>{\raggedright\arraybackslash}p{9cm}|}
\hline
\rowcolor{headercolor}
\textcolor{white}{\textbf{ID}} & \textcolor{white}{\textbf{Requisito}} & \textcolor{white}{\textbf{Descripción}} \\
\hline
\endfirsthead

\multicolumn{3}{c}{\textit{Continuación de la tabla anterior}} \\
\hline
\rowcolor{headercolor}
\textcolor{white}{\textbf{ID}} & \textcolor{white}{\textbf{Requisito}} & \textcolor{white}{\textbf{Descripción}} \\
\hline
\endhead

\hline
\multicolumn{3}{r}{\textit{Continúa en la siguiente página}} \\
\endfoot

\hline
\endlastfoot

RNF-046 & Protección de Datos & El sistema debe cumplir con regulaciones de protección de datos personales. \\
\hline
\rowcolor{rowgray}
RNF-047 & Consentimiento & El sistema debe solicitar consentimiento explícito para tratamiento de datos. \\
\hline
RNF-048 & Derecho al Olvido & El sistema debe permitir eliminación completa de datos personales. \\
\hline
\rowcolor{rowgray}
RNF-049 & Registro de Consentimientos & El sistema debe mantener registro de consentimientos otorgados. \\
\hline
RNF-050 & Transparencia & Debe existir política de privacidad clara y accesible. \\
\hline
\caption{Requisitos No Funcionales - Cumplimiento Legal}
\end{longtable}

\newpage

\section{Resumen de Requisitos}

\subsection{Totales por Categoría}

\begin{table}[h]
\centering
\begin{tabular}{|l|c|}
\hline
\rowcolor{headercolor}
\textcolor{white}{\textbf{Categoría}} & \textcolor{white}{\textbf{Cantidad}} \\
\hline
Requisitos Funcionales (RF) & 199 \\
\hline
Requisitos No Funcionales (RNF) & 50 \\
\hline
\rowcolor{rowgray}
\textbf{TOTAL} & \textbf{249} \\
\hline
\end{tabular}
\caption{Resumen cuantitativo de requisitos del sistema}
\end{table}

\subsection{Distribución de Requisitos Funcionales}

\begin{table}[h]
\centering
\begin{tabular}{|l|l|c|}
\hline
\rowcolor{headercolor}
\textcolor{white}{\textbf{Módulo}} & \textcolor{white}{\textbf{Rango}} & \textcolor{white}{\textbf{Total}} \\
\hline
Autenticación y Seguridad & RF-001 a RF-010 & 10 \\
\hline
\rowcolor{rowgray}
Gestión de Usuarios & RF-011 a RF-021 & 11 \\
\hline
Gestión Institucional & RF-022 a RF-030 & 9 \\
\hline
\rowcolor{rowgray}
Disciplinas Deportivas & RF-031 a RF-040 & 10 \\
\hline
Gestión de Torneos & RF-041 a RF-052 & 12 \\
\hline
\rowcolor{rowgray}
Equipos y Participantes & RF-053 a RF-068 & 16 \\
\hline
Escenarios y Horarios & RF-069 a RF-078 & 10 \\
\hline
\rowcolor{rowgray}
Generación de Fixture & RF-079 a RF-091 & 13 \\
\hline
Ejecución de Partidos & RF-092 a RF-107 & 16 \\
\hline
\rowcolor{rowgray}
Resultados y Tablas & RF-108 a RF-119 & 12 \\
\hline
Justicia Deportiva & RF-120 a RF-132 & 13 \\
\hline
\rowcolor{rowgray}
Notificaciones & RF-133 a RF-146 & 14 \\
\hline
Portal Público & RF-147 a RF-158 & 12 \\
\hline
\rowcolor{rowgray}
Reportes y Exportación & RF-159 a RF-173 & 15 \\
\hline
Analítica y KPIs & RF-174 a RF-185 & 12 \\
\hline
\rowcolor{rowgray}
Configuración General & RF-186 a RF-199 & 14 \\
\hline
\textbf{TOTAL} & & \textbf{199} \\
\hline
\end{tabular}
\caption{Distribución de requisitos funcionales por módulo}
\end{table}



\newpage

\section{Matriz de Trazabilidad}

Esta sección establece la trazabilidad bidireccional entre requisitos, casos de uso, componentes arquitectónicos y casos de prueba. La matriz de trazabilidad es fundamental para:

\begin{itemize}[noitemsep]
    \item Verificar que todos los requisitos estén implementados en componentes específicos.
    \item Asegurar que cada requisito tenga casos de prueba asociados.
    \item Facilitar el análisis de impacto ante cambios de requisitos.
    \item Garantizar cobertura completa en validación y verificación.
    \item Mantener consistencia entre documentación de análisis, diseño y pruebas.
\end{itemize}

\subsection{Trazabilidad de Módulos Principales}

\begin{longtable}{|>{\centering\arraybackslash}p{1cm}|>{\raggedright\arraybackslash}p{3.5cm}|>{\raggedright\arraybackslash}p{3.5cm}|>{\raggedright\arraybackslash}p{3.5cm}|>{\centering\arraybackslash}p{2.5cm}|}
\hline
\rowcolor{headercolor}
\textcolor{white}{\textbf{\#}} & \textcolor{white}{\textbf{Módulo}} & \textcolor{white}{\textbf{Requisitos}} & \textcolor{white}{\textbf{Componente}} & \textcolor{white}{\textbf{Prioridad}} \\
\hline
\endfirsthead
\hline
\rowcolor{headercolor}
\textcolor{white}{\textbf{\#}} & \textcolor{white}{\textbf{Módulo}} & \textcolor{white}{\textbf{Requisitos}} & \textcolor{white}{\textbf{Componente}} & \textcolor{white}{\textbf{Prioridad}} \\
\hline
\endhead
\hline
\endlastfoot

1 & Autenticación y Seguridad & RF-001 a RF-010 & AuthService, AuthController, SecurityMiddleware & Alta \\
\hline
\rowcolor{rowgray}
2 & Gestión de Usuarios & RF-011 a RF-021 & UserService, UserController, RoleManager & Alta \\
\hline
3 & Gestión Institucional & RF-022 a RF-032 & InstitutionService, FacultyController & Alta \\
\hline
\rowcolor{rowgray}
4 & Disciplinas Deportivas & RF-033 a RF-043 & SportService, DisciplineController & Alta \\
\hline
5 & Gestión de Torneos & RF-044 a RF-058 & TournamentService, TournamentController & Crítica \\
\hline
\rowcolor{rowgray}
6 & Equipos y Participantes & RF-059 a RF-073 & TeamService, PlayerService, TeamController & Alta \\
\hline
7 & Escenarios y Horarios & RF-074 a RF-084 & VenueService, ScheduleService, VenueController & Media \\
\hline
\rowcolor{rowgray}
8 & Generación de Fixture & RF-085 a RF-093 & FixtureService, FixtureGenerator, AlgorithmEngine & Crítica \\
\hline
9 & Ejecución de Partidos & RF-094 a RF-107 & MatchService, MatchController, DigitalActa & Crítica \\
\hline
\rowcolor{rowgray}
10 & Resultados y Tablas & RF-108 a RF-119 & ResultService, StandingsService, StatsCalculator & Crítica \\
\hline
11 & Justicia Deportiva & RF-120 a RF-133 & JusticeService, SanctionController, AppealManager & Alta \\
\hline
\rowcolor{rowgray}
12 & Notificaciones & RF-134 a RF-144 & NotificationService, EmailService, NotifController & Media \\
\hline
13 & Portal Público & RF-145 a RF-158 & PublicController, PublicService, NewsService & Alta \\
\hline
\rowcolor{rowgray}
14 & Reportes y Exportación & RF-159 a RF-171 & ReportService, ExportService, PDFGenerator & Media \\
\hline
15 & Analítica y KPIs & RF-172 a RF-184 & AnalyticsService, KPICalculator, DashboardController & Media \\
\hline
\rowcolor{rowgray}
16 & Configuración General & RF-185 a RF-199 & ConfigService, SystemController, ParameterManager & Alta \\
\hline
\caption{Trazabilidad de Módulos a Componentes}
\end{longtable}

\subsection{Mapeo Requisitos-Casos de Uso}

\begin{longtable}{|>{\centering\arraybackslash}p{2cm}|>{\raggedright\arraybackslash}p{5cm}|>{\raggedright\arraybackslash}p{7.5cm}|}
\hline
\rowcolor{headercolor}
\textcolor{white}{\textbf{Caso de Uso}} & \textcolor{white}{\textbf{Descripción}} & \textcolor{white}{\textbf{Requisitos Relacionados}} \\
\hline
\endfirsthead
\hline
\rowcolor{headercolor}
\textcolor{white}{\textbf{Caso de Uso}} & \textcolor{white}{\textbf{Descripción}} & \textcolor{white}{\textbf{Requisitos Relacionados}} \\
\hline
\endhead
\hline
\endlastfoot

CU-001 & Iniciar Sesión en el Sistema & RF-001, RF-003, RF-004, RF-008, RNF-001, RNF-002 \\
\hline
\rowcolor{rowgray}
CU-002 & Recuperar Contraseña Olvidada & RF-005, RF-010, RNF-003, RNF-016 \\
\hline
CU-003 & Crear Torneo Deportivo & RF-044, RF-045, RF-046, RF-185, RNF-001, RNF-016 \\
\hline
\rowcolor{rowgray}
CU-004 & Inscribir Equipo en Torneo & RF-059, RF-060, RF-061, RF-062, RNF-016, RNF-020 \\
\hline
CU-005 & Generar Fixture Automático & RF-085, RF-086, RF-087, RF-088, RNF-004, RNF-038 \\
\hline
\rowcolor{rowgray}
CU-006 & Registrar Resultado de Partido & RF-094, RF-095, RF-108, RF-109, RNF-001, RNF-006 \\
\hline
CU-007 & Consultar Tabla de Posiciones & RF-108, RF-111, RF-145, RNF-004, RNF-018 \\
\hline
\rowcolor{rowgray}
CU-008 & Gestionar Sanción Disciplinaria & RF-120, RF-121, RF-122, RF-123, RNF-001, RNF-036 \\
\hline
CU-009 & Generar Reporte de Torneo & RF-159, RF-160, RF-161, RNF-004, RNF-034 \\
\hline
\rowcolor{rowgray}
CU-010 & Visualizar Dashboard Analítico & RF-172, RF-173, RF-174, RNF-004, RNF-018 \\
\hline
CU-011 & Configurar Parámetros del Sistema & RF-185, RF-186, RF-187, RNF-001, RNF-034 \\
\hline
\rowcolor{rowgray}
CU-012 & Gestionar Usuario del Sistema & RF-011, RF-012, RF-013, RF-014, RNF-001 \\
\hline
CU-013 & Administrar Escenarios Deportivos & RF-074, RF-075, RF-076, RNF-016, RNF-034 \\
\hline
\rowcolor{rowgray}
CU-014 & Crear Acta Digital de Partido & RF-094, RF-095, RF-096, RF-097, RNF-001, RNF-006 \\
\hline
CU-015 & Enviar Notificación Masiva & RF-134, RF-135, RF-136, RNF-005, RNF-024 \\
\hline
\caption{Mapeo de Casos de Uso a Requisitos}
\end{longtable}
\newpage
\subsection{Requisitos No Funcionales por Atributo de Calidad}

\begin{longtable}{|>{\raggedright\arraybackslash}p{4cm}|>{\centering\arraybackslash}p{3cm}|>{\raggedright\arraybackslash}p{7cm}|}
\hline
\rowcolor{headercolor}
\textcolor{white}{\textbf{Atributo ISO 25010}} & \textcolor{white}{\textbf{Requisitos}} & \textcolor{white}{\textbf{Mecanismo de Verificación}} \\
\hline
\endfirsthead
\hline
\rowcolor{headercolor}
\textcolor{white}{\textbf{Atributo ISO 25010}} & \textcolor{white}{\textbf{Requisitos}} & \textcolor{white}{\textbf{Mecanismo de Verificación}} \\
\hline
\endhead
\hline
\endlastfoot

Eficiencia de Desempeño & RNF-001 a RNF-005 & Pruebas de carga con JMeter, métricas APM \\
\hline
\rowcolor{rowgray}
Seguridad & RNF-006 a RNF-015 & Auditoría de seguridad, pruebas de penetración, OWASP ZAP \\
\hline
Usabilidad & RNF-016 a RNF-023 & Pruebas de usuario, evaluación heurística, métricas UX \\
\hline
\rowcolor{rowgray}
Fiabilidad & RNF-024 a RNF-028 & Monitoreo de uptime, simulación de fallos, logs de incidentes \\
\hline
Escalabilidad & RNF-029 a RNF-033 & Pruebas de volumen, análisis de crecimiento, benchmarking \\
\hline
\rowcolor{rowgray}
Mantenibilidad & RNF-034 a RNF-038 & Análisis estático de código, revisión de documentación, cobertura de pruebas \\
\hline
Compatibilidad & RNF-039 a RNF-043 & Cross-browser testing, pruebas multi-dispositivo \\
\hline
\rowcolor{rowgray}
Portabilidad & RNF-044 a RNF-047 & Despliegue en múltiples ambientes, containerización \\
\hline
Conformidad Legal & RNF-048 a RNF-050 & Revisión legal, auditoría de cumplimiento normativo \\
\hline
\caption{Trazabilidad de RNF a Atributos de Calidad ISO 25010}
\end{longtable}

\subsection{Matriz de Cobertura de Pruebas}

\begin{longtable}{|>{\centering\arraybackslash}p{2cm}|>{\raggedright\arraybackslash}p{3cm}|>{\raggedright\arraybackslash}p{3cm}|>{\raggedright\arraybackslash}p{6cm}|}
\hline
\rowcolor{headercolor}
\textcolor{white}{\textbf{Suite}} & \textcolor{white}{\textbf{Módulo Objetivo}} & \textcolor{white}{\textbf{Tipo Prueba}} & \textcolor{white}{\textbf{Requisitos Cubiertos}} \\
\hline
\endfirsthead
\hline
\rowcolor{headercolor}
\textcolor{white}{\textbf{Suite}} & \textcolor{white}{\textbf{Módulo Objetivo}} & \textcolor{white}{\textbf{Tipo Prueba}} & \textcolor{white}{\textbf{Requisitos Cubiertos}} \\
\hline
\endhead
\hline
\endlastfoot

TS-001 & Autenticación & Funcional & RF-001, RF-002, RF-003, RF-005, RF-008 \\
\hline
\rowcolor{rowgray}
TS-002 & Gestión Usuarios & Funcional & RF-011, RF-012, RF-013, RF-014, RF-016 \\
\hline
TS-003 & Torneos & Funcional & RF-044 a RF-058 \\
\hline
\rowcolor{rowgray}
TS-004 & Fixture Generator & Integración & RF-085 a RF-093 \\
\hline
TS-005 & Acta Digital & Funcional & RF-094 a RF-107 \\
\hline
\rowcolor{rowgray}
TS-006 & Resultados & Funcional & RF-108 a RF-119 \\
\hline
TS-007 & Portal Público & E2E & RF-145 a RF-158 \\
\hline
\rowcolor{rowgray}
TS-008 & Seguridad & No Funcional & RNF-006 a RNF-015 \\
\hline
TS-009 & Rendimiento & No Funcional & RNF-001 a RNF-005 \\
\hline
\rowcolor{rowgray}
TS-010 & Usabilidad & No Funcional & RNF-016 a RNF-023 \\
\hline
\caption{Suites de Prueba y Cobertura de Requisitos}
\end{longtable}

\subsection{Análisis de Impacto: Requisitos Críticos}

Los siguientes requisitos han sido identificados como críticos debido a su alto impacto en la funcionalidad core del sistema:

\begin{longtable}{|>{\centering\arraybackslash}p{2cm}|>{\raggedright\arraybackslash}p{4.5cm}|>{\raggedright\arraybackslash}p{7.5cm}|}
\hline
\rowcolor{headercolor}
\textcolor{white}{\textbf{Requisito}} & \textcolor{white}{\textbf{Justificación}} & \textcolor{white}{\textbf{Dependencias}} \\
\hline
\endfirsthead
\hline
\rowcolor{headercolor}
\textcolor{white}{\textbf{Requisito}} & \textcolor{white}{\textbf{Justificación}} & \textcolor{white}{\textbf{Dependencias}} \\
\hline
\endhead
\hline
\endlastfoot

RF-001 & Base de seguridad de todo el sistema & RF-003, RF-004, RNF-001, RNF-006 \\
\hline
\rowcolor{rowgray}
RF-044 & Funcionalidad principal: creación de torneos & RF-033, RF-045, RF-046, RF-074, RF-085 \\
\hline
RF-085 & Automatización central del sistema & RF-074, RF-044, RF-059, RF-087 \\
\hline
\rowcolor{rowgray}
RF-094 & Núcleo de la ejecución digital de eventos & RF-108, RF-109, RF-095, RF-096 \\
\hline
RF-108 & Procesamiento automático de resultados & RF-094, RF-109, RF-111, RF-112 \\
\hline
\rowcolor{rowgray}
RNF-001 & Rendimiento aceptable del sistema completo & Todos los RF que implican consultas DB \\
\hline
RNF-006 & Seguridad de datos sensibles & RF-001, RF-011, RNF-007, RNF-009 \\
\hline
\caption{Análisis de Requisitos Críticos y Dependencias}
\end{longtable}

\newpage

\subsection{Matriz de Trazabilidad Detallada}

La siguiente matriz establece la trazabilidad completa de todos los requisitos funcionales del sistema SIGED, identificando el origen de cada requisito, su prioridad, criterios de aceptación, vinculación con la EDT y estado actual. Esta matriz es fundamental para la gestión de requisitos, análisis de impacto y control de cambios durante todo el ciclo de vida del proyecto.

\textbf{Leyenda de Origen de Requisitos:}
\begin{itemize}[noitemsep]
    \item \textbf{LI:} Lluvia de Ideas (sesiones de brainstorming con equipo técnico)
    \item \textbf{INT:} Interesado (requerimientos directos del profesor/cliente)
    \item \textbf{BM:} Benchmarking (análisis de sistemas similares)
    \item \textbf{NOR:} Normativa (cumplimiento legal o estándares)
\end{itemize}

\textbf{Leyenda de Prioridades:}
\begin{itemize}[noitemsep]
    \item \textbf{C:} Crítica (esencial para funcionamiento core del sistema)
    \item \textbf{A:} Alta (importante para operación eficiente)
    \item \textbf{M:} Media (mejora la experiencia pero no es indispensable)
    \item \textbf{B:} Baja (deseable pero puede posponerse)
\end{itemize}

\textbf{Leyenda de Estados:}
\begin{itemize}[noitemsep]
    \item \textbf{P:} Pendiente (no iniciado)
    \item \textbf{D:} En Desarrollo (en construcción)
    \item \textbf{I:} Implementado (codificado y desplegado)
    \item \textbf{V:} Validado (probado y aceptado)
\end{itemize}

\textbf{Referencias a EDT:}

\begin{itemize}[noitemsep]
    \item \textbf{3.6:} Diseño de Base de Datos
    \item \textbf{4.1.1:} Backend - Autenticación y Seguridad
    \item \textbf{4.1.2:} Backend - Gestión de Usuarios y Permisos
    \item \textbf{4.1.3:} Backend - Lógica de Negocio (Torneos, Partidos, Fixture)
    \item \textbf{4.1.4:} Backend - Registro de Eventos y Actas
    \item \textbf{4.1.5:} Backend - Reportes y Analítica
    \item \textbf{4.2.1:} Frontend - Componentes UI Compartidos
    \item \textbf{4.2.2:} Frontend - Interfaces Operacionales
    \item \textbf{4.2.3:} Frontend - Portal Público
    \item \textbf{4.3.3:} Integración - APIs de Terceros
    \item \textbf{6.2:} Configuración de Servidores
    \item \textbf{6.6:} Configuración de Monitoreo
    \item \textbf{6.7:} Configuración de Respaldos
    \item \textbf{7.4:} Documentación
\end{itemize}


\clearpage

{\small
\begin{longtable}{|>{\centering\arraybackslash}p{1.5cm}|>{\raggedright\arraybackslash}p{3cm}|>{\centering\arraybackslash}p{1cm}|>{\centering\arraybackslash}p{1cm}|>{\centering\arraybackslash}p{1.2cm}|>{\raggedright\arraybackslash}p{3.5cm}|>{\centering\arraybackslash}p{1cm}|>{\centering\arraybackslash}p{1.3cm}|}
\hline
\rowcolor{headerblue}
\textcolor{white}{\textbf{ID}} &
\textcolor{white}{\textbf{Requisito}} &
\textcolor{white}{\textbf{Tipo}} &
\textcolor{white}{\textbf{Prior.}} &
\textcolor{white}{\textbf{Origen}} &
\textcolor{white}{\textbf{Criterio de Aceptación}} &
\textcolor{white}{\textbf{EDT}} &
\textcolor{white}{\textbf{Estado}} \\
\hline
\endfirsthead

\multicolumn{8}{c}{\textit{Continuación de Matriz de Trazabilidad}} \\
\hline
\rowcolor{headerblue}
\textcolor{white}{\textbf{ID}} &
\textcolor{white}{\textbf{Requisito}} &
\textcolor{white}{\textbf{Tipo}} &
\textcolor{white}{\textbf{Prior.}} &
\textcolor{white}{\textbf{Origen}} &
\textcolor{white}{\textbf{Criterio de Aceptación}} &
\textcolor{white}{\textbf{EDT}} &
\textcolor{white}{\textbf{Estado}} \\
\hline
\endhead

\hline
\multicolumn{8}{r}{\textit{Continúa en la siguiente página}} \\
\endfoot

\hline
\endlastfoot

\multicolumn{8}{|l|}{\cellcolor{headerblue}\textcolor{white}{\textbf{Módulo 1: Autenticación y Seguridad (RF-001 a RF-010)}}} \\
\hline
RF-001 & Login de Usuario & F & C & INT & Autenticación con email/contraseña en <2 seg & 4.1.1 & P \\
\hline
\rowcolor{rowgray}
RF-002 & Logout de Usuario & F & A & INT & Invalidar sesión y token JWT al cerrar sesión & 4.1.1 & P \\
\hline
RF-003 & Gestión de Sesiones & F & C & INT & Mantener sesión activa con JWT expirable & 4.1.1 & P \\
\hline
\rowcolor{rowgray}
RF-004 & Roles y Permisos (RBAC) & F & C & INT & Implementar jerarquía de 4 roles con permisos & 4.1.2 & P \\
\hline
RF-005 & Recuperación de Contraseña & F & A & INT & Enviar enlace de recuperación por email & 4.1.1 & P \\
\hline
\rowcolor{rowgray}
RF-006 & Cambio de Contraseña & F & A & INT & Permitir cambio con validación de contraseña actual & 4.1.2 & P \\
\hline
RF-007 & Auditoría de Acciones & F & A & NOR & Registrar acciones críticas con timestamp y usuario & 4.1.1 & P \\
\hline
\rowcolor{rowgray}
RF-008 & Bloqueo por Intentos Fallidos & F & A & NOR & Bloquear cuenta tras 5 intentos fallidos & 4.1.1 & P \\
\hline
RF-009 & Autenticación de Dos Factores & F & M & BM & Enviar código 2FA por email (opcional) & 4.1.1 & P \\
\hline
\rowcolor{rowgray}
RF-010 & Validación de Contraseñas & F & A & NOR & Exigir mín. 8 caracteres, mayús, minús y números & 4.1.1 & P \\
\hline

\multicolumn{8}{|l|}{\cellcolor{headerblue}\textcolor{white}{\textbf{Módulo 2: Gestión de Usuarios (RF-011 a RF-021)}}} \\
\hline
RF-011 & Creación de Administradores & F & A & INT & Super Admin crea cuentas admin con credenciales & 4.1.2 & P \\
\hline
\rowcolor{rowgray}
RF-012 & Creación de Comités & F & A & INT & Admins crean cuentas de miembros del comité & 4.1.2 & P \\
\hline
RF-013 & Creación de Delegados & F & A & INT & Comité crea cuentas de delegados de facultades & 4.1.2 & P \\
\hline
\rowcolor{rowgray}
RF-014 & Asignación de Roles & F & A & INT & Asignar y modificar roles según jerarquía & 4.1.2 & P \\
\hline
RF-015 & Gestión de Permisos & F & A & INT & Personalizar permisos específicos por rol & 4.1.2 & P \\
\hline
\rowcolor{rowgray}
RF-016 & Bloqueo/Desbloqueo de Cuentas & F & A & INT & Bloquear y desbloquear cuentas manualmente & 4.1.2 & P \\
\hline
RF-017 & Listado de Usuarios & F & M & INT & Mostrar usuarios con filtros por rol y estado & 4.1.2 & P \\
\hline
\rowcolor{rowgray}
RF-018 & Búsqueda de Usuarios & F & M & INT & Buscar por nombre, correo o rol & 4.1.2 & P \\
\hline
RF-019 & Edición de Datos & F & M & INT & Editar información básica de usuarios & 4.1.2 & P \\
\hline
\rowcolor{rowgray}
RF-020 & Eliminación Lógica & F & M & INT & Desactivar usuarios sin eliminar historial & 4.1.2 & P \\
\hline
RF-021 & Historial de Actividad & F & M & INT & Mostrar historial de acciones por usuario & 4.1.2 & P \\
\hline

\multicolumn{8}{|l|}{\cellcolor{headerblue}\textcolor{white}{\textbf{Módulo 3: Gestión Institucional (RF-022 a RF-030)}}} \\
\hline
RF-022 & Registro de Facultades & F & A & INT & Registrar facultades con código, nombre y logo & 3.6 & P \\
\hline
\rowcolor{rowgray}
RF-023 & Registro de Escuelas & F & A & INT & Registrar escuelas asociadas a facultades & 3.6 & P \\
\hline
RF-024 & Gestión de Logos & F & M & INT & Subir y gestionar logos institucionales & 4.2.1 & P \\
\hline
\rowcolor{rowgray}
RF-025 & Colores Institucionales & F & B & LI & Configurar colores para identificación visual & 4.2.1 & P \\
\hline
RF-026 & Asignación de Delegados & F & A & INT & Asignar delegados deportivos a cada facultad & 4.1.2 & P \\
\hline
\rowcolor{rowgray}
RF-027 & Datos Institucionales & F & M & INT & Gestionar información de contacto institucional & 3.6 & P \\
\hline
RF-028 & Activar/Desactivar & F & M & INT & Cambiar estado de facultades/escuelas & 4.1.3 & P \\
\hline
\rowcolor{rowgray}
RF-029 & Listado Institucional & F & M & INT & Listar facultades y escuelas con filtros & 4.2.1 & P \\
\hline
RF-030 & Estadísticas Institucionales & F & B & LI & Mostrar estadísticas de participación por facultad & 4.1.5 & P \\
\hline

\multicolumn{8}{|l|}{\cellcolor{headerblue}\textcolor{white}{\textbf{Módulo 4: Disciplinas Deportivas (RF-031 a RF-040)}}} \\
\hline
RF-031 & Crear Disciplina Deportiva & F & A & INT & Registrar nueva disciplina con nombre y tipo & 3.6 & P \\
\hline
\rowcolor{rowgray}
RF-032 & Configurar Reglas & F & A & NOR & Definir reglas específicas por disciplina & 3.6 & P \\
\hline
RF-033 & Sistema de Puntuación & F & C & NOR & Configurar sistema de puntos/victorias/goles & 4.1.3 & P \\
\hline
\rowcolor{rowgray}
RF-034 & Criterios de Desempate & F & A & NOR & Definir criterios automáticos de desempate & 4.1.3 & P \\
\hline
RF-035 & Plantillas de Disciplinas & F & M & BM & Usar plantillas predefinidas para deportes comunes & 3.6 & P \\
\hline
\rowcolor{rowgray}
RF-036 & Categorías por Género & F & A & NOR & Configurar categorías (masculino/femenino/mixto) & 3.6 & P \\
\hline
RF-037 & Tipos de Competencia & F & C & INT & Definir formato (todos vs todos, eliminatoria, grupos) & 4.1.3 & P \\
\hline
\rowcolor{rowgray}
RF-038 & Configuración de Eventos & F & M & INT & Configurar eventos múltiples en atletismo & 3.6 & P \\
\hline
RF-039 & Código de Vestimenta & F & B & NOR & Especificar requisitos de uniforme por disciplina & 3.6 & P \\
\hline
\rowcolor{rowgray}
RF-040 & Estado de Disciplinas & F & M & INT & Activar/desactivar disciplinas para torneos & 4.1.3 & P \\
\hline

\multicolumn{8}{|l|}{\cellcolor{headerblue}\textcolor{white}{\textbf{Módulo 5: Gestión de Torneos (RF-041 a RF-052)}}} \\
\hline
RF-041 & Crear Torneo & F & C & INT & Registrar torneo con nombre, disciplina y fechas & 4.1.3 & P \\
\hline
\rowcolor{rowgray}
RF-042 & Configurar Fechas & F & C & INT & Definir fechas de inicio, fin e inscripciones & 4.1.3 & P \\
\hline
RF-043 & Tipo de Torneo & F & C & INT & Seleccionar formato: round-robin, eliminación o grupos & 4.1.3 & P \\
\hline
\rowcolor{rowgray}
RF-044 & Estados de Torneo & F & C & INT & Gestionar estados: inscripción, en curso, finalizado & 4.1.3 & P \\
\hline
RF-045 & Publicar Torneo & F & A & INT & Publicar torneo para inscripciones públicas & 4.1.3 & P \\
\hline
\rowcolor{rowgray}
RF-046 & Cierre de Inscripciones & F & A & INT & Cerrar inscripciones automática o manualmente & 4.1.3 & P \\
\hline
RF-047 & Número de Equipos & F & A & NOR & Validar mínimo y máximo de equipos por torneo & 4.1.3 & P \\
\hline
\rowcolor{rowgray}
RF-048 & Comité Organizador & F & A & INT & Asignar miembros del comité al torneo & 4.1.3 & P \\
\hline
RF-049 & Premios y Reconocimientos & F & M & INT & Definir premiación por posiciones & 4.1.3 & P \\
\hline
\rowcolor{rowgray}
RF-050 & Duplicar Torneo & F & M & LI & Copiar configuración de torneo anterior & 4.1.3 & P \\
\hline
RF-051 & Cancelar Torneo & F & M & INT & Cancelar torneo con notificación a participantes & 4.1.3 & P \\
\hline
\rowcolor{rowgray}
RF-052 & Historial de Torneos & F & M & INT & Consultar torneos anteriores con estadísticas & 4.1.3 & P \\
\hline

\multicolumn{8}{|l|}{\cellcolor{headerblue}\textcolor{white}{\textbf{Módulo 6: Equipos y Participantes (RF-053 a RF-067)}}} \\
\hline
RF-053 & Registro de Equipos & F & A & INT & Inscribir equipo con nombre, facultad y delegado & 4.1.3 & P \\
\hline
\rowcolor{rowgray}
RF-054 & Inscripción de Jugadores & F & A & INT & Registrar jugadores con datos académicos & 4.1.3 & P \\
\hline
RF-055 & Validación de Alumnos & F & A & NOR & Verificar matrícula vigente en base institucional & 4.3.3 & P \\
\hline
\rowcolor{rowgray}
RF-056 & Asignar Capitán & F & M & INT & Designar capitán del equipo & 4.1.3 & P \\
\hline
RF-057 & Delegado de Equipo & F & A & INT & Asignar delegado responsable del equipo & 4.1.3 & P \\
\hline
\rowcolor{rowgray}
RF-058 & Validación Académica & F & A & NOR & Verificar créditos mínimos y promedio & 4.3.3 & P \\
\hline
RF-059 & Restricciones de Género & F & A & NOR & Aplicar reglas de género según categoría & 4.1.3 & P \\
\hline
\rowcolor{rowgray}
RF-060 & Número de Jugadores & F & A & NOR & Validar cantidad mín/máx de jugadores por equipo & 4.1.3 & P \\
\hline
RF-061 & Fotografía de Jugador & F & M & INT & Subir foto de identificación del jugador & 4.1.3 & P \\
\hline
\rowcolor{rowgray}
RF-062 & Documentos de Jugador & F & M & INT & Adjuntar documentos de verificación & 4.1.3 & P \\
\hline
RF-063 & Estado de Jugador & F & A & INT & Gestionar estado: activo, suspendido, retirado & 4.1.3 & P \\
\hline
\rowcolor{rowgray}
RF-064 & Transferencia de Jugador & F & M & INT & Permitir cambio de equipo con aprobación & 4.1.3 & P \\
\hline
RF-065 & Listado de Equipos & F & M & INT & Mostrar equipos inscritos con filtros & 4.2.1 & P \\
\hline
\rowcolor{rowgray}
RF-066 & Buscar Jugador & F & M & INT & Buscar jugadores por nombre o código & 4.2.1 & P \\
\hline
RF-067 & Historial de Jugador & F & M & INT & Ver participación histórica en torneos & 4.1.5 & P \\
\hline

\multicolumn{8}{|l|}{\cellcolor{headerblue}\textcolor{white}{\textbf{Módulo 7: Escenarios y Horarios (RF-068 a RF-077)}}} \\
\hline
RF-068 & Registro de Canchas & F & M & INT & Registrar escenarios deportivos con nombre y tipo & 3.6 & P \\
\hline
\rowcolor{rowgray}
RF-069 & Disponibilidad de Horarios & F & M & INT & Configurar horarios disponibles por cancha & 3.6 & P \\
\hline
RF-070 & Capacidad de Espectadores & F & B & INT & Especificar aforo de cada escenario & 3.6 & P \\
\hline
\rowcolor{rowgray}
RF-071 & Tipo de Superficie & F & B & INT & Definir superficie: cesped, sintético, parquet, etc & 3.6 & P \\
\hline
RF-072 & Detección de Conflictos & F & A & NOR & Detectar solapamiento de horarios automáticamente & 4.1.3 & P \\
\hline
\rowcolor{rowgray}
RF-073 & Reserva de Canchas & F & A & INT & Reservar escenarios para partidos del fixture & 4.1.3 & P \\
\hline
RF-074 & Periodos de Mantenimiento & F & M & INT & Programar períodos de mantenimiento de canchas & 4.1.3 & P \\
\hline
\rowcolor{rowgray}
RF-075 & Características de Cancha & F & B & INT & Detallar servicios: iluminación, graderías, etc & 3.6 & P \\
\hline
RF-076 & Ubicación Geográfica & F & B & LI & Mostrar ubicación en mapa interactivo & 4.2.1 & P \\
\hline
\rowcolor{rowgray}
RF-077 & Calendario de Uso & F & M & INT & Visualizar calendario de ocupación de canchas & 4.2.1 & P \\
\hline

\multicolumn{8}{|l|}{\cellcolor{headerblue}\textcolor{white}{\textbf{Módulo 8: Generación de Fixture (RF-079 a RF-091)}}} \\
\hline
RF-079 & Generación Automática & F & C & INT & Generar fixture automáticamente según formato & 4.1.3 & P \\
\hline
\rowcolor{rowgray}
RF-080 & Algoritmo Round-Robin & F & C & NOR & Generar todos contra todos con ida y vuelta & 4.1.3 & P \\
\hline
RF-081 & Algoritmo Eliminación & F & C & NOR & Generar llaves eliminatorias con cuadro de llaves & 4.1.3 & P \\
\hline
\rowcolor{rowgray}
RF-082 & Restricciones de Horario & F & A & INT & Respetar horarios disponibles de canchas & 4.1.3 & P \\
\hline
RF-083 & Restricciones de Equipos & F & A & INT & Evitar confrontación prematura entre equipos & 4.1.3 & P \\
\hline
\rowcolor{rowgray}
RF-084 & Distribución Equitativa & F & A & NOR & Balancear carga de partidos entre equipos & 4.1.3 & P \\
\hline
RF-085 & Reprogramar Partidos & F & A & INT & Cambiar fecha/hora/cancha de partido & 4.1.3 & P \\
\hline
\rowcolor{rowgray}
RF-086 & Validar Conflictos & F & A & NOR & Detectar conflictos en fixture generado & 4.1.3 & P \\
\hline
RF-087 & Fechas de Descanso & F & M & INT & Programar fechas libres entre jornadas & 4.1.3 & P \\
\hline
\rowcolor{rowgray}
RF-088 & Publicar Fixture & F & A & INT & Publicar calendario oficial de partidos & 4.1.3 & P \\
\hline
RF-089 & Notificar Fixture & F & A & INT & Enviar fixture por email a delegados y equipos & 4.1.3 & P \\
\hline
\rowcolor{rowgray}
RF-090 & Fechas Dobles & F & M & LI & Permitir múltiples partidos en mismo día & 4.1.3 & P \\
\hline
RF-091 & Grupos + Eliminación & F & A & INT & Generar fase de grupos seguida de eliminatoria & 4.1.3 & P \\
\hline

\multicolumn{8}{|l|}{\cellcolor{headerblue}\textcolor{white}{\textbf{Módulo 9: Ejecución de Partidos (RF-092 a RF-107)}}} \\
\hline
RF-092 & Registro en Tiempo Real & F & C & INT & Registrar eventos del partido en tiempo real & 4.1.4 & P \\
\hline
\rowcolor{rowgray}
RF-093 & Cronómetro de Partido & F & A & INT & Controlar tiempo de juego con inicio/pausa/fin & 4.2.2 & P \\
\hline
RF-094 & Registro de Goles/Puntos & F & C & NOR & Anotar goles con jugador, minuto y asistencia & 4.1.4 & P \\
\hline
\rowcolor{rowgray}
RF-095 & Registro de Tarjetas & F & C & NOR & Registrar tarjetas (amarilla/roja) con motivo & 4.1.4 & P \\
\hline
RF-096 & Registro de Cambios & F & A & NOR & Registrar sustituciones con jugador sale/entra & 4.1.4 & P \\
\hline
\rowcolor{rowgray}
RF-097 & Registro de Faltas & F & M & NOR & Contabilizar faltas por equipo y jugador & 4.1.4 & P \\
\hline
RF-098 & Firma Digital de Acta & F & A & NOR & Firmar acta electrónica por capitanes y árbitro & 4.1.4 & P \\
\hline
\rowcolor{rowgray}
RF-099 & Estados de Partido & F & C & INT & Gestionar: programado, en juego, finalizado, suspendido & 4.1.4 & P \\
\hline
RF-100 & Alineación de Equipos & F & A & NOR & Registrar jugadores titulares y suplentes & 4.1.4 & P \\
\hline
\rowcolor{rowgray}
RF-101 & Observaciones del Partido & F & M & INT & Agregar notas y observaciones al acta & 4.1.4 & P \\
\hline
RF-102 & Registro de Lesiones & F & M & NOR & Documentar lesiones ocurridas en partido & 4.1.4 & P \\
\hline
\rowcolor{rowgray}
RF-103 & Tiempo Extra/Penales & F & A & NOR & Registrar prórroga y lanzamiento de penales & 4.1.4 & P \\
\hline
RF-104 & Verificación de Asistencia & F & A & NOR & Confirmar asistencia de equipos antes del partido & 4.1.4 & P \\
\hline
\rowcolor{rowgray}
RF-105 & Walkover Automático & F & A & NOR & Dar victoria por ausencia si equipo no se presenta & 4.1.4 & P \\
\hline
RF-106 & Editar Acta Finalizada & F & M & INT & Permitir corrección de acta con justificación & 4.1.4 & P \\
\hline
\rowcolor{rowgray}
RF-107 & Notificación de Inicio & F & M & INT & Enviar recordatorio antes del partido & 4.1.4 & P \\
\hline

\multicolumn{8}{|l|}{\cellcolor{headerblue}\textcolor{white}{\textbf{Módulo 10: Resultados y Tablas (RF-108 a RF-119)}}} \\
\hline
RF-108 & Cálculo Automático & F & C & NOR & Calcular puntos, goles, diferencia automáticamente & 4.1.4 & P \\
\hline
\rowcolor{rowgray}
RF-109 & Tabla de Posiciones & F & C & NOR & Mostrar clasificación ordenada por puntos y desempate & 4.1.5 & P \\
\hline
RF-110 & Criterios de Desempate & F & A & NOR & Aplicar desempate: diferencia, goles, enfrentamiento & 4.1.5 & P \\
\hline
\rowcolor{rowgray}
RF-111 & Clasificación Automática & F & C & NOR & Determinar equipos clasificados a siguiente fase & 4.1.5 & P \\
\hline
RF-112 & Estadísticas Generales & F & A & INT & Mostrar promedios, totales y comparativas & 4.1.5 & P \\
\hline
\rowcolor{rowgray}
RF-113 & Tabla de Goleadores & F & A & INT & Ranking de jugadores por goles anotados & 4.1.5 & P \\
\hline
RF-114 & Estadísticas de Tarjetas & F & M & INT & Contabilizar tarjetas por jugador y equipo & 4.1.5 & P \\
\hline
\rowcolor{rowgray}
RF-115 & Diferencia de Goles & F & A & NOR & Calcular diferencia acumulada por equipo & 4.1.5 & P \\
\hline
RF-116 & Rachas de Equipos & F & M & LI & Mostrar rachas de victorias/derrotas consecutivas & 4.1.5 & P \\
\hline
\rowcolor{rowgray}
RF-117 & Resultados Históricos & F & M & INT & Consultar resultados de temporadas anteriores & 4.1.5 & P \\
\hline
RF-118 & Comparación de Equipos & F & M & LI & Comparar estadísticas entre equipos & 4.1.5 & P \\
\hline
\rowcolor{rowgray}
RF-119 & Predicciones de Resultados & F & B & LI & Sugerir posibles resultados según historial (ML) & 4.1.5 & P \\
\hline

\multicolumn{8}{|l|}{\cellcolor{headerblue}\textcolor{white}{\textbf{Módulo 11: Justicia Deportiva (RF-120 a RF-132)}}} \\
\hline
RF-120 & Presentar Reclamo & F & A & NOR & Registrar reclamo con descripción y evidencias & 4.1.4 & P \\
\hline
\rowcolor{rowgray}
RF-121 & Adjuntar Evidencias & F & A & NOR & Subir documentos, fotos, actas como pruebas & 4.1.4 & P \\
\hline
RF-122 & Tribunal Deportivo & F & A & NOR & Asignar comité para evaluar reclamos & 4.1.4 & P \\
\hline
\rowcolor{rowgray}
RF-123 & Estados de Reclamo & F & A & INT & Seguimiento: presentado, en revisión, resuelto & 4.1.4 & P \\
\hline
RF-124 & Notificar Resoluciones & F & A & NOR & Enviar resolución del tribunal a las partes & 4.1.4 & P \\
\hline
\rowcolor{rowgray}
RF-125 & Sanciones Automáticas & F & A & NOR & Aplicar sanciones según tarjetas (2 amarillas = suspensión) & 4.1.4 & P \\
\hline
RF-126 & Sanciones Manuales & F & A & NOR & Tribunal aplica sanciones: multas, suspensiones & 4.1.4 & P \\
\hline
\rowcolor{rowgray}
RF-127 & Suspensión de Jugador & F & A & NOR & Suspender jugador por partidos específicos & 4.1.4 & P \\
\hline
RF-128 & Suspensión de Equipo & F & A & NOR & Suspender equipo completo de torneo & 4.1.4 & P \\
\hline
\rowcolor{rowgray}
RF-129 & Historial Disciplinario & F & M & INT & Ver antecedentes de jugadores/equipos & 4.1.5 & P \\
\hline
RF-130 & Apelaciones & F & M & NOR & Permitir apelar decisiones dentro de plazo & 4.1.4 & P \\
\hline
\rowcolor{rowgray}
RF-131 & Cumplimiento de Sanciones & F & A & NOR & Verificar automáticamente cumplimiento de sanciones & 4.1.4 & P \\
\hline
RF-132 & Puntos Disciplinarios & F & B & LI & Sistema de puntos acumulativos por infracciones & 4.1.4 & P \\
\hline

\multicolumn{8}{|l|}{\cellcolor{headerblue}\textcolor{white}{\textbf{Módulo 12: Notificaciones (RF-133 a RF-146)}}} \\
\hline
RF-133 & Notificación por Email & F & A & INT & Enviar notificaciones sistemáticas por correo & 4.1.5 & P \\
\hline
\rowcolor{rowgray}
RF-134 & Notificaciones In-App & F & A & INT & Mostrar alertas dentro de la aplicación & 4.2.1 & P \\
\hline
RF-135 & Alertas de Cambios & F & A & INT & Notificar modificaciones en fixture/resultados & 4.1.5 & P \\
\hline
\rowcolor{rowgray}
RF-136 & Recordatorios & F & A & INT & Envíar recordatorios de partidos y plazos & 4.1.5 & P \\
\hline
RF-137 & Notificar Resultados & F & M & INT & Avisar cuando se publiquen resultados de partidos & 4.1.5 & P \\
\hline
\rowcolor{rowgray}
RF-138 & Confirmación de Inscripción & F & M & INT & Confirmar inscripción exitosa a torneos & 4.1.5 & P \\
\hline
RF-139 & Notificar Sanciones & F & A & NOR & Informar de sanciones aplicadas & 4.1.5 & P \\
\hline
\rowcolor{rowgray}
RF-140 & Notificar Reclamos & F & A & NOR & Avisar cambios en estado de reclamos & 4.1.5 & P \\
\hline
RF-141 & Preferencias de Notificación & F & M & INT & Configurar canales y tipos de notificaciones & 4.2.1 & P \\
\hline
\rowcolor{rowgray}
RF-142 & Historial de Notificaciones & F & M & INT & Visualizar notificaciones anteriores & 4.1.5 & P \\
\hline
RF-143 & Notificaciones Push & F & B & BM & Enviar push a dispositivos móviles (futuro) & 4.3.3 & P \\
\hline
\rowcolor{rowgray}
RF-144 & Notificar Mantenimiento & F & M & INT & Alertar de mantenimiento programado del sistema & 6.6 & P \\
\hline
RF-145 & Notificación por SMS & F & B & BM & Enviar SMS urgentes (opcional) & 4.3.3 & P \\
\hline
\rowcolor{rowgray}
RF-146 & Notificar Logros & F & B & LI & Notificar premios y reconocimientos (gamificación) & 4.1.5 & P \\
\hline

\multicolumn{8}{|l|}{\cellcolor{headerblue}\textcolor{white}{\textbf{Módulo 13: Portal Público (RF-147 a RF-158)}}} \\
\hline
RF-147 & Resultados en Vivo & F & A & INT & Mostrar marcadores actualizados en tiempo real & 4.2.3 & P \\
\hline
\rowcolor{rowgray}
RF-148 & Consultar Fixture Público & F & A & INT & Ver calendario completo de partidos sin login & 4.2.3 & P \\
\hline
RF-149 & Visualizar Tablas Públicas & F & A & INT & Mostrar tabla de posiciones actual & 4.2.3 & P \\
\hline
\rowcolor{rowgray}
RF-150 & Sección de Noticias & F & M & LI & Publicar noticias y comunicados deportivos & 4.2.3 & P \\
\hline
RF-151 & Galería de Fotos & F & B & LI & Mostrar fotos de eventos deportivos & 4.2.3 & P \\
\hline
\rowcolor{rowgray}
RF-152 & Estadísticas Públicas & F & A & INT & Ver estadísticas de equipos y jugadores & 4.2.3 & P \\
\hline
RF-153 & Calendario de Eventos & F & M & INT & Mostrar próximos eventos deportivos & 4.2.3 & P \\
\hline
\rowcolor{rowgray}
RF-154 & Buscar Equipos Público & F & M & INT & Buscar equipos participantes sin autenticación & 4.2.3 & P \\
\hline
RF-155 & Buscar Jugadores Público & F & M & INT & Buscar perfil de jugadores & 4.2.3 & P \\
\hline
\rowcolor{rowgray}
RF-156 & Diseño Responsive & F & A & NOR & Adaptación automática a móvil y tablet & 4.2.3 & P \\
\hline
RF-157 & Compartir en Redes Sociales & F & M & BM & Botón para compartir resultados y estadísticas & 4.2.3 & P \\
\hline
\rowcolor{rowgray}
RF-158 & Descargar Fixture PDF & F & M & INT & Exportar fixture a PDF desde portal público & 4.1.5 & P \\
\hline

\multicolumn{8}{|l|}{\cellcolor{headerblue}\textcolor{white}{\textbf{Módulo 14: Reportes (RF-159 a RF-173)}}} \\
\hline
RF-159 & Generar PDFs & F & A & INT & Generar reportes en formato PDF & 4.1.5 & P \\
\hline
\rowcolor{rowgray}
RF-160 & Exportar a Excel & F & A & INT & Exportar datos a hojas de cálculo & 4.1.5 & P \\
\hline
RF-161 & Reporte de Actas & F & A & NOR & Generar PDF de acta oficial de partido & 4.1.5 & P \\
\hline
\rowcolor{rowgray}
RF-162 & Reporte de Tablas & F & M & INT & Generar reporte de tabla de posiciones & 4.1.5 & P \\
\hline
RF-163 & Reporte de Fixture & F & M & INT & Exportar fixture completo a PDF/Excel & 4.1.5 & P \\
\hline
\rowcolor{rowgray}
RF-164 & Estadísticas Individuales & F & M & INT & Reporte de desempeño por jugador & 4.1.5 & P \\
\hline
RF-165 & Estadísticas de Equipos & F & M & INT & Reporte consolidado por equipo & 4.1.5 & P \\
\hline
\rowcolor{rowgray}
RF-166 & Reporte Institucional & F & M & INT & Reporte ejecutivo para autoridades & 4.1.5 & P \\
\hline
RF-167 & Reporte de Participación & F & M & INT & Estadísticas de participación por facultad & 4.1.5 & P \\
\hline
\rowcolor{rowgray}
RF-168 & Uso de Instalaciones & F & M & INT & Reporte de ocupación de canchas & 4.1.5 & P \\
\hline
RF-169 & Reporte de Sanciones & F & M & NOR & Listado de sanciones aplicadas & 4.1.5 & P \\
\hline
\rowcolor{rowgray}
RF-170 & Generar Certificados & F & M & INT & Certificados de participación/premiación & 7.4 & P \\
\hline
RF-171 & Ranking de Goleadores & F & M & INT & Tabla de máximos anotadores exportable & 4.1.5 & P \\
\hline
\rowcolor{rowgray}
RF-172 & Reporte Comparativo & F & B & LI & Comparar temporadas/torneos anteriores & 4.1.5 & P \\
\hline
RF-173 & Exportar Actas Históricas & F & M & INT & Descargar actas de torneos pasados & 4.1.5 & P \\
\hline

\multicolumn{8}{|l|}{\cellcolor{headerblue}\textcolor{white}{\textbf{Módulo 15: Analítica y KPIs (RF-174 a RF-185)}}} \\
\hline
RF-174 & Dashboard Analítico & F & M & INT & Panel con indicadores clave de rendimiento & 4.1.5 & P \\
\hline
\rowcolor{rowgray}
RF-175 & Participación por Facultad & F & M & INT & KPI de participación por unidad académica & 4.1.5 & P \\
\hline
RF-176 & Uso de Canchas (KPI) & F & M & INT & Métrica de ocupación de instalaciones & 4.1.5 & P \\
\hline
\rowcolor{rowgray}
RF-177 & Rendimiento de Torneos & F & M & INT & Métricas de éxito por tipo de torneo & 4.1.5 & P \\
\hline
RF-178 & Análisis de Tendencias & F & B & LI & Proyecciones y tendencias históricas & 4.1.5 & P \\
\hline
\rowcolor{rowgray}
RF-179 & Análisis de Asistencia & F & B & LI & Tasa de cumplimiento de equipos a partidos & 4.1.5 & P \\
\hline
RF-180 & Disciplinas Populares & F & M & INT & Ranking de deportes más practicados & 4.1.5 & P \\
\hline
\rowcolor{rowgray}
RF-181 & Tiempo Promedio de Partidos & F & B & LI & Métrica de duración por disciplina & 4.1.5 & P \\
\hline
RF-182 & Tasa de Reclamos & F & M & INT & KPI de conflictos por torneo & 4.1.5 & P \\
\hline
\rowcolor{rowgray}
RF-183 & Cumplimiento de Horarios & F & M & INT & Métrica de puntualidad en partidos & 4.1.5 & P \\
\hline
RF-184 & Exportar Analítica & F & M & INT & Descargar datos de KPIs a Excel/PDF & 4.1.5 & P \\
\hline
\rowcolor{rowgray}
RF-185 & Comparación Interanual & F & B & LI & Comparar métricas entre años académicos & 4.1.5 & P \\
\hline

\multicolumn{8}{|l|}{\cellcolor{headerblue}\textcolor{white}{\textbf{Módulo 16: Configuración (RF-186 a RF-199)}}} \\
\hline
RF-186 & Parámetros del Sistema & F & A & INT & Configurar valores globales del sistema & 6.2 & P \\
\hline
\rowcolor{rowgray}
RF-187 & Periodos Académicos & F & A & NOR & Definir semestres y ciclos académicos & 3.6 & P \\
\hline
RF-188 & Límites Globales & F & M & NOR & Configurar límites de equipos, jugadores, etc & 3.6 & P \\
\hline
\rowcolor{rowgray}
RF-189 & Tiempos de Sesión & F & M & NOR & Configurar timeout de sesión y tokens & 6.2 & P \\
\hline
RF-190 & Plantillas de Email & F & M & INT & Personalizar plantillas de notificaciones & 7.4 & P \\
\hline
\rowcolor{rowgray}
RF-191 & Backup Automático & F & A & NOR & Programar respaldos automáticos de datos & 6.7 & P \\
\hline
RF-192 & Restaurar Backup & F & A & NOR & Recuperar sistema desde respaldo & 6.7 & P \\
\hline
\rowcolor{rowgray}
RF-193 & Logs del Sistema & F & A & NOR & Visualizar y exportar registros de actividad & 6.6 & P \\
\hline
RF-194 & Roles Personalizados & F & M & LI & Crear roles con permisos customizados & 4.1.2 & P \\
\hline
\rowcolor{rowgray}
RF-195 & Mantenimiento Programado & F & M & INT & Programar ventanas de mantenimiento & 6.2 & P \\
\hline
RF-196 & Gestión de Versiones & F & M & INT & Control de versiones y releases del sistema & 6.2 & P \\
\hline
\rowcolor{rowgray}
RF-197 & Configurar Integraciones & F & M & BM & Configurar APIs externas (OAuth, LDAP, etc) & 4.3.3 & P \\
\hline
RF-198 & Personalización UI & F & B & LI & Configurar logo, colores y tema institucional & 4.2.1 & P \\
\hline
\rowcolor{rowgray}
RF-199 & Configuración Regional & F & M & NOR & Definir idioma, zona horaria y formato de fecha & 6.2 & P \\
\hline

\caption{Matriz de Trazabilidad de Requisitos Funcionales - SIGDE}
\end{longtable}
}
\newpage
\subsection{Notas sobre la Matriz de Trazabilidad}

\textbf{Estructura de la Matriz:}
La Matriz de Trazabilidad presenta 8 columnas que documentan cada requisito funcional:
\begin{itemize}[noitemsep]
    \item \textbf{ID:} Identificador único del requisito (RF-001 a RF-199)
    \item \textbf{Requisito:} Nombre descriptivo del requisito funcional
    \item \textbf{Tipo:} Clasificación del requisito (F = Funcional, actualmente todos son tipo F)
    \item \textbf{Prior.:} Prioridad (C = Crítica, A = Alta, M = Media, B = Baja)
    \item \textbf{Origen:} Fuente del requisito según metodología aplicada
    \item \textbf{Criterio de Aceptación:} Condición verificable para considerar el requisito completo
    \item \textbf{EDT:} Referencia al paquete de trabajo en la Estructura de Desglose del Trabajo
    \item \textbf{Estado:} Estado actual del requisito (P = Pendiente, D = En Desarrollo, I = Implementado, V = Validado)
\end{itemize}



\textbf{Códigos de Prioridad:}
\begin{itemize}[noitemsep]
    \item \textbf{C (Crítica):} Requisito esencial para el funcionamiento del sistema, debe implementarse en MVP
    \item \textbf{A (Alta):} Requisito muy importante que afecta funcionalidad principal
    \item \textbf{M (Media):} Requisito deseable que mejora la experiencia pero no es crítico
    \item \textbf{B (Baja):} Requisito opcional que puede implementarse en fases posteriores
\end{itemize}


\textbf{Uso de la Matriz:}
\begin{itemize}[noitemsep]
    \item \textbf{Planificación de Sprints:} Agrupar requisitos por módulo y prioridad para definir iteraciones
    \item \textbf{Asignación de Trabajo:} Mapear requisitos con paquetes EDT para distribución de tareas
    \item \textbf{Verificación de Completitud:} Validar que todos los requisitos tengan criterios de aceptación claros
    \item \textbf{Seguimiento de Progreso:} Actualizar estado de requisitos conforme avanza el desarrollo
    \item \textbf{Análisis de Impacto:} Al modificar un requisito, identificar rápidamente módulos relacionados
    \item \textbf{Gestión de Cambios:} Documentar y rastrear modificaciones en requisitos con trazabilidad
\end{itemize}

\subsection{Matriz de Trazabilidad de Requisitos No Funcionales}

La siguiente matriz documenta los 50 requisitos no funcionales del sistema, organizados por atributo de calidad según el estándar ISO/IEC 25010. Cada requisito incluye su criterio de verificación, prioridad, métrica de aceptación y referencia al paquete EDT correspondiente.

\clearpage

{\small
\begin{longtable}{|>{\centering\arraybackslash}p{1.2cm}|>{\raggedright\arraybackslash}p{3.5cm}|>{\centering\arraybackslash}p{1.8cm}|>{\centering\arraybackslash}p{1cm}|>{\raggedright\arraybackslash}p{4.5cm}|>{\centering\arraybackslash}p{1cm}|>{\centering\arraybackslash}p{1cm}|}
\hline
\rowcolor{headerblue}
\textcolor{white}{\textbf{ID}} &
\textcolor{white}{\textbf{Requisito}} &
\textcolor{white}{\textbf{Atributo}} &
\textcolor{white}{\textbf{Prior.}} &
\textcolor{white}{\textbf{Métrica de Verificación}} &
\textcolor{white}{\textbf{EDT}} &
\textcolor{white}{\textbf{Estado}} \\
\hline
\endfirsthead

\multicolumn{7}{c}{\textit{Continuación de la tabla anterior}} \\
\hline
\rowcolor{headerblue}
\textcolor{white}{\textbf{ID}} &
\textcolor{white}{\textbf{Requisito}} &
\textcolor{white}{\textbf{Atributo}} &
\textcolor{white}{\textbf{Prior.}} &
\textcolor{white}{\textbf{Métrica de Verificación}} &
\textcolor{white}{\textbf{EDT}} &
\textcolor{white}{\textbf{Estado}} \\
\hline
\endhead

\hline
\multicolumn{7}{r}{\textit{Continúa en la siguiente página}} \\
\endfoot

\hline
\endlastfoot

\multicolumn{7}{|l|}{\cellcolor{headerblue}\textcolor{white}{\textbf{Rendimiento (RNF-001 a RNF-005)}}} \\
\hline
RNF-001 & Tiempo de Respuesta & REND & C & Respuesta \textless 2 seg en 95\% de consultas & 4.1.3 & P \\
\hline
\rowcolor{rowgray}
RNF-002 & Carga Concurrente & REND & C & 500 usuarios simultáneos sin degradación & 6.2 & P \\
\hline
RNF-003 & Tiempo de Carga & REND & A & Páginas públicas cargan \textless 3 seg (LCP) & 4.2.3 & P \\
\hline
\rowcolor{rowgray}
RNF-004 & Procesamiento Fixture & REND & A & Generación 20 equipos \textless10 seg & 4.1.3 & P \\
\hline
RNF-005 & Actualización Tiempo Real & REND & A & Latencia marcadores \textless5 seg & 4.1.4 & P \\
\hline

\multicolumn{7}{|l|}{\cellcolor{headerblue}\textcolor{white}{\textbf{Seguridad (RNF-006 a RNF-015)}}} \\
\hline
RNF-006 & Encriptación HTTPS & SEG & C & 100\% comunicaciones con TLS 1.3 & 6.2 & P \\
\hline
\rowcolor{rowgray}
RNF-007 & Encriptación Contraseñas & SEG & C & Bcrypt/Argon2 con salt único & 4.1.1 & P \\
\hline
RNF-008 & Protección SQL Injection & SEG & C & Consultas parametrizadas en 100\% queries & 4.1.3 & P \\
\hline
\rowcolor{rowgray}
RNF-009 & Protección CSRF & SEG & C & Tokens CSRF en formularios críticos & 4.1.1 & P \\
\hline
RNF-010 & Validación de Entrada & SEG & C & Validación server-side en 100\% endpoints & 4.1.3 & P \\
\hline
\rowcolor{rowgray}
RNF-011 & Control de Acceso RBAC & SEG & C & Verificación permisos en operaciones críticas & 4.1.2 & P \\
\hline
RNF-012 & Auditoría de Seguridad & SEG & A & Log de intentos fallidos con IP y timestamp & 4.1.1 & P \\
\hline
\rowcolor{rowgray}
RNF-013 & Sesiones Seguras & SEG & A & Tokens expiran en 24h, refresh en 4h & 4.1.1 & P \\
\hline
RNF-014 & Validación de Archivos & SEG & A & Verificar tipo MIME y tamaño \textless10MB & 4.1.3 & P \\
\hline
\rowcolor{rowgray}
RNF-015 & Backups Encriptados & SEG & A & Respaldos con AES-256 & 6.7 & P \\
\hline

\multicolumn{7}{|l|}{\cellcolor{headerblue}\textcolor{white}{\textbf{Usabilidad (RNF-016 a RNF-023)}}} \\
\hline
RNF-016 & Interfaz Intuitiva & USAB & A & \textless1 hora capacitación para operaciones básicas & 4.2.1 & P \\
\hline
\rowcolor{rowgray}
RNF-017 & Navegación Clara & USAB & A & Máximo 3 clics para alcanzar funcionalidad & 4.2.1 & P \\
\hline
RNF-018 & Diseño Responsive & USAB & C & Adaptación 320px - 2560px sin scroll horizontal & 4.2.1 & P \\
\hline
\rowcolor{rowgray}
RNF-019 & Accesibilidad WCAG & USAB & A & Cumplir WCAG 2.1 nivel AA (87+ score Lighthouse) & 4.2.1 & P \\
\hline
RNF-020 & Mensajes de Error & USAB & M & Mensajes específicos con sugerencia de solución & 4.2.1 & P \\
\hline
\rowcolor{rowgray}
RNF-021 & Retroalimentación Visual & USAB & M & Feedback \textless200ms en interacciones (spinners) & 4.2.1 & P \\
\hline
RNF-022 & Consistencia UI & USAB & M & Sistema de diseño único (colores, tipografía) & 4.2.1 & P \\
\hline
\rowcolor{rowgray}
RNF-023 & Ayuda Contextual & USAB & B & Tooltips en campos complejos & 4.2.1 & P \\
\hline

\multicolumn{7}{|l|}{\cellcolor{headerblue}\textcolor{white}{\textbf{Disponibilidad (RNF-024 a RNF-028)}}} \\
\hline
RNF-024 & Uptime 99.5\% & DISP & C & Máximo 3.6 horas inactividad mensual & 6.2 & P \\
\hline
\rowcolor{rowgray}
RNF-025 & Recuperación Rápida & DISP & A & Recuperación de fallos menores \textless5 min & 6.2 & P \\
\hline
RNF-026 & Ventana Mantenimiento & DISP & M & Mantenimientos \textless4h en horario bajo uso & 6.2 & P \\
\hline
\rowcolor{rowgray}
RNF-027 & Backup Automático & DISP & C & Respaldo diario automático a las 02:00 AM & 6.7 & P \\
\hline
RNF-028 & Plan de Contingencia & DISP & A & RTO 24h, RPO 24h documentado & 6.7 & P \\
\hline

\multicolumn{7}{|l|}{\cellcolor{headerblue}\textcolor{white}{\textbf{Escalabilidad (RNF-029 a RNF-033)}}} \\
\hline
RNF-029 & Crecimiento de Datos & ESCA & A & Soportar +1TB sin degradación \textgreater 10\% & 3.6 & P \\
\hline
\rowcolor{rowgray}
RNF-030 & Escalabilidad Horizontal & ESCA & A & Arquitectura permite agregar nodos & 6.2 & P \\
\hline
RNF-031 & Carga de Archivos & ESCA & M & Archivos hasta 200MB con progress bar & 4.1.3 & P \\
\hline
\rowcolor{rowgray}
RNF-032 & Torneos Simultáneos & ESCA & A & Mínimo 50 torneos activos concurrentes & 4.1.3 & P \\
\hline
RNF-033 & Optimización Storage & ESCA & M & Compresión imágenes (WebP) y logs rotativos & 6.2 & P \\
\hline

\multicolumn{7}{|l|}{\cellcolor{headerblue}\textcolor{white}{\textbf{Mantenibilidad (RNF-034 a RNF-038)}}} \\
\hline
RNF-034 & Código Documentado & MANT & A & Comentarios JSDoc/XML en funciones públicas & 7.4 & P \\
\hline
\rowcolor{rowgray}
RNF-035 & Arquitectura Modular & MANT & C & Separación capas: presentación, lógica, datos & 4.1.3 & P \\
\hline
RNF-036 & Logs Estructurados & MANT & A & Logs en formato JSON con niveles (info, warn, error) & 6.6 & P \\
\hline
\rowcolor{rowgray}
RNF-037 & Control de Versiones & MANT & C & Git con branching strategy (GitFlow) & 7.4 & P \\
\hline
RNF-038 & Pruebas Automatizadas & MANT & A & Cobertura ≥ 70\% en lógica negocio & 7.4 & P \\
\hline

\multicolumn{7}{|l|}{\cellcolor{headerblue}\textcolor{white}{\textbf{Compatibilidad (RNF-039 a RNF-042)}}} \\
\hline
RNF-039 & Navegadores Modernos & COMP & C & Chrome, Firefox, Safari, Edge (últimas 2 versiones) & 4.2.1 & P \\
\hline
\rowcolor{rowgray}
RNF-040 & Sistemas Operativos & COMP & A & Windows, macOS, Linux, iOS, Android & 4.2.1 & P \\
\hline
RNF-041 & Resoluciones Pantalla & COMP & A & Soporte 320px móvil hasta 2560px 4K & 4.2.1 & P \\
\hline
\rowcolor{rowgray}
RNF-042 & Formatos de Archivo & COMP & M & PDF, JPG, PNG, XLSX, CSV import/export & 4.1.5 & P \\
\hline

\multicolumn{7}{|l|}{\cellcolor{headerblue}\textcolor{white}{\textbf{Portabilidad (RNF-043 a RNF-045)}}} \\
\hline
RNF-043 & Base de Datos Portable & PORT & A & Compatible PostgreSQL, MySQL, SQL Server & 3.6 & P \\
\hline
\rowcolor{rowgray}
RNF-044 & Configuración Externa & PORT & A & Variables entorno y config files (appsettings.json) & 6.2 & P \\
\hline
RNF-045 & Contenedorización & PORT & A & Dockerfile y docker-compose funcionales & 6.2 & P \\
\hline

\multicolumn{7}{|l|}{\cellcolor{headerblue}\textcolor{white}{\textbf{Cumplimiento Legal (RNF-046 a RNF-050)}}} \\
\hline
RNF-046 & Protección de Datos & LEGAL & C & Cumplir ley protección datos personales nacional & 4.1.2 & P \\
\hline
\rowcolor{rowgray}
RNF-047 & Consentimiento Usuario & LEGAL & C & Checkbox explícito para uso de datos & 4.1.2 & P \\
\hline
RNF-048 & Derecho al Olvido & LEGAL & A & Eliminación completa de datos en \textless30 días & 4.1.2 & P \\
\hline
\rowcolor{rowgray}
RNF-049 & Registro Consentimientos & LEGAL & A & Auditoría de consentimientos con timestamp & 4.1.2 & P \\
\hline
RNF-050 & Política de Privacidad & LEGAL & C & Política accesible desde todas las páginas & 4.2.3 & P \\
\hline

\caption{Matriz de Trazabilidad de Requisitos No Funcionales - SIGDE}
\end{longtable}
}

\subsection{Notas sobre la Matriz de RNF}

\textbf{Códigos de Atributos de Calidad:}
\begin{itemize}[noitemsep]
    \item \textbf{REND:} Rendimiento - Tiempo de respuesta y capacidad de procesamiento
    \item \textbf{SEG:} Seguridad - Protección de datos y prevención de amenazas
    \item \textbf{USAB:} Usabilidad - Facilidad de uso y experiencia de usuario
    \item \textbf{DISP:} Disponibilidad - Uptime y recuperación ante fallos
    \item \textbf{ESCA:} Escalabilidad - Capacidad de crecimiento del sistema
    \item \textbf{MANT:} Mantenibilidad - Facilidad de mantenimiento y evolución
    \item \textbf{COMP:} Compatibilidad - Soporte multiplataforma
    \item \textbf{PORT:} Portabilidad - Capacidad de despliegue en diferentes entornos
    \item \textbf{LEGAL:} Cumplimiento Legal - Conformidad con normativas
\end{itemize}

\textbf{Niveles de Prioridad:}
\begin{itemize}[noitemsep]
    \item \textbf{C (Crítica):} 26 requisitos - Obligatorios para MVP, impactan funcionamiento básico
    \item \textbf{A (Alta):} 19 requisitos - Muy importantes para calidad del producto
    \item \textbf{M (Media):} 4 requisitos - Deseables pero no bloqueantes
    \item \textbf{B (Baja):} 1 requisito - Mejoras futuras
\end{itemize}

\textbf{Estrategia de Verificación:}
Los requisitos no funcionales se verificarán mediante:
\begin{itemize}[noitemsep]
    \item \textbf{Pruebas de Rendimiento:} JMeter/k6 para carga y respuesta (RNF-001 a RNF-005)
    \item \textbf{Auditorías de Seguridad:} OWASP ZAP, análisis estático con SonarQube (RNF-006 a RNF-015)
    \item \textbf{Pruebas de Usabilidad:} Google Lighthouse, tests con usuarios reales (RNF-016 a RNF-023)
    \item \textbf{Monitoreo de Disponibilidad:} Uptime monitoring con Prometheus/Grafana (RNF-024 a RNF-028)
    \item \textbf{Pruebas de Escalabilidad:} Load testing incremental (RNF-029 a RNF-033)
    \item \textbf{Análisis de Código:} SonarQube para cobertura y complejidad (RNF-034 a RNF-038)
    \item \textbf{Pruebas Cross-browser:} BrowserStack o selenium grid (RNF-039 a RNF-042)
    \item \textbf{Revisión Legal:} Auditoría de cumplimiento normativo (RNF-046 a RNF-050)
\end{itemize}

\subsection{Versiones y Control de Cambios de Requisitos}

\begin{longtable}{|>{\centering\arraybackslash}p{1.2cm}|>{\centering\arraybackslash}p{1.5cm}|>{\raggedright\arraybackslash}p{2.5cm}|>{\raggedright\arraybackslash}p{6cm}|}
\hline
\rowcolor{headercolor}
\textcolor{white}{\textbf{Versión}} & \textcolor{white}{\textbf{Fecha}} & \textcolor{white}{\textbf{Autor}} & \textcolor{white}{\textbf{Cambios Realizados}} \\
\hline
\endfirsthead
\hline
\rowcolor{headercolor}
\textcolor{white}{\textbf{Versión}} & \textcolor{white}{\textbf{Fecha}} & \textcolor{white}{\textbf{Autor}} & \textcolor{white}{\textbf{Cambios Realizados}} \\
\hline
\endhead
\hline
\endlastfoot

1.0 & 2026-02-11 & Equipo SIGED & Versión inicial completa de ERS con 249 requisitos \\
\hline
\rowcolor{rowgray}
 &  &  & Inclusión de 16 módulos funcionales principales \\
\hline
 &  &  & Especificación de requisitos del proyecto y del producto \\
\hline
\rowcolor{rowgray}
 &  &  & Matriz de trazabilidad completa \\
\hline
\caption{Historial de Versiones del Documento}
\end{longtable}

\textit{Nota: Los cambios futuros en requisitos deben documentarse en esta tabla, incluyendo fecha, responsable, descripción del cambio y análisis de impacto en otros requisitos, componentes y casos de prueba.}

\newpage

\section{Conclusiones y Recomendaciones}

El presente documento de Especificación de Requisitos de Software (ERS) ha establecido de manera exhaustiva y sistemática la totalidad de requisitos que conforman el Sistema de Gestión de Eventos Deportivos (SIGED). La especificación comprende 8 Requisitos de Proyecto (RP), 199 Requisitos Funcionales (RF) distribuidos en 16 módulos funcionales, 50 Requisitos No Funcionales (RNF) organizados en 9 atributos de calidad según ISO/IEC 25010, y una Matriz de Trazabilidad completa que vincula requisitos, casos de uso, componentes arquitectónicos y casos de prueba.

Este documento ha sido elaborado siguiendo los lineamientos y mejores prácticas definidos por IEEE 830-1998 (Recommended Practice for Software Requirements Specifications), ISO/IEC 25010:2011 (Modelo de calidad de software SQuaRE), SWEBOK v3.0 (Área de conocimiento de ingeniería de requisitos), y PMI PMBOK para requisitos de gestión de proyecto.

La arquitectura modular propuesta, basada en 16 módulos funcionales independientes e interoperables, garantiza mantenibilidad mediante cambios localizados sin impacto en otros módulos, escalabilidad horizontal para incorporar nuevos módulos sin rediseño arquitectónico, reutilización de componentes y servicios compartidos entre módulos, y testabilidad con pruebas independientes por módulo. Los requisitos no funcionales especificados establecen umbrales claros y medibles: tiempos de respuesta inferiores a 2 segundos para operaciones estándar, conformidad con estándares ISO 27001 y OWASP Top 10 en seguridad, SLA del 99.5\% de disponibilidad durante períodos de competencia activa, y curva de aprendizaje inferior a 1 hora para usuarios básicos.

Para maximizar las probabilidades de éxito del proyecto, se recomienda implementar desarrollo iterativo e incremental con sprints de 3 semanas y entregables funcionales validables, realizar revisiones de sprint con usuarios finales para validación continua, priorizar estratégicamente los módulos críticos (Autenticación, Torneos, Fixture, Partidos, Resultados), establecer pipeline CI/CD desde sprint 1 con cobertura mínima del 70\%, validar interfaces de usuario mediante prototipos antes de implementación backend, incorporar auditorías de seguridad en cada fase del desarrollo, mantener documentación técnica actualizada en paralelo con desarrollo, e iniciar formación de usuarios finales al completar 60\% del desarrollo.

Los riesgos identificados en la sección de requisitos del proyecto requieren monitoreo activo durante todo el ciclo de desarrollo mediante el establecimiento de comité de control de cambios para gestionar modificaciones a requisitos, implementación de plan de retención de talento para minimizar rotación del equipo técnico, realización de pruebas de carga progresivas desde sprints tempranos, y ejecución de auditorías de seguridad al menos cada 6 semanas.

El proyecto se considerará exitoso cuando se satisfagan los siguientes criterios: implementación completa de requisitos funcionales de prioridad Alta y Media (100\%), cumplimiento verificable de todos los requisitos no funcionales críticos, superación satisfactoria de pruebas de aceptación de usuario (UAT) sin defectos críticos, documentación técnica completa y aprobada por comité de revisión, sistema desplegado en producción y operativo con disponibilidad del 99.5\%, usuarios clave capacitados con certificación de competencia, y cobertura de código de al menos 95\% en pruebas automatizadas.

Completada la especificación de requisitos, los siguientes pasos del ciclo de vida del proyecto son: revisión y aprobación formal de este documento por todos los stakeholders clave, elaboración del Documento de Arquitectura (SAD) con diseño de arquitectura de software detallada, diseño detallado con diagramas UML (casos de uso, secuencia, clases, componentes, despliegue), diseño de base de datos con modelo entidad-relación y optimización de índices, prototipos de UI/UX con wireframes y mockups para validación, plan detallado de pruebas basado en matriz de trazabilidad, configuración de entorno de desarrollo (repositorio, CI/CD, herramientas, estándares de código), y Sprint 0 con configuración inicial de proyecto y autenticación MVP.

El equipo de ingeniería de requisitos certifica que el presente documento cumple con los estándares IEEE 830-1998 e ISO/IEC 25010:2011, ha sido elaborado con participación activa de stakeholders representativos, refleja fielmente las necesidades institucionales identificadas en fase de elicitación, establece bases sólidas y verificables para las fases subsecuentes del proyecto, y proporciona trazabilidad completa desde requisitos hasta implementación y pruebas.

\vspace{1cm}

\textit{Este documento constituye la línea base oficial de requisitos del proyecto SIGED. Cualquier modificación posterior debe ser gestionada mediante el proceso formal de control de cambios establecido en la metodología del proyecto, documentando el cambio en la matriz de trazabilidad (Sección 6) e impactos en componentes relacionados.}

\vspace{1.5cm}

\noindent\rule{\textwidth}{0.5pt}

\vspace{0.5cm}

\begin{center}
\textbf{Elaborado por:} Equipo de Ingeniería de Requisitos - SIGED\\
\textbf{Fecha de Emisión:} 11 de Febrero de 2026\\
\textbf{Versión del Documento:} 1.0\\
\textbf{Estado:} Aprobado para Revisión de Stakeholders
\end{center}

\vspace{0.5cm}

\noindent\rule{\textwidth}{0.5pt}
\newpage
\section{Referencias}

\begin{itemize}[noitemsep]
    \item IEEE Std 830-1998: IEEE Recommended Practice for Software Requirements Specifications.
    \item ISO/IEC 25010:2011: Systems and software engineering — Systems and software Quality Requirements and Evaluation (SQuaRE).
    \item SWEBOK v3.0: Guide to the Software Engineering Body of Knowledge.
    \item ISO/IEC 27001:2013: Information technology — Security techniques — Information security management systems.
    \item WCAG 2.1: Web Content Accessibility Guidelines.
    \item PMI PMBOK Guide: A Guide to the Project Management Body of Knowledge.
\end{itemize}

\end{document}
// Separar datos en services/data
// services/liveMatchesService.ts
export const liveMatchesService = {
  getLiveMatches: async () => { /* API call */ },
  getMockLiveMatches: () => [...]  // Para desarrollo
}

// ✅ Componente limpio
export default function Index() {
  const { data: liveMatches } = useLiveMatches(); // Custom hook
  return <HomeView matches={liveMatches} />
}